% DER_II-b: Grothendieck's text with the commentary set in, then the overview, inside a group that applies the chunk's own definitions
\dgrpart{1}{II-b}{Part 2 of 2, manuscript pages 54--99}{II-b\enspace\textperiodcentered\enspace Chapter II, part 2 of 2\enspace\textperiodcentered\enspace manuscript pages 54--99}{Chapter II \textperiodcentered\ \textup{II-b}}
\dgrchunkbegin
% DER_II-b: the definitions of the chunk's preamble, applied inside its group
\providecommand{\IfDocumentMetadataT}[1]{}
\let\printedpage\relax\newcommand{\printedpage}[1]{\marginpar{\footnotesize p.~#1}}
\def\dcprintedpending{}
\let\dcjoinedprintedpage\relax\newcommand{\dcjoinedprintedpage}[1]{\def\dcprintedpending{#1}}
\let\msp\relax\newcommand{\msp}[1]{\marginpar{\footnotesize\raggedright
  \ifx\dcprintedpending\empty\else p.~\dcprintedpending\par\fi
  [#1]}\def\dcprintedpending{}}
\let\Hom\relax\newcommand{\Hom}{\operatorname{Hom}}
\let\uHom\relax\newcommand{\uHom}{\underline{\operatorname{Hom}}}
\let\Ob\relax\newcommand{\Ob}{\operatorname{Ob}}
\let\bs\relax\newcommand{\bs}{\backslash}
\let\cM\relax\newcommand{\cM}{\mathcal{M}}
\let\dbl\relax\newcommand{\dbl}[2]{\mathrel{\overset{#1}{\underset{#2}{\rightrightarrows}}}}
\let\defeq\relax\newcommand{\defeq}{\overset{\text{déf}}{=}}
\let\bD\relax\newcommand{\bD}{\mathbf{D}}
\let\bDelta\relax\newcommand{\bDelta}{\boldsymbol{\Delta}}
\let\D\relax\newcommand{\D}{\mathbf{D}}
\let\cP\relax\newcommand{\cP}{\mathcal{P}}
\let\cQ\relax\newcommand{\cQ}{\mathcal{Q}}
\let\cR\relax\newcommand{\cR}{\mathcal{R}}
\let\fP\relax\newcommand{\fP}{\mathfrak{P}}
\let\sR\relax\newcommand{\sR}{{}^{s}\mathcal{R}}
\let\circled\relax\newcommand{\circled}[1]{\tikz[baseline=(c.base)]\node[draw,rounded corners=2pt,inner sep=1.5pt](c){#1};}
\let\Fl\relax\newcommand{\Fl}{\mathrm{Fl}}
\let\Pf\relax\newcommand{\Pf}{\mathfrak{P}}
\let\Is\relax\newcommand{\Is}{\mathrm{Is}}
\let\Fib\relax\newcommand{\Fib}{\mathrm{Fib}}
\let\margmark\relax\newcommand{\margmark}[1]{\marginpar{\footnotesize #1}}
\let\PS\relax\newcommand{\PS}{\mathrm{PS}}
\let\tl\relax\newcommand{\tl}[1]{\begin{tabular}{@{}l@{}}#1\end{tabular}}
\let\tls\relax\newcommand{\tls}[1]{\text{\footnotesize\begin{tabular}{@{}l@{}}#1\end{tabular}}}

% ---- from Der-ChII-234_p036-040.tex
\printedpage{34}
% head: A. Grothendieck : Dérivateurs, Ch. II, Édition des Universités de Montpellier II et Paris VII
Ceci montre l'équivalence de d, d$'$, qui se voit aussi par l'equivalence de chacune de ces deux conditions avec d$''$, par le critère de cofinalité du paragraphe précedent. Soit D cette condition sur $u$. On a par (4.8) D $\Longrightarrow$ A, donc % [sic] l'equivalence -> l'équivalence, accent missing; précedent -> précédent, accent missing
\[
\mathrm{D} \Longrightarrow \mathrm{A} \Longrightarrow \mathrm{c},
\]
il reste à prouver que c $\Longrightarrow$ D, plus précisément c $\Longrightarrow$ d$''$ et c$'$ $\Longrightarrow$ d$''$ \texttt{[c$'$?]}. Explicitons le morphisme d'adjonction (pour $i$ dans $I$)
\[
u_!\underbrace{u^*i} \longrightarrow i,
\]
i.e. pour $j$ dans $I$
\[
\underbrace{(u_!u^*i)(j)}_{\pi_0(j\backslash(I'/u^*(i)))} \longrightarrow i(j) = \Hom(j,i).
\]

\msp{page 54}Or $I'/u^*(i) \simeq I'/i$, donc on a la flèche
\[
(u_!u^*(i))(j) \simeq \pi_0(j\backslash(I'/i)) \longrightarrow \Hom_I(j,i). \tag*{(4.10)}
\]
Or on voit aussitôt que pour $\alpha : j \longrightarrow i$ dans $\Hom(j,i)$, la fibre de (4.10) en $u$ s'identifie à $\pi_0(j\backslash\underset{\alpha}{I'}/i\ )$, de façon plus précise, $j\backslash\underset{\alpha}{I'}/i$\ \ est la sous-catégorie ouverte et fermée de $J = j\backslash(I'/i)$ correspondant à la partie du $\pi_0(J)$ qui est au dessus de $\alpha$ dans (4.10). Cela achève la démonstration de 4.1.

\dgrblock{dgrb117}{Reading the two-sided criterion}
\begin{dcomment}{exposition}{p54:Cela achève la démonstration de 4.1.}{Reading the two-sided criterion}
Grothendieck has now converted full faithfulness of restriction into a condition on individual arrows. This is the point at which the cofinality theorem becomes a tool for reconstructing an entire diagram. The commentary at manuscript p.~53 (``Biasphericity, lax epimorphisms, and factorization categories'') already discusses these pages; we retain its notation and add two clarifications needed by the argument that follows.

For $u:A\to B$ and $\gamma:b_0\to b_1$, write $\operatorname{Fact}_u(\gamma)$ for the category whose objects are triples $(a,r,s)$ with $r:b_0\to u(a)$, $s:u(a)\to b_1$, and $sr=\gamma$. A morphism to $(a',r',s')$ is $h:a\to a'$ with $u(h)r=r'$ and $s'u(h)=s$. The map of connected components in (4.10) has precisely $\pi_0\operatorname{Fact}_u(\gamma)$ as its fibre over $\gamma$. Its being a bijection therefore says that every arrow has a factorization through $u$, unique up to a finite zigzag of such factorizations. This is stronger than merely obtaining a connected comma category at each object.

The direction in 4.1(d$'$) must accordingly be read with care. For $B/b_1$, the relevant comma category is $\gamma\backslash(A/b_1)$, so its projection must be \emph{cofinal} in the convention of Definition 3.3. For $b_0\backslash B$, the relevant category is $(b_0\backslash A)/\gamma$, so the dual projection must be \emph{co-cofinal}. The edition prints ``cofinal'' in both places. Taken literally, the second condition is too weak: the inclusion $\{1\}\to[1]$ satisfies it, although restriction of presheaves to the terminal object is not fully faithful. The corrected orientation is also the one used in Proposition 4.4(b).

There is a separate existence issue in 4.1(a$'$), already isolated in the commentary at manuscript p.~51 (``The missing pointwise limit hypothesis''). The left Kan extension requires the colimits over the appropriate comma categories, not merely colimits of shapes $I$ and $I'$. Working in a cocomplete coefficient category is sufficient. Once the adjoints exist, the full-faithfulness equivalences are formal, and the representable counit test is valid because both $u_!u^*$ and the identity preserve colimits. No comparable argument makes the representable \emph{unit} test sufficient: its exact meaning is different.
\end{dcomment}


NB Une autre façon de prouver que c $\Longrightarrow$ A c'est de noter que les deux membres dans le morphisme d'adjonction $u_!u^*F \longrightarrow F$ dans $I^{\wedge}$ sont des foncteurs $I^{\wedge} \longrightarrow I^{\wedge}$ qui commutent aux $\varinjlim$ quelconques. Donc pour vérifier que c'est \texttt{[un]} isomorphisme, il suffit de le vérifier pour $F$ dans $I$, puisque $I$ engendre $I^{\wedge}$ par $\varinjlim$ . Par contre, un argument de ce type ne semble pas prouver que pour vérifier que l'on a un \textit{isomorphisme} $F \longrightarrow u_*u^*F$, il suffise de le faire pour $F$ représentable, i.e. de vérifier b$''$ pour $F$ représentable.

Il resterait à voir si c$'$ $\Longrightarrow$ d$''$ \texttt{[c$'$?]}.

\msp{page 55}Explicitons le morphisme d'adjonction
\[
i \longrightarrow u_*u^*i
\]
(pour $i \in \Ob I$) par ses valeurs sur $j$
\[
\underbrace{i(j)}_{\Hom(j,i)} \longrightarrow \Gamma(I'/j, u^*i). \tag*{$(\ast)$}
\]
Soit donc
\[
u_j : I'/j \longrightarrow I,
\]
% foot: Version 8 juillet 2001
\printedpage{35}
% head: A. Grothendieck : Dérivateurs, Ch. II, Édition des Universités de Montpellier II et Paris VII
on a
\[
\begin{aligned}
\Gamma(I'/j, \underbrace{u^*(i)|I'/j}_{u_j^*(i)}) &\simeq \Hom(e_{I'/j}, u^*(i)) \\ % ?? uncertain: extent of the underbrace -- over "u^*(i)|I'/j" or including the closing parenthesis
&\overset{\mathrm{adj}}{\simeq} \Hom(u_{j,!}e_{I'/j}, i),
\end{aligned}
\]
et la flèche $(\ast)$, fonctorielle en $i$, est déduite d'une flèche
\[
u_{j,!}(e_{I'/j}) \longrightarrow j. \tag*{$(\ast\ast)$}
\]
On voit tout de suite, comme ci-dessus, que $(\ast\ast)$ est \texttt{[un]} isomorphisme si et seulement si pour tout $\alpha : k \longrightarrow j$, la catégorie $k\backslash\underset{\alpha}{I'}/j$ est 0-connexe, et si on l'exige pour tout $j$, $k$, $\alpha$, c'est la condition d$''$. Mais il n'est pas clair que pour qu'on ait $(\ast\ast)$ isomorphisme, il suffise que le Hom des deux membres à valeurs dans tout $i$ dans $I$ (et non tout $F$ dans $I^{\wedge}$) soit \texttt{[un]} isomorphisme. La difficulté provient du fait que $u_{j,!}(e_{I'/j})$ n'est pas forcément représentable, et qu'un $F$ quelconque dans $I^{\wedge}$ \msp{page 56}n'est pas forcément $\varprojlim$ de foncteurs représentables. Donc il est douteux que c$'$ \texttt{[c$'$?]} suffise à entraîner que $u$ est $W_0$-connexe.

\dgrblock{dgrb118}{The unit on representables tests density}
\begin{dcomment}{examples}{p56:Donc il est douteux que}{The unit on representables tests density}
Grothendieck's hesitation on pages 54--56 identifies a genuine distinction. He asks whether knowing every representable presheaf can be recovered by right Kan extension already recovers all presheaves. The condition has a familiar modern interpretation: it says that $u$ is dense. Density reconstructs the objects of the original category from their restricted Yoneda images; it does not reconstruct its entire presheaf category.

Indeed, let $y_B(b)=B(-,b)$. By adjunction, the component at $c$ of the unit on $y_B(b)$ is the map
\[
 B(c,b)\longrightarrow
 \operatorname{Nat}\bigl(B(u-,c),B(u-,b)\bigr).
\]
Thus the unit is invertible on all representables exactly when the restricted Yoneda functor $B\to\operatorname{Set}^{A^{\mathrm{op}}}$ is fully faithful. This is the definition of a dense functor; see \citet[Definition 6.5.11]{Riehl2016}. It has a different quantifier from the full faithfulness of $u^*$ on \emph{all} presheaves.

A finite example settles the question without any existence qualifications. Let $B$ be the full subcategory of sets on the objects $1=\{*\}$ and $2=\{0,1\}$, and let $u:e\to B$ select $1$. For $b,c\in\{1,2\}$, evaluation on elements gives
\[
 B(c,b)\xrightarrow{\sim}
 \operatorname{Set}\bigl(B(1,c),B(1,b)\bigr).
\]
Both sides are simply the set of functions $c\to b$, and the displayed map is the identity under the indicated identifications. Hence every unit on a representable is an isomorphism. On the other hand, $\operatorname{Fact}_u(\mathrm{id}_2)$ is empty: any composite $2\to1\to2$ is constant, whereas $\mathrm{id}_2$ is not. Proposition 4.1 therefore shows that $u$ is not $W_0$-connected.

One can see the failure directly at the counit. The presheaf $u_!u^*y_B(2)$ evaluated at $2$ parametrizes factorizations $2\to1\to2$; its image in $B(2,2)$ consists of the two constant functions and misses the identity and the transposition. The representable unit has nevertheless passed every test. The two adjunction maps carry different information.

This is an instance of the standard distinction between density and absolute density. \citet{AdamekElBashirSobralVelebil2001} identify lax epimorphisms with absolutely dense functors and with the connected-factorization condition. The construction above applies that distinction to the question left here; we do not claim a new density theorem. It also explains the failed proof: representables generate a presheaf category by colimits, whereas the proposed right-adjoint argument would require an appropriate cogeneration statement, which is absent.
\end{dcomment}


\textbf{Corollaire 4.2.} \textit{Pour que $u$ soit $W_0$-connexe, il} suffit \textit{qu'il soit une $W_0$-équivalence universelle, i.e. (cf. 3.1) que les fibres soient \textup{0}-connexes, et que pour toute flèche non identique dans $I$, définissant $\bDelta^1 \longrightarrow I$, $I' \times_I \bDelta^1$ soit \textup{0}-connexe.}

4.3. Cette condition de $W_0$-équivalence universelle est cependant \textit{strictement} plus forte que la $W_0$-connexité, comme on le voit sur l'exemple où $u$ est une \textit{équivalence de catégories} non surjective sur les objets (p.ex. l'inclusion d'un sommet dans un groupoïde 1-connexe ayant plus d'un objet). Plus généralement, tout foncteur de localisation est $W_0$-connexe, alors que ses fibres ne sont pas forcément connexes, et a fortiori ce n'est pas une $W_0$-équivalence universelle.

Cependant, si tout endomorphisme $p$ d'un $i$ dans $I$, tel que $p^2 = p$ (projecteur) est l'identité, et si \msp{page 57}$I' \longrightarrow I$ est `transportable', p.ex. si dans $I$ tout isomorphisme est une identité, alors $u$ $W_0$-connexe implique que les fibres de $u$ sont 0-connexes, comme on voit en exprimant d$''$ pour $\alpha$ flèche identique. Les conditions envisagées sur $I$ sont satisfaites si $I$ est ordonnée. Mais même si $I'$ et $I$ sont toutes deux ordonnées, $u$ peut être $W_0$-connexe, sans être une $W_0$-équivalence universelle. Voici un exemple avec $I = \bDelta^2$ :
\[
\begin{array}{r@{\;}l@{\quad}c}
I' = & \left\{\begin{tikzcd}[ampersand replacement=\&]
 \& 1' \arrow[r] \arrow[d] \& 2 \\
0 \arrow[r] \& 1 \&
\end{tikzcd}\right\} & \\[1ex]
 & & \Big\downarrow \\[1ex]
I = & \left\{\begin{tikzcd}[ampersand replacement=\&]
0 \arrow[r] \& 1 \arrow[r] \& 2
\end{tikzcd}\right\}. &
\end{array}
\]
% foot: Version 8 juillet 2001

\printedpage{36}
% head: A. Grothendieck : Dérivateurs, Ch. II, Édition des Universités de Montpellier II et Paris VII
Les catégories bilocalisées de $I$ sont les trois sommets $\{0\}$, $\{1\}$, $\{2\}$, et les trois segments $[0,1]$, $[1,2]$ et $[0,2] = I$, les images inverses sont respectivement
\[
\bullet\ ,\quad
\begin{tikzcd}[row sep=normal] {} \arrow[d] \\ {} \end{tikzcd}\ ,\quad
\bullet\ ,\quad
\begin{tikzcd}[row sep=normal] & {} \arrow[d] \\ {} \arrow[r] & {} \end{tikzcd}\ ,\quad
\begin{tikzcd}[row sep=normal] {} \arrow[r] \arrow[d] & {} \\ {} & \end{tikzcd}\ ,\quad
I'\quad ,
\]
elles sont toutes 0-connexes, i.e. $W_0$-asphériques, et même contractiles donc $W_\omega$-asphériques (donc $u$ est même $W_\omega$-connexe, i.e. $W_\omega$-biasphérique). Mais ce n'est pas une $W_0$-équivalence universelle, car l'image inverse de $I'$ par $\bDelta^1 \xrightarrow{\partial} \bDelta^2$, $\partial(0) = 0$, $\partial(1) = 2$, \msp{page 58}est la catégorie disconnexe $\varepsilon = \{\bullet\ \ \bullet\}$.

J'avais espéré un moment que les foncteurs $W_0$-connexes $I' \xrightarrow{u} I$ seraient ceux qui se factorisent en $I' \xrightarrow{u_0} I'' \xrightarrow{v} I$, avec $u_0$ une équivalence, et $v$ une $W_0$-équivalence universelle. Mais l'exemple précédent le met en défaut, avec $u$ même $W_\omega$-connexe. En effet, $I''$ serait déduit d'un ensemble préordonné, et $v$ se factoriserait par l'ensemble ordonné correspondant, savoir $I'$, donc serait de la forme
\[
I'' \longrightarrow I' \xrightarrow{u} I,
\]
où $I'' \longrightarrow I'$ est un passage au quotient, et l'image inverse de $\bDelta^1 \xrightarrow{\partial} I$ serait encore disconnexe (un $\pi_0$ de cardinal 2). Donc c'est vraiment sans espoir.

\dgrblock{dgrb119}{A contractible replacement which loses an edge}
\begin{dcomment}{examples}{p58:Donc c'est vraiment sans espoir.}{A contractible replacement which loses an edge}
Grothendieck next tests whether full faithfulness might conceal a stronger, universally preserved equivalence. His four-object example separates these requirements even for finite posets. It deserves to be read as a complete computation, since the obstruction lies in a strict pullback which the homotopy type of the total category cannot detect.

Write $A$ for the poset with precisely the nonidentity arrows
\[
 0\longrightarrow 1,\qquad 1'\longrightarrow1,
 \qquad 1'\longrightarrow2,
\]
and let $u:A\to[2]$ identify $1'$ with $1$. For the three identity arrows of $[2]$, the factorization categories are respectively the point $\{0\}$, the interval $1'\to1$, and the point $\{2\}$. For $0\to1$ and $1\to2$, they are the full subposets $0\to1\leftarrow1'$ and $1\leftarrow1'\to2$. The former has a terminal object and the latter an initial object. For $0\to2$, the factorization category is all of $A$: its nerve has no nondegenerate simplices above dimension one, and its realization is the tree with successive vertices $0,1,1',2$. All six nerves are contractible. Thus this is a homotopical epimorphism, and restriction is fully faithful for every derivator; the later criterion is \citet[Theorem 6.10]{GrothStovicek2016Trees}.

Now pull back along the edge $[1]\to[2]$ selecting $0\to2$. The only objects left are $0$ and $2$, and $A$ contains no arrow between them. The pullback is the discrete two-point category, whose map to $[1]$ is not a $W_0$-equivalence. Full faithfulness has survived every factorization test and failed this base-change test.

The same computation rules out the proposed factorization through an equivalence followed by a universal equivalence. Suppose $A\xrightarrow{v}C\xrightarrow{w}[2]$ were such a factorization, with $v$ an equivalence. Since $A$ is a poset, $C$ is a preorder, and its poset of isomorphism classes identifies with $A$. Any functor from $C$ to $[2]$ is constant on isomorphism classes. The inverse image of the edge $0\to2$ consequently has two nonempty components, corresponding to the objects $0$ and $2$ of $A$, and no morphism joining them. It cannot be connected. This argument uses the skeletal quotient only to count components; it does not assume that an arbitrary strict pullback preserves an equivalence.
\end{dcomment}


Cependant

\textbf{Corollaire 4.3.} \textit{Supposons que $u$ soit $W_0$-lisse, ou $W_0$-propre, ou une $W_0$-fibration faible. Alors $u$ est $W_0$-connexe si et seulement si c'est une $W_0$-équivalence universelle, i.e. si et seulement si a) ses fibres sont \textup{0}-connexes et b) $I' \times_I \bDelta^1$} \texttt{[est]} \textit{\textup{0}-connexe pour tout $\bDelta^1 \longrightarrow I$ (correspondant à une flèche non identique).}

\msp{page 59}Dans les trois cas on sait que $u$ est une $W_0$-équivalence universelle si et seulement si ses fibres sont $W_0$-connexes. Or la $W_0$-lissité implique que $\forall i \in \Ob I$, l'inclusion
\[
I'_i \longrightarrow i\backslash I'
\]
est une $W_0$-équivalence, d'autre part $i\backslash I' \longrightarrow i\backslash I$ étant $W_0$-asphérique, donc une $W_0$-équivalence par hypothèse, $i\backslash I'$ est 0-connexe, donc $I'_i$ aussi. Le cas $u$ $W_0$-propre est dual. Dans le cas \texttt{[où]} $u$ \texttt{[est]} une $W_0$-fibration faible, on sait encore que $I'_i \longrightarrow i\backslash I'$ est une $W_0$-équivalence, et on conclut comme dans le cas lisse (ou, au choix, via $I'_i \longrightarrow I'/i$).

Mais l'idée me vient que les foncteurs $W_0$-connexes pourraient être exactement les foncteurs de localisation ? Il revient au même de dire qu'un foncteur $u : I' \longrightarrow I$ \textit{conservatif} (i.e. $u(\alpha)$ isomorphisme $\Longrightarrow$ $\alpha$ isomorphisme) qui est $W_0$-connexe, est une équivalence. Mais c'est faux, comme on voit en prenant $I = e$, $I'$ un groupoïde 0-connexe et non 1-connexe : $I' \longrightarrow e$ est $W_0$-connexe (et même une $W_0$-équivalence universelle), mais ce n'est pas une équivalence (et pas un foncteur de localisation).
% foot: Version 8 juillet 2001

\dgrblock{dgrb120}{The quantifier in homotopical epimorphisms}
\begin{dcomment}{later}{p59:mais ce n'est pas une équivalence (et pas un foncteur de localisation).}{The quantifier in homotopical epimorphisms}
Grothendieck has separated three operations that are easy to conflate: localizing a category, embedding its diagram theory by restriction, and preserving an equivalence after every strict base change. The example $BG\to e$ at the end of this page separates the first two in dimension zero. Constant $G$-sets form a full subcategory of $G$-sets, so restriction from sets is fully faithful; yet a localization of a groupoid at any collection of its arrows is equivalent to that groupoid, since all those arrows are already invertible. For $G\ne1$, its projection to $e$ is therefore not a localization.

The modern term \emph{homotopical epimorphism} adds a quantifier absent from the fixed-$\D$ problem about to begin: restriction must be fully faithful for \emph{every} derivator. \citet[Definition 6.1 and Theorem 6.10]{GrothStovicek2016Trees}, in the arXiv v3 numbering, characterize this condition by weak contractibility of all the factorization categories. Their Remark 6.3(ii) expressly distinguishes full faithfulness for one specified derivator. The space-valued theory detects the universal condition; a particular theory of coefficients can forget more.

This development follows the path from Heller's homotopy theories and Cisinski's comparison results to a calculus of homotopy exact squares. For the present pages, the decisive feature is concrete: connected components in 4.1 are replaced by weakly contractible nerves. The application in tilting theory is equally concrete: a homotopical epimorphism identifies a category of coherent diagrams with a full subcategory of another one, described by unit or counit conditions. No claim that every such functor is an ordinary localization is required.
\end{dcomment}


\printedpage{37}
% head: A. Grothendieck : Dérivateurs, Ch. II, Édition des Universités de Montpellier II et Paris VII
\msp{page 60}Je voudrais généraliser la notion de flèche $W_0$-biasphérique au cas d'un dérivateur $\bD$ quelconque.

\dgrblock{dgrb121}{Recovering a global diagram from its slices}
\begin{dcomment}{exposition}{p60:Je voudrais généraliser la notion}{Recovering a global diagram from its slices}
Grothendieck now wants the factorization criterion to work with an arbitrary theory of coefficients. The difficulty is no longer the adjunction: a right adjoint is fully faithful exactly when its counit is invertible, and a left adjoint exactly when its unit is invertible. The issue is whether the tests visible from global diagrams detect every diagram on a slice. Proposition 4.4 isolates a sufficient detection condition instead of assuming such a property implicitly.

Put $q=u/i:A/i\to B/i$, $\alpha:B/i\to B$, and let $\xi$ lie in $\D(e)$. With the variance of the manuscript, the comparison of homotopy limits on $B/i$ and $A/i$ is induced, after applying $\operatorname{Hom}(\xi,-)$, by the counit
\[
 c_\xi:q_!q^*(\xi_{B/i})\longrightarrow\xi_{B/i}
 \qquad\hbox{in }\D(B/i).
\]
The subscripts in the edition's $(***)$ should be $B/i$ on the object to which $q^*$ is applied and on the target. The printed $I'/i$ in both places makes that display ill-typed; the argument itself determines the correction.

If all $c_\xi$ are invertible, then for every $G\in\D(B/i)$ the maps $\operatorname{Hom}(c_\xi,G)$ are bijections. By adjunction and Yoneda, the induced comparison of homotopy limits is invertible. Taking $G=\alpha^*F$ and using the pointwise right Kan extension formula proves that $F\to u_*u^*F$ is invertible at $i$. As this holds for every $i$, the derivator axiom detecting isomorphisms pointwise gives full faithfulness. This proves the sufficient direction without requiring every slice diagram to extend globally.

For the converse, full faithfulness supplies only
\[
 \operatorname{Hom}_{\D(B/i)}(c_\xi,\alpha^*F)
 \quad\hbox{bijective for every }F\in\D(B).
\]
Adjunction turns this into bijectivity of $\operatorname{Hom}_{\D(B)}(\alpha_!c_\xi,F)$ for every $F$, so Yoneda says exactly that $\alpha_!c_\xi$ is invertible. Conservativity of $\alpha_!$ then proves that $c_\xi$ is invertible. This is condition $3^o$; the proof needs conservativity only on this family of counits, although the stronger formulation is easier to use.

Condition $2^o$ implies $3^o$ by an elementary adjunction argument. If $L\dashv R$ and $R$ is essentially surjective, then an arrow $h$ for which $Lh$ is invertible induces bijections on Hom into every $Rd$, hence into every object; therefore $h$ is invertible. If $B$ is a preorder, $B/i\to B$ is fully faithful. Its left Kan extension is fully faithful, so the unit $G\to\alpha^*\alpha_!G$ is invertible and $\alpha^*$ is essentially surjective. Thus $1^o$ implies $2^o$, which implies $3^o$.

The hypothesis on coproducts in 4.7 will make $3^o$ available in many coefficient theories, including pointed ones. It is a property of those coefficients, not a formal property of every derivator. The modern axioms used here---pointwise calculation of Kan extensions and pointwise detection of isomorphisms---are recorded in \citet[Definition 1.5]{Groth2013}; their different roles will matter for the parenthetical assertion about faithfulness.
\end{dcomment}


\textbf{Proposition 4.4.} \textit{Soit $\bD$ un dérivateur sur} Diag, \textit{et $u : I' \longrightarrow I$ une flèche dans} Diag. \textit{Supposons que $u_*$ et $u_!$ existent, et se calculent de façon standard, pour tout flèche $u$ dans} Diag. \textit{(En fait, il suffit que ce soit le cas pour $u$ et ses localisées $u/i : I'/i \longrightarrow I/i$, ou $u$ et ses colocalisées $i\backslash u : i\backslash I' \longrightarrow i\backslash I$.) Considérons les conditions :} % [sic] pour tout flèche -> pour toute flèche, gender agreement
\begin{itemize}[leftmargin=4em]
\item[\textit{a)}] $u^* : \bD(I) \longrightarrow \bD(I')$ \texttt{[est]} \textit{pleinement fidèle.}
\item[\textit{a$'$)}] \textit{$\forall$ $F$ dans $\bD(I)$, la flèche}
\[
u_!u^*F \longrightarrow F
\]
\textit{est un isomorphisme.}
\item[\textit{a$''$)}] \textit{$\forall$ $F$ dans $\bD(I)$, la flèche}
\[
u_*u^*F \longleftarrow F
\]
\textit{est un isomorphisme.}
\item[\textit{b)}] \textit{$\forall$ $i$ dans $I$, la flèche $i\backslash I' \longrightarrow i\backslash I$ est $\bD$-asphérique.}
\item[\textit{b$'$)}] \textit{$\forall$ $i$ dans $I$, la flèche $I'/i \longrightarrow I/i$ est $\bD$-coasphérique.}
\item[\textit{b$''$)}] \textit{$\forall$ $\alpha : i_0 \longrightarrow i_1$ dans $I$, la catégorie $i_0\backslash\underset{\alpha}{I'}/i_1$ est $\bD$-asphérique.}
\end{itemize}
\textit{On a alors les implications}
\[
(b \Longleftrightarrow b' \Longleftrightarrow b'') \Longrightarrow (a \Longleftrightarrow a' \Longleftrightarrow a'').
\]
\textit{On a $a \Longrightarrow b$ (i.e. toutes les conditions licites envisagées sont équivalentes)} \msp{page 61}\textit{dans chacun des cas suivants.}
\begin{itemize}[leftmargin=4em]
\item[$1^o$] \textit{$I$ est préordonné.}
\item[$2^o$] \textit{$\forall$ $i$ dans $I$, $\alpha_i^* : \bD(I) \longrightarrow \bD(I/i)$ est essentiellement surjectif.}
\item[$3^o$] \textit{$\forall$ $i$ dans $I$, considérant $\alpha_i : I/i \longrightarrow I$, le foncteur $\alpha_{i,!} : \bD(I/i) \longrightarrow \bD(I)$ est conservatif.}
\item[$4^o$] $\mathcal{D} = \mathrm{HOT}_{W_0} \approx \bD_{\mathrm{Ens}}$.
\end{itemize}
\textsc{Démonstration}. L'équivalence de a, a$'$, a$''$ est claire, soit $A$. De même l'equivalence de b, b$'$, b$''$, soit $B$. % [sic] l'equivalence -> l'équivalence, accent missing

Prouvons
\[
B \Longrightarrow A
\]
en supposant p.ex. l'existence de $u_*$. On admet tacitement, avec l'existence de $u_*$ : $\bD(I') \longrightarrow \bD(I)$, le mode de calcul des fibres de $u_*(-)$. Donc $A$, via a$''$, équivaut à
\[
H^{\bullet}_{\bD}(I'/i, u^*F) \xleftarrow{\mathit{iso}} H^{\bullet}_{\bD}(I/i, F) \simeq F(i)
\]
% foot: Version 8 juillet 2001
\printedpage{38}
% head: A. Grothendieck : Dérivateurs, Ch. II, Édition des Universités de Montpellier II et Paris VII
pour tout $F$ dans $\bD(I)$, tout $i$ dans $I$.

Soit
\[
\alpha_i : I/i \longrightarrow F, \qquad F_i = \alpha_i^*(F), % [sic] I/i \longrightarrow F -> I/i \longrightarrow I, alpha_i is defined as I/i -> I in 3^o
\]
on doit donc exprimer que pour tout $F_i$ provenant d'un $F$, on a
\[
H^{\bullet}_{\bD}(I/i, F_i) \xrightarrow{\mathrm{iso}} H^{\bullet}_{\bD}(I'/i, (u/i)^*(F_i)). \tag*{$(\ast)$}
\]
Or sous réserve d'existence de $(u/i)_!$, dire que $(\ast)$ est \texttt{[un]} isomorphisme pour \textit{tout} $F_i$ de $\bD(I/i)$, signifie que
\[
u/i : I'/i \longrightarrow I/i
\]
\msp{page 62}est $\bD$-coasphérique, d'où b$'$ $\Longrightarrow$ a$''$. Il est clair que l'implication inverse est valable si $\alpha_i^*$ est essentiellement surjectif, c'est le cas $2^o$. Or $2^o$ est satisfait si $I$ est préordonné, d'où le cas $1^o$. Le cas $4^o$ a déjà été traité (4.1). Il faut encore traiter le cas $3^o$, qui contient d'ailleurs le cas $2^o$, donc le cas $1^o$ (et aussi le cas $4^o$, cf. plus bas).

Il faut exprimer l'isomorphie de la flèche $(\ast)$ pour $F_i$ de la forme $\alpha_i^*(F)$. Soit $\xi$ \texttt{[un]} objet de $\mathcal{A} = \mathcal{D}(e)$, il faut exprimer l'isomorphie de $\Hom(\xi, -)$ dans les deux membres de $(\ast)$, i.e. de
\[
\begin{aligned}
\Hom(I/i, F_i) \quad &\longrightarrow \quad \Hom(\xi'_{I'/i}, (u/i)^*(F_i)) \\
&\overset{\text{adj.}}{\simeq} \quad \Hom((u/i)_!(\xi'_{I'/i}), F_i),
\end{aligned}
\tag*{$(\ast\ast)$}
\]
\texttt{[plutôt} $\xi_{I'/i}$\texttt{]} flèche qui se déduit de
\[
(u/i)_!(u/i)^*(\xi_{I'/i}) \xrightarrow{\text{adj.}} \xi_{I'/i} \tag*{$(\ast\ast\ast)$}
\]
en prenant le Hom dans $F_i$. La condition $u/i$ coasphérique (i.e. b$'$) signifie que $(\ast\ast\ast)$ est un isomorphisme pour tout $\xi$ dans $\mathcal{A}$. \textit{Suffit}-il pour ceci que son Hom dans un $F_i$ soit bijectif, quand on restreint $F_i$ dans $\bD(I/i)$ à être dans l'image de $\bD(I)$ par $\alpha_i^*$ ? Or cette flèche $(\ast\ast)$, pour $F_i = \alpha_i^*F$, s'explicite par la formule d'adjonction relative \msp{page 63}au foncteur $\alpha_i$, comme
\[
\Hom(\alpha_{i,!}(\xi_{I/i}), F) \longrightarrow \Hom(\alpha_{i,!}(\ldots), F).
\]
Dire que c'est un isomorphisme pour tout $F$, signifie donc \textit{exactement} que la transformée de la flèche $(\ast\ast\ast)$ par $\alpha_!$ est un isomorphisme. Pour en conclure que $(\ast\ast\ast)$ elle-même est \texttt{[un]} isomorphisme, la condition naturelle à imposer est donc que $\alpha_{i,!} : \bD(I/i) \longrightarrow \bD(I)$ soit conservatif. C'est le cas $3^o$. Cela achève la démonstration de 4.4.

4.5. Montrons que la condition $3^o$ sur $I$ n'est pas toujours satisfaite, dans le cas de $\bD = (\mathrm{HOT}_{W_0})^{\mathrm{opp}} \approx (\bD_{\mathrm{Ens}})^{\mathrm{opp}}$, bien que dans ce cas, il soit pourtant vrai que pour toute $I'$ sur
% foot: Version 8 juillet 2001

% ---- from Der-ChII-234_p041-045.tex
\printedpage{39}
% head: A. Grothendieck : Dérivateurs, Ch. II, Édition des Universités de Montpellier II et Paris VII

$I$ satisfaisant la condition a$''$ (autoduale), la condition b$'$ (elle aussi autoduale) est aussi satisfaite (par cas $4^o$ ci-dessus). Je dis donc qu'il existe $I$ et $i$ dans $I$, tels que
\[
\beta_{i,*} : (i\backslash I)^\wedge \longrightarrow I^\wedge
\]
ne soit pas conservatif. (\textbf{NB} Par contre $\alpha_{i,!} : (I/i)^\wedge \longrightarrow I^\wedge$ est toujours conservatif, ce qui redonne le cas $4^o$ comme cas particulier de $3^o$). Comme $\beta_i : i\backslash I \longrightarrow I$ est $W_0$-propre (et même $W_\omega$-propre !), $\beta_{i,*}$ `se calcule fibre par fibre', donc on a
\msp{page 64}
\[
\beta_{i,*}(G)(j) \simeq \prod_{\alpha\in\mathrm{Hom}(i,j)} G(\alpha),
\]
la fibre en question en $j$ étant discrète. Prenons
\[
I = 0 \underset{\beta}{\overset{\alpha}{\rightrightarrows}} 1, \qquad\quad \text{avec } \alpha \neq \beta,
\]
alors $0\backslash I$ est $\Psi = {}$\begin{tikzcd}[ampersand replacement=\&,row sep=0.1em, column sep=1.2em, baseline=(\tikzcdmatrixname-2-1.base)]
 \& 1 \\
0 \ar[ur, "\alpha'"] \ar[dr, "\beta'"'] \& \\
 \& 1'
\end{tikzcd}, et $\Psi^\wedge$ est la catégorie des diagrammes d'ensembles $G = {}$\begin{tikzcd}[ampersand replacement=\&,row sep=0.1em, column sep=1.2em, baseline=(\tikzcdmatrixname-2-1.base)]
 \& G_1 \ar[dl, "\lambda"'] \\
G_0 \& \\
 \& G_{1'} \ar[ul, "\mu"]
\end{tikzcd}. Le foncteur $\beta_*$ déduit de $\beta : i\backslash I \longrightarrow I$ transforme ce diagramme en
\[
G_1 \times G_{1'} \underset{\lambda\,\mathrm{pr}_2}{\overset{\lambda\,\mathrm{pr}_1}{\rightrightarrows}} G_0.
\]
\texttt{[plutôt} $\mu\,\mathrm{pr}_2$\texttt{]}. Si on se borne aux diagrammes $G$ pour lesquels $G_{1'} = \emptyset$, $\beta_*$ sur ceux-ci donne
\[
\emptyset \rightrightarrows G_0,
\]
donc la restriction de $\beta_*$ à cette sous-catégorie (équivalente à $(\boldsymbol{\Delta}^1)^\wedge$) n'est pas fidèle, ni conservative, ne tenant compte que de la composante $G_0$, et non pas de $G_1$.

\dgrblock{dgrb122}{A product erases a component}
\begin{dcomment}{examples}{p64:ne tenant compte que de la composante}{A product erases a component}
Grothendieck's example explains why the conservative-slice hypothesis cannot simply be imposed on one chosen side of every derivator. A right Kan extension over a discrete fibre is a product, and the empty set can erase every other factor. This is an obstruction in ordinary presheaf theory, before any passage to homotopy categories.

Let $B$ have two distinct arrows $\alpha,\beta:0\rightrightarrows1$, and let $p:0\backslash B\to B$ be the projection. Its domain has the shape $a\to b$, $a\to c$, with $p(a)=0$ and $p(b)=p(c)=1$. A presheaf $G$ is a diagram $G_b\xrightarrow{\lambda}G_a\xleftarrow{\mu}G_c$, and
\[
 (p_*G)(0)=G_a,\qquad
 (p_*G)(1)=G_b\times G_c.
\]
The two restriction maps are $\lambda\operatorname{pr}_1$ and $\mu\operatorname{pr}_2$. Take $G_a=1$, $G_b=\varnothing$, $G_c=\varnothing$, and $H_a=1$, $H_b=1$, $H_c=\varnothing$. The evident map $G\to H$ is not an isomorphism, but $p_*G\to p_*H$ is the identity of the presheaf $\varnothing\rightrightarrows1$. To see nonfaithfulness as well, take $G_b=1$, $H_b=2$ while keeping the $a$-values equal to $1$ and the $c$-values empty; the two distinct maps on the $b$-component have identical direct images.

For the other variance, $B/i\to B$ is a discrete fibration and its left Kan extension is a disjoint union over each fibre. In sets, a map between two such unions that preserves the indexing is bijective exactly when each component is bijective. This proves the conservativity used in 4.4 for set coefficients. Dualizing the derivator exchanges these two mechanisms, which is why failure of the particular sufficient condition does not refute the equivalence of full faithfulness and biasphericity for sets or their opposite theory.
\end{dcomment}


\msp{page 63bis}

\textbf{Definition 4.6.} \textit{On dit que $u : I' \longrightarrow I$ dans} Diag \textit{est} $\D$-biasphérique \textit{si elle satisfait les conditions équivalentes b, b$'$, b$''$ de 4.4,} $\D$-connexe \textit{si elle satisfait la condition a, ou encore (sous réserve d'existence de $u_!$ resp. de $u_*$) les conditions équivalentes a$'$, a$''$.} % [sic] Definition -> Définition, accent missing in print

Ainsi, $\D$-biasphérique implique $\D$-connexe, et la réciproque est vraie quand $I$ satisfait (par rapport à $\D$) la condition $3^o$ de 4.4 (donc aussi, si elle satisfait l'une des conditions plus faibles $2^o$, $1^o$), ou la condition duale de celle-ci (vu que les notions $\D$-biasphérique, $\D$-connexe sont autoduales). À cet égard, notons le

\textbf{Corollaire 4.7.} \textit{Supposons que $\mathcal{A} = \D(e)$ satisfasse la condition suivante : le foncteur}
\[
\left\{
\begin{array}{lcccl}
\mathcal{A} & \times & \mathcal{A} & \longrightarrow & \mathcal{A} \\[1ex]
(\xi & , & \eta) & \longrightarrow & \xi \amalg \eta
\end{array}
\right.
\]
\textit{est conservatif (resp. fidèle). Alors pour toute flèche étale $u : J \longrightarrow I$ dans} Diag \textit{(et en particulier pour $I/i \longrightarrow I$), le foncteur $u_! : \D(J) \longrightarrow \D(I)$ est conservatif (resp. fidèle)} (\footnote{\textbf{NB} Cette condition est satisfaite si $\forall\, \xi, \eta \in \mathcal{A}$ on a $\mathrm{Hom}(\xi,\eta) \neq \emptyset$, (car on se ramène à ceci : si $u_1$,
% footnote continues at the foot of p. 40
$u_2$ sont deux flèches dans Ens entre ensembles \textit{non vides,} et si $u_1 \times u_2$ est \texttt{[un]} isomorphisme, $u_1$, $u_2$ sont \texttt{[des]} isomorphismes). Cela s'applique en particulier si $\mathcal{A}$ admet un objet neutre, p.ex. si $\mathcal{A}$ est additive.}). \textit{Si cette condition, ou la condition duale,}
\printedpage{40}
% head: A. Grothendieck : Dérivateurs, Ch. II, Édition des Universités de Montpellier II et Paris VII
\msp{page 64bis}
\textit{impliquant les produits $\xi\times\eta$ ($\xi$, $\eta$ dans $\mathcal{A}$), est satisfaite, alors pour toute flèche $u : I' \longrightarrow I$ dans} Diag, \textit{$u$ est $\D$-connexe si et seulement si $u$ est $\D$-biasphérique.}

\textsc{Démonstration.} L'hypothèse de conservativité, ou de fidélité se transporte aux sommes indexées par un ensemble quelconque (sous réserve d'existence), comme on voit en écrivant $\coprod\limits_{\lambda} \xi_\lambda = \xi_{\lambda_0} \amalg \coprod\limits_{\lambda\neq\lambda_0} \xi_\lambda$. L'assertion faite sur $u_!$ résulte alors du calcul de $u_!$ fibre par fibre (qui est valable vu que $u$ est étale donc lisse), en tenant compte que les fibres sont discrètes.

\dgrblock{dgrb123}{Faithfulness does not follow from the fibre calculation}
\begin{dcomment}{development}{p64bis:en tenant compte que les fibres sont discrètes.}{Faithfulness does not follow from the fibre calculation}
Grothendieck uses disjoint fibres to promote a property of coproducts in $\D(e)$ to a property of $u_!$ on coherent diagrams. For conservativity, the argument is complete: an invertible $u_!h$ gives an invertible component $h_j$ at every $j$, and the derivator axiom then makes $h$ invertible. For faithfulness, pointwise evaluation does not provide the analogous conclusion. A small derived example shows that the parenthetical assertion is false.

\textbf{An étale direct image killing a nonzero extension.}\ours{}
Let $k$ be a field and take $\D(K)=D(\operatorname{Vect}_k^{K^{\mathrm{op}}})$. There is a discrete fibration $p:J\to I$ for which $p_!$ is not faithful, although binary coproducts in $\D(e)=D(k)$ are faithful.

\emph{Proof.} Take $I=(0\to1)$ and let $J$ have objects $a,b,c$ and the two nonidentity arrows $a\to b$, $a\to c$. Send $a$ to $0$ and both $b,c$ to $1$. The unique lifting condition for arrows with a fixed target verifies that $p$ is a discrete fibration. A presheaf on $J$ is written $(B\to A\leftarrow C)$. Its left Kan extension is
\[
       p_!(B\xrightarrow{f}A\xleftarrow{g}C)
       =(B\oplus C\xrightarrow{(f,g)}A).
\]
This functor is exact, since it is computed by finite direct sums at each vertex. Its derived functor is consequently given by the same formula on complexes.

Consider the three representations
\[
 M=(k\to0\leftarrow0),\qquad
 N=(0\to k\xleftarrow{\mathrm{id}}k),\qquad
 E=(k\xrightarrow{\mathrm{id}}k\xleftarrow{\mathrm{id}}k).
\]
There is a short exact sequence $0\to N\to E\to M\to0$: the inclusion is the identity at $a,c$ and zero at $b$, and the quotient is the identity at $b$ and zero elsewhere. It does not split. A section $M\to E$ would have to be the identity at $b$ and zero at $a$; commutativity along $a\to b$ for the contravariant diagrams would then assert $\mathrm{id}_k=0$. Therefore its class
\[
       \delta\in\operatorname{Ext}^1_{\operatorname{Vect}_k^{J^{\mathrm{op}}}}(M,N)
       =\operatorname{Hom}_{\D(J)}(M,N[1])
\]
is nonzero. The equality is the usual identification of extension classes with degree-one derived morphisms; here all objects are representations of a finite acyclic quiver.

After applying $p_!$, the middle representation is $(k\oplus k\xrightarrow{(1,1)}k)$ and the quotient is $(k\to0)$. The map at the source sending $t$ to $(t,-t)$, together with the zero map at the target, gives a section of that quotient. Hence the image extension splits and $p_!\delta=0$. Binary direct sum on $D(k)$ is faithful because each summand of a morphism is recovered by composing with the canonical inclusion and projection. This proves the assertion. \hfill$\square$

The example concerns only ``resp. fidèle''. It leaves the conservative conclusion, and therefore the final equivalence in 4.7, intact. For represented derivators the evaluation functors are jointly faithful and the parenthetical argument does work; it also works under an explicit joint-faithfulness hypothesis on $\D(J)\to\D(e)^{\operatorname{Ob}J}$. We checked the distinction between coherent and incoherent diagram categories in \citet[Introduction and Definition 1.5]{Groth2013} and the treatment of homotopical epimorphisms in \citet[\S6]{GrothStovicek2016Trees}; we did not find this particular counterexample to 4.7 there. The calculation is included as a repair of the statement, without a claim of priority.
\end{dcomment}


Je ne connais pas d'exemple d'un couple $(\D, u)$ d'un dérivateur $\D$ et d'une flèche $u$ dans Diag, pour lesquels $u$ soit $\D$-connexe sans être $\D$-asphérique. On a donné des conditions sur $\D$, qui assurent équivalence des deux notions. Si $\D$ est de la forme $\D_{\mathcal{M}}$, et si $\mathcal{M}$ satisfait la condition dite dans 4.7 ou sa duale, alors par 4.7 c'est OK. Donc un contrexemple avec un $\D_{\mathcal{M}}$ devrait en tenir compte. Je ne vais pas m'obstiner pour tirer au clair ce qui en est.

\dgrblock{dgrb124}{A connected restriction with an empty factorization category}
\begin{dcomment}{development}{p64bis:Je ne vais pas m'obstiner pour tirer au clair ce qui en est.}{A connected restriction with an empty factorization category}
Grothendieck asks whether the sufficient conditions just proved leave a genuine distinction between $\D$-connectedness and $\D$-biasphericity. The edition says ``sans être $\D$-asphérique'', but that literal question has an immediate negative answer: full faithfulness makes the unit invertible on constant objects, hence implies asphericity. The surrounding discussion and the next diagram concern \emph{biasphericity}. A represented derivator supplies the requested separating example.

\textbf{A two-valued coefficient theory.}\ours{}
Let $L=\{0<1\}$ be the complete lattice, regarded as a category, and let $\D_L(K)=\operatorname{Fun}(K^{\mathrm{op}},L)$. Let $B$ be the free category on the directed cycle $x\xrightarrow{s}y\xrightarrow{t}x$, and let $u:e\to B$ select $y$. Then $u$ is $\D_L$-connected and is not $\D_L$-biaspherical.

\emph{Proof.} The complete lattice $L$ has every small limit and colimit, given respectively by infima and suprema. Thus its represented prederivator is a derivator: restriction has both pointwise Kan extensions, disjoint unions of indexing categories give products of functor categories, and natural transformations are invertible exactly when their components are. It is also strong, since currying gives the required coherent-arrow identifications. These are the standard represented examples recalled in \citet[Examples 1.2(1), 1.6 and Definitions 1.5, 1.8]{Groth2013}, with the manuscript's contravariant diagram convention.

An object $F$ of $\D_L(B)$ gives inequalities $F(y)\leq F(x)$ and $F(x)\leq F(y)$. Antisymmetry forces equality. All arrows of $B$ then have the unique endomorphism of that value as image. Every diagram is therefore constant, and a natural transformation between the two possible diagrams is precisely an inequality in $L$. Evaluation at $y$ identifies $u^*:\D_L(B)\to L$ with an isomorphism of categories. In particular it is fully faithful.

For the identity of the omitted object, however,
\[
        \operatorname{Fact}_u(\mathrm{id}_x)=\varnothing.
\]
Indeed, a factorization through $y$ would be a path $x\to y$ followed by a path $y\to x$. Both have positive length, and in a free category their composite cannot be the length-zero identity. The empty category is not $\D_L$-aspherical: its homotopy colimit of the constant value $1$ is the initial object $0$, and the comparison $0\to1$ is not invertible. Equivalently, $L\to\D_L(\varnothing)=e$ is not fully faithful. Condition 4.4(b$''$) fails. \hfill$\square$

Both potential escapes supplied by 4.7 genuinely fail here. Coproduct is maximum, so the map $(0,1)\to(1,1)$ becomes an identity after taking coproducts although its first component is not invertible. Product is minimum, so $(0,0)\to(1,0)$ gives the dual failure. The counterexample therefore does not contradict the conservative-coproduct theorem, nor the preorder-base case: the indexing category $B$ is a free category with nontrivial endomorphisms, not a preorder.

The example isolates the information lost by the coefficient theory. It sees mutual reachability of objects but cannot distinguish the identity from a nonidentity cycle; the factorization condition still distinguishes them. This coefficient theory is also the classical derivator of truth values, or homotopy $(-1)$-types; see \citet[p.~33]{Shulman2023Setoids}. We checked that discussion, the density and lax-epimorphism literature cited above, the fixed-derivator qualification in \citet[Remark 6.3(ii)]{GrothStovicek2016Trees}, and searches for biasphericity in the derivator literature. We could not find this counterexample there. The claim here is the explicit proposition and proof, rather than an assertion that it has no earlier occurrence.
\end{dcomment}


\msp{page 65}

\textbf{Corollaire 4.8.} \textit{Une flèche $u$ de} Diag \textit{qui est une $\D$-équivalence universelle est $\D$-biasphérique, et \textup{a fortiori} $\D$-connexe. Inversement, si $u$ est soit $\D$-lisse, soit $\D$-propre, soit une $\D$-fibration faible, et si $u$ est $\D$-biasphérique, c'est une $\D$-équivalence universelle.}

\dgrblock{dgrb125}{What the seven classes retain}
\begin{dcomment}{exposition}{p65:c'est une $\D$-équivalence universelle.}{What the seven classes retain}
Grothendieck's diagram organizes the preceding tests by the amount of coefficient information they preserve. Its arrows remain useful, but the two observations below are needed to read the diagram after the preceding counterexample.

First, $\D$-connectedness is stable under composition. If $u:A\to B$ and $v:B\to C$ are $\D$-connected, then
\[
       (vu)^*\simeq u^*v^*
\]
is a composite of fully faithful functors and is fully faithful. Thus the exception printed in the footnote cannot be an exception to closure under composition. This is a formal fact, independent of any Kan-extension existence hypothesis. The footnote's further assertion that all seven classes depend only on $W_\D$ requires an independent argument for the fixed-coefficient connectedness class; it does not follow from the stated factorization criterion, whose converse has just failed.

Second, connectedness always implies both asphericity and coasphericity. The adjunction unit and counit are invertible on every $F\in\D(B)$, hence on constant $F$ in particular. The distinction established above is between this global full faithfulness and the two-sided slice tests defining biasphericity. The label ``équivalence dans tous les cas connus'' records the author's experience at this stage, not an additional theorem.

For universal equivalences the factorization test still gives the implication in 4.8. The category of factorizations of a fixed arrow $\gamma:b_0\to b_1$ through arbitrary objects of $B$ has an initial object $(b_0,\mathrm{id},\gamma)$ and a terminal object $(b_1,\gamma,\mathrm{id})$. Pulling $u$ back to that category produces $\operatorname{Fact}_u(\gamma)$. A universal $W_\D$-equivalence therefore makes this category $\D$-aspherical. Under the smooth, proper, or weak-fibration hypotheses, the fibre-to-comma comparisons give the converse exactly as in 4.3. These hypotheses are what connect strict base changes to the homotopy calculation; the four-object example shows why they cannot simply be omitted.
\end{dcomment}


La démonstration est celle du cas particulier 4.3. On a donc un diagramme d'implications :
\[
\begin{tikzcd}[column sep=small, row sep=large]
\text{Cat-équivalence} \ar[dr, Rightarrow] & \D\text{-équivalence universelle } (W_{\mathrm{univ}}) \ar[d, Rightarrow, "{\text{\footnotesize\begin{tabular}{@{}l@{}}implication stricte ; équivalence\\ dans le cas où $u$ est soit $\mathbf{D}$-propre,\\ soit $\mathbf{D}$-lisse, soit une $\mathbf{D}$-fibration faible.\end{tabular}}}"] & \\
 & \D\text{-biasphérique} \ar[dl, Rightarrow] \ar[d, Rightarrow, "{\text{\footnotesize\begin{tabular}{@{}l@{}}équivalence\\ dans tous les\\ cas connus\end{tabular}}}"] \ar[dr, Rightarrow] & \\
\D\text{-asphérique} \ar[dr, Rightarrow] & \D\text{-connexe} \ar[d, Rightarrow] & \D\text{-coasphérique} \ar[dl, Rightarrow] \\
 & \D\text{-équivalences } (W) &
\end{tikzcd}
\]
(\footnote{\textbf{NB} Ces sept classes de morphismes dans Cat sont stables par composition sauf la notion de $\D$-connexité, toutes ces notions ne dépendent que de l'ensemble $W = W_{\D}$ des $\D$-équivalences.}). Rappelons que pour $u : I' \longrightarrow I$ donné
\begin{enumerate}[leftmargin=*, labelindent=1.5em]
\item[a)] $u$ est $\D$-asphérique si et seulement si $\forall\, F$ \textit{constant} dans $\D(I)$
\[
F \xrightarrow[\sim]{\text{adj.}} u_*u^*F \qquad\quad \text{est \texttt{[un]} isomorphisme.} \tag*{$(\ast)$}
\]
\printedpage{41}
% head: A. Grothendieck : Dérivateurs, Ch. II, Édition des Universités de Montpellier II et Paris VII
\item[b)] $u$ est $\D$-coasphérique si et seulement si $\forall\, F$ \textit{constant} dans $\D(I)$
\[
F \xleftarrow[\sim]{\text{adj.}} u_!u^*F \qquad\quad \text{est \texttt{[un]} isomorphisme.} \tag*{$(\ast\ast)$}
\]
\item[c)] $u$ est $\D$-connexe si et seulement si pour \textit{tout} $F$ dans $\D(I)$ (pas forcément constant, ni localement constant), $(\ast)$ est \texttt{[un]} isomorphisme, ou ce qui revient au même, pour tout $F$ $(\ast\ast)$ est \texttt{[un]} isomorphisme.
\end{enumerate}

\msp{page 66}

Rappelons que la classe $W_{\mathrm{asph}}$ \texttt{[des]} flèches $W$-asphériques et celle, $\mathrm{Liss}_W$, des flèches $W$-lisses, se déterminent mutuellement, ainsi
\[
\left\{
\begin{array}{lcl}
(f : X \longrightarrow Y) \in W_{\mathrm{asph}} & \Longleftrightarrow & \forall\, (Y' \xrightarrow{\ u\ } Y) \in \mathrm{Liss}_W, \\
 & & X' = X \times_Y Y' \longrightarrow X' \text{ est } \in W. % [sic] X' = X \times_Y Y' \longrightarrow X' -> X' = X \times_Y Y' \longrightarrow Y', the target must be Y' (compare (4.12))
\\
(f : Y' \longrightarrow Y) \in \mathrm{Liss}_W & \Longleftrightarrow & \text{Le foncteur } X \longmapsto X' = X \times_Y Y' \\
 & & \text{de } \mathrm{Cat}/Y \text{ dans } \mathrm{Cat}/Y' \text{ transforme} \\
 & & \text{flèches dans } W_{\mathrm{asph}} \text{ dans \texttt{[plutôt en]}} \\
 & & \text{flèches dans } W_{\mathrm{asph}}.
\end{array}
\right.
\tag*{(4.10)}
\]

Pour toute partie $A$ de Fl(Cat) stable par isomorphismes, on peut associer une partie $\mathbf{L}(A)$ de Fl(Cat), par la relation
\[
\begin{array}{rcl}
(u : Y' \longrightarrow Y) \in \mathbf{L}(A) & \Longleftrightarrow & u^* : \mathrm{Cat}/Y \longrightarrow \mathrm{Cat}/Y' \\
 & & \text{transforme flèches} \\
 & & \in A \text{ en flèches } \in A.
\end{array}
\tag*{(4.11)}
\]

On voit que $\mathbf{L}(A)$ contient les isomorphismes, est stable par composition, et est stable aussi par changement de base (comme c'est le cas de $\mathrm{Liss}_W$). Inversement, à une partie $L$ de Fl(Cat), stable par isomorphismes, on associe une partie $\mathbf{A}_W(L)$ par
\[
\begin{array}{rcl}
(f : X \longrightarrow Y) \in \mathbf{A}_W(L) & \Longleftrightarrow & \forall\, (Y' \longrightarrow Y) \text{ dans } L, \\
 & & X' = X \times_Y Y' \longrightarrow Y' \\
 & & \text{est dans } W.
\end{array}
\tag*{(4.12)}
\]

Si $L$ est stable par changement de base, alors $\mathbf{A}_W(L)$ est stable par composition. Il contient les isomorphismes (car $W$ les contient), est contenu dans $W$ si $L$ contient les identités. Mais
\msp{page 67}
aucune raison qu'il soit stable par changement de base quelconque (et en fait, $W_{\mathrm{asph}}$ ne l'est pas).

\dgrblock{dgrb126}{What survives the cancelled argument}
\begin{dcomment}{exposition}{p66}{What survives the cancelled argument}
Grothendieck is seeking a way to recover the classes of smooth and proper maps from the weak equivalences alone. The two tests in (4.10) already suggest an order-theoretic correspondence, but they use different notions of preservation. This distinction explains the cancellation on the next page; it does not make every formula in the cancelled passage false.

To see exactly what is involved, write $S(A)$ for the stabilizer of $A$: a map $p:Y'\to Y$ belongs to $S(A)$ when its pullback functor on slices sends every arrow of $A$ to another arrow of $A$. Write $N(L)=\mathbf A_W(L)$ for the arrows whose pullbacks along maps of $L$ belong to $W$. Here $W$ contains isomorphisms and is closed under composition. Classes are invariant under isomorphism of arrows, and the ambient category has pullbacks. The quantifier over arrows \emph{in the slice} matters: one tests an arrow $X_1\to X_2$ together with any structure map $X_2\to Y$, not merely arrows with codomain exactly $Y$.

The class $S(A)$ contains isomorphisms and is closed under composition and pullback, by the identity and composition laws for pullback functors. If $L$ is closed under pullback, then $N(L)$ is closed under composition: pull back a composite $gf$ along $p\in L$, and factor its pullback as the pullback of $f$ along a pullback of $p$, followed by the pullback of $g$ along $p$. Both factors lie in $W$. If $L$ also contains identities, $N(L)\subseteq W$; its isomorphisms belong automatically. These are the assertions preceding the cancellation, with their implicit hypotheses made explicit.

For $A\subseteq W$, there are also two useful inclusions:
\[
 A\subseteq N(S(A)),\qquad L\subseteq S(N(L)),
\]
where the second requires $L$ to be closed under pullback and composition. For its proof, let $p\in L$ and $f\in N(L)$ be over the target of $p$. The pullback $p^*f$ lies in $W$. Every further test of $p^*f$ along $q\in L$ is a test of $f$ along the composite of $q$ with the appropriate pullback of $p$. That composite is in $L$, so $p^*f\in N(L)$.

For classes $L$ containing identities, these inclusions even prove the two displayed idempotence conclusions on page 67. Since $N$ reverses inclusions,
\[
 NSNS(A)\subseteq NS(A),\qquad NSN(L)\subseteq N(L),
\]
and the reverse inclusions follow from $A\subseteq NS(A)$, applied to $NS(A)$ and to $N(L)$; the identity test ensures $N(L)\subseteq W$. Thus $NS$ is idempotent and $SNSN=SN$. However, $S$ need not reverse inclusions. Consequently these calculations do not supply the order-reversing adjunction that the next section needs. In particular, the stronger justification printed there, $S(A)^\flat=S(A)$, does not follow from them and can fail. The repair in \S6 is to replace $S(A)$ by a transporter into the fixed class $W$. That operation is genuinely order-reversing.

There is a further empty-class detail in the cancelled argument. Formula (4.14) does not explicitly require identities in $L$, whereas (4.12) does not explicitly restrict $N(L)$ to $W$. For $L=\varnothing$, one has $N(L)=\operatorname{Fl}(\mathrm{Cat})$, so the auxiliary assertion $N(L)^\natural=N(L)$ is false when $W$ is proper: its left side is $W_{\mathrm{univ}}$. Section 6 repairs the domain issue by inserting $f\in W$ into (6.6), and later by requiring isomorphisms in $L$.

The idempotence of $SN$ itself nevertheless survives in $\mathrm{Cat}$ for a fundamental localizer and every class $L$ closed under pullback and composition. If every member of $L$ has empty domain, every test is an identity of the empty category, so $N(L)$ is all functors and $SN(L)$ is all functors; applying $SN$ again gives the same class, since universal $W$-equivalences are stable under all pullbacks. Otherwise choose $p\in L$ with nonempty domain and pull it back to an object of its target lying in its image. This gives $C\to e$ in $L$ with $C$ nonempty. Further pullback gives $B\times C\to B$ in $L$ for every $B$. Thus $f\in N(L)$ implies $f\times C\in W$. Choosing an object of $C$ makes $f$ a retract of $f\times C$ in the arrow category, so $f\in W$ by closure under retracts; for fundamental localizers this is \citet[Proposition 4.2.4, p.~183]{Cisinski2006Prefaisceaux}. Hence $N(L)\subseteq W$, and the earlier proof applies. This records the surviving conclusion without endorsing either incorrect auxiliary justification.
\end{dcomment}


(\footnote{Canulé \texttt{[jusqu'à la fin du} \S4\texttt{]}, cf. \S6.}). On a, pour un $A \subseteq W \subseteq \mathrm{Fl(Cat)}$ donné, stable par isomorphismes,
\[
A \subseteq \mathbf{A}_W(\mathbf{L}(A)) \overset{\text{déf}}{=} A^\natural \subseteq W, \qquad \text{\footnotesize\begin{tabular}{@{}l@{}}partie stable par composition,\\ contient \texttt{[les]} isomorphismes,\end{tabular}}
\tag*{(4.13)}
\]
\printedpage{42}
% head: A. Grothendieck : Dérivateurs, Ch. II, Édition des Universités de Montpellier II et Paris VII
et pour toute partie $L$ de Fl(Cat) qui est stable par changement de base, et par composition
\[
L \subseteq \mathbf{L}(\mathbf{A}_W(L)) \overset{\text{déf}}{=} L^\flat, \qquad \text{\footnotesize\begin{tabular}{@{}l@{}}partie de Fl\,Cat stable\\ par composition et par changement\\ de base, contenant les isomorphismes.\end{tabular}}
\tag*{(4.14)}
\]
On a
\[
\left\{
\begin{array}{lcl}
A^{\natural\natural} & = & A^\natural, \quad \text{car } \mathbf{A}_W(L)^\natural = \mathbf{A}_W(L) \\[1ex]
L^{\flat\flat} & = & L^\flat, \quad \text{car } \mathbf{L}(A)^\flat = \mathbf{L}(A).
\end{array}
\right.
\]
Exemples :
\begin{enumerate}[leftmargin=*, labelindent=1em]
\item[$1^o$)] $L = \text{Ét}$, $\mathbf{A}_W(L) = W_{\mathrm{asph}}$, $\mathbf{L}(\mathbf{A}_W(L)) = L^\flat = \mathrm{Liss}_W$.
\item[$2^o$)] $L = \text{Coét}$, $\mathbf{A}_W(L) = W_{\mathrm{coas}}$, $\mathbf{L}(\mathbf{A}_W(L)) = L^\flat = \mathrm{Prop}_W$.
\item[$3^o$)] $A = W$, $\mathbf{L}(A) = \mathrm{Fib}_W$, $A^\natural = W$.
\item[$4^o$)] $L = \mathrm{Fl\,Cat}$, $\mathbf{A}_W(L) = W_{\mathrm{univ}}$, $L^\flat = L$.
\end{enumerate}
\msp{page 68 vide}

\clearpage
\printedpage{43}
% head: A. Grothendieck : Dérivateurs, Ch. II, Édition des Universités de Montpellier II et Paris VII
\msp{page 69}

\dgrblock{dgrb127}{A cancelled justification tested in categories}
\begin{dcomment}{development}{p67}{A cancelled justification tested in categories}
The cancelled lines are best understood by testing the precise identity used in their justification. The following example keeps the ambient category equal to $\mathrm{Cat}$ and uses a fundamental localizer, so the difficulty is not an artefact of passing to an unrelated category.

\textbf{A stabilizer need not be a closed right-hand class.}\ours{}
Let $W$ be the maximal fundamental localizer, consisting of all functors, and let $A=\mathrm{Equ}$ be the class of equivalences of categories. With $S$ and $N$ as above, $S(A)^\flat\ne S(A)$.

Indeed, every pullback functor takes functors to functors, so $N(L)=\operatorname{Fl}(\mathrm{Cat})$ for every $L$. Consequently
\[
 S(A)^\flat=S\bigl(N(S(A))\bigr)
 =S\bigl(\operatorname{Fl}(\mathrm{Cat})\bigr)
 =\operatorname{Fl}(\mathrm{Cat}).
\]
To show that $S(A)$ is smaller, let $J$ be the groupoid with two objects $0,1$ and precisely one morphism between each ordered pair. Each point inclusion $e_i\to J$ is fully faithful and essentially surjective, hence an equivalence. Their strict pullback is empty, because its objects would have to map to both $0$ and $1$. Pulling $e_1\to J$ back along $e_0\to J$ therefore gives $\varnothing\to e_0$, which is not essentially surjective. Thus $e_0\to J\notin S(A)$, proving the claim.

The same calculation shows why the anticipated polarity fails: although $A\subseteq B=\operatorname{Fl}(\mathrm{Cat})$, one has $S(A)\subsetneq S(B)$, whereas an order-reversing correspondence would require the opposite inclusion. This does not contradict the idempotence of either composite proved in the preceding block. For the present $W$ those composites are particularly degenerate: $NS$ is constant with value all functors, and $SN$ is likewise constant.

This is an elementary diagnostic for the manuscript's cancellation. We checked the subsequent transporter construction in \S6 and the preservation definition of smooth functors in \citet[Definition 4.6]{Maltsiniotis2005Asphericite}; we did not find this particular counterexample stated there. The mark above records our verification, not a claim of priority. The later block ``The stabilizer of categorical equivalences'' identifies the whole class $S(\mathrm{Equ})$.
\end{dcomment}


\section*{5\quad Dualité formelle}

Soient $\mathcal{P}$, $\mathcal{Q}$ deux ensembles ordonnés. J'essaye de dégager une notion de `dualité' (formelle) reliant $\mathcal{P}$ et $\mathcal{Q}$. Ce sera une relation entre $\mathcal{P}$ et $\mathcal{Q}$
\[
\mathcal{R} \subseteq \mathcal{P} \times \mathcal{Q}
\tag*{(5.1)}
\]
ayant les propriétés suivantes
\begin{enumerate}[leftmargin=*, labelindent=0.5em]
\item[(D)] Pour toute partie $A$ de $\mathcal{P}$, l'ensemble
\end{enumerate}
\[
A^\cup \overset{\text{déf}}{=} \{ q \in \mathcal{Q} \mid (p,q) \in \mathcal{R} \text{ pour tout } p \in A,\ \text{ i.e. } A \times \{q\} \subseteq \mathbf{R} \} % [sic] \mathbf{R} -> \mathcal{R}, the relation is written \mathcal{R} everywhere else
\]
est une partie ouverte de $\mathcal{Q}$ ayant un plus grand élément $\beta(A) : A^\cup = \mathcal{Q}/\beta(A)$. Dualement, pour toute partie $B$ de $\mathcal{Q}$, l'ensemble
\[
B^\cap \overset{\text{déf}}{=} \{ p \in \mathcal{P} \mid (p,q) \in \mathcal{R} \text{ pour tout } q \in B,\ \text{ i.e. } \{p\} \times B \subseteq \mathcal{R} \}
\]
est une partie ouverte de $\mathcal{P}$ ayant un plus grand élément $\alpha(B)$, i.e. $B^\cap = \mathcal{P}/\alpha(B)$.

Le plus souvent, $\mathcal{P}$ et $\mathcal{Q}$ seront stables par sup quelconques (donc aussi par inf quelconques). Dans ce cas, (D) équivaut, du côté des $A^\cup$ (pour $A \in \mathfrak{P}(\mathcal{P})$) à ceci :
\begin{enumerate}[leftmargin=*, labelindent=0em]
\item[(D$_0$)] Pour tout $p \in \mathcal{P}$, $\mathcal{R}(p) = \{ q \in Q \mid (p,q) \in \mathcal{R} \}$ est une partie ouverte \textit{filtrante} de $\mathcal{Q}$, ou ce qui revient au même, une partie non vide ($\Longleftrightarrow$ contenant le plus petit élément $\emptyset_{\mathcal{Q}}$ de $\mathcal{Q}$) stable par sup de deux éléments. De plus, $\mathcal{R}(p)$ est stable par sup filtrants, ou ce qui revient au même (étant lui-même filtrant), il contient un plus grand élément. Il revient % [sic] Q -> \mathcal{Q}, plain italic Q printed in "q \in Q" instead of the script \mathcal{Q}
\msp{page 70}
au même de dire que $\mathcal{R}(p)$ est un \textit{ouvert} de $\mathcal{Q}$ stable par sup quelconques (ce qui implique qu'il contient $\emptyset_{\mathcal{Q}}$). Dualement, pour tout $q \in \mathcal{Q}$, ${}^{s}\mathcal{R}(q) = \{ p \in \mathcal{P} \mid (p,q) \in \mathcal{R} \}$ est une partie \textit{ouverte} de $\mathcal{P}$, stable par sup quelconques, i.e. une partie ouverte ayant un plus grand élément. % doubt resolved against PDF leaf 15: raised s, transpose relation
\end{enumerate}

\dgrblock{dgrb128}{The relation comes before the order}
\begin{dcomment}{exposition}{p69}{The relation comes before the order}
Grothendieck pauses to isolate the formal mechanism that a successful transporter construction must satisfy. The first part of the section concerns an arbitrary relation, even though the initial presentation starts with ordered sets. Separating those two levels makes the notation more transparent and prevents the later confusion between a closed subset and an element of an ordered set.

For a relation $R\subseteq P\times Q$, put
\[
 A^\cup=\{q\mid (p,q)\in R\text{ for every }p\in A\},
 \qquad
 B^\cap=\{p\mid (p,q)\in R\text{ for every }q\in B\}.
\]
The decisive equivalence is
\[
 B\subseteq A^\cup\quad\Longleftrightarrow\quad
 A\times B\subseteq R\quad\Longleftrightarrow\quad
 A\subseteq B^\cap.
\]
Both operations reverse inclusions. Every $p\in A$ is related to every member of $A^\cup$, so $A\subseteq A^{\cup\cap}$. Applying the inclusion-reversing operation gives $A^{\cup\cap\cup}\subseteq A^\cup$; applying extensivity on the other side gives the reverse inclusion. Thus the triple operation equals the single operation, and the double operation is idempotent. This proves (5.5)--(5.7) without any order on $P$ or $Q$.

The fixed subsets form complete lattices. Arbitrary intersections of fixed subsets are fixed: if $A=\bigcap_i A_i$, monotonicity gives $A^{\cup\cap}\subseteq A_i$ for every $i$, and extensivity gives the opposite inclusion $A\subseteq A^{\cup\cap}$. Their join is the closure of their union. The two polar operations exchange these closed-set lattices by inverse order-reversing bijections, proving Proposition 5.1 and Corollary 5.2. Empty families cause no exception: the largest closed subset of $P$ is $P$, while the smallest is $Q^\cap$, which need not be empty.

The ordered version asks for more. An ``open'' subset in this passage is a lower set, and the additional requirement is that each polar lower set be principal. For complete lattices $P,Q$, this gives maps $\beta_0:P\to Q$ and $\alpha_0:Q\to P$ characterized by
\[
 q\leq\beta_0(p)\quad\Longleftrightarrow\quad p\leq\alpha_0(q).
\]
Equivalently, $\alpha_0:Q\to P^{\mathrm{op}}$ is left adjoint to $\beta_0:P^{\mathrm{op}}\to Q$. It follows that $\beta_0$ sends joins in $P$ to meets in $Q$, including the empty join. Conversely, if it has this property, set
\[
 \alpha_0(q)=\bigvee\{p\in P:q\leq\beta_0(p)\}.
\]
The set under the join is a lower set and is closed under arbitrary joins, so its join belongs to it; this proves the displayed equivalence and constructs the adjoint. The ensuing fixed-point correspondence is the usual antitone Galois connection. A systematic later account is \citet[\S0.4, Propositions 7--8 and Theorem 2]{GanterWille1999}; the arguments needed here are supplied above, independently of terminology.
\end{dcomment}


En d'autres termes, quand on postule la condition (D) pour $A$ resp. $B$ réduit à un élément, alors elle reste vérifiée pour toute partie. Cela provient des formules
\[
\left\{
\begin{array}{lcl}
(\bigcup_i A_i)^\cup & = & \bigcap_i A_i^\cup \\[1ex]
(\bigcup_j B_j)^\cap & = & \bigcap_j B_j^\cap,
\end{array}
\right.
\tag*{(5.2)}
\]
et du fait que l'intersection de parties ouvertes de la forme $\mathcal{P}/\alpha_i$ dans $\mathcal{P}$ est de la forme $\mathcal{P}/\alpha$, où $\alpha = \inf_i \alpha_i$, et de même dans $\mathcal{Q}$. Ainsi on trouve
\[
\left\{
\begin{array}{lcl}
\beta(\bigcup_i A_i) & = & \inf_i \beta(A_i) \\[1ex]
\alpha(\bigcup_j B_j) & = & \inf_j \alpha(B_j).
\end{array}
\right.
\tag*{(5.3)}
\]

% ---- from Der-ChII-234_p046-050.tex
\printedpage{44}
% head: A. Grothendieck : Dérivateurs, Ch. II, Édition des Universités de Montpellier II et Paris VII

On posera, pour une partie $A$ de $\cP$ resp. $B$ de $\cQ$
\[
A^{\natural} = (A^{\cup})^{\cap}, \qquad B^{\flat} = (B^{\cap})^{\cup}
\tag*{(5.4)}
\]
de sorte qu'on aura
\[
A \subseteq A^{\natural}, \qquad B \subseteq B^{\flat}
\tag*{(5.5)}
\]
On dira que $A$ est $\cR$-\textit{clos} si $A = A^{\natural}$ (ou encore si $A$ est de la forme $B'^{\cap}$ pour $B' \subseteq \cQ$, cf. 5.6), que
\msp{71}
$B$ est $\cR$-\textit{clos} si $B = B^{\flat}$. On aura pour tout $A \subseteq \cP$, $B \subseteq \cQ$
\[
\left\{
\begin{array}{rcll}
(A^{\cup})^{\flat} & = & A^{\cup}, & \text{i.e. } A^{\cup} \text{ est } \cR\text{-clos,}\\[1ex]
(B^{\cap})^{\natural} & = & B^{\cap}, & \text{i.e. } B^{\cap} \text{ est } \cR\text{-clos.}
\end{array}
\right.
\tag*{(5.6)}
\]
En particulier
\[
(A^{\natural})^{\natural} = A^{\natural}, \qquad (B^{\flat})^{\flat} = B^{\flat},
\tag*{(5.7)}
\]
i.e. la $\cR$-clôture $A^{\natural}$ de $A \subseteq \cP$ est $\cR$-clos, et de même la $\cR$-cloture $B^{\flat}$ de $B \subseteq \cQ$. % [sic] cloture -> clôture, accent circonflexe manquant (cf. « clôture » dans la même ligne)

En fait, je m'aperçois que je suis en train de mélanger deux espèces de faits différents. Les relations (5.2), et (5.4) à (5.7), ne dépendent pas de la donné de relations d'ordre sur $\cP$, $\cQ$, ni d'axiomes reliant celles-ci à la relation donnée $\cR$, mais seulement de cette relation, laquelle n'a pas à être soumise à aucune condition. Notons que % [sic] la donné -> la donnée, accord manquant
\[
\emptyset^{\cup} = \cQ, \qquad \emptyset^{\cap} = \cP,
\tag*{(5.8)}
\]
tandis que
\[
\emptyset^{\natural} = \cQ^{\cap}, \qquad \emptyset^{\flat} = \cP^{\cup}
\tag*{(5.9)}
\]
sont respectivement la plus petite partie $\cR$-close de $\cP$ et de $\cQ$.

L'application
\[
A \longmapsto A^{\natural}
\tag*{(5.10)}
\]
est une application \textit{croissante} (et \textit{idempotente}) de $\fP(\cP)$ dans lui-même, ayant donc la propriété
\msp{72}
\[
(\bigcup_i A_i)^{\natural} = (\bigcup_i A_i^{\natural})^{\natural}.
\]

\printedpage{45}
% head: A. Grothendieck : Dérivateurs, Ch. II, Édition des Universités de Montpellier II et Paris VII

Pour tout $A \subseteq \cP$, $A^{\natural}$ est la plus petite partie $\cR$-close de $\cP$ contenant $A$, et de même pour tout $B \subseteq \cQ$, $B^{\flat}$ est la plus petite partie $\cR$-close de $Q$ contentant $B$. On a de plus (c'était même à dire avant de dire que $A \longmapsto A^{\natural}$ est croissante) % [sic] de $Q$ -> de $\mathcal{Q}$, lettre droite au lieu de la lettre calligraphique ; [sic] contentant -> contenant, faute de frappe
\[
\left\{
\begin{array}{rcl}
A \subseteq A' & \Longrightarrow & A'^{\cup} \subseteq A^{\cup}\\[1ex]
B \subseteq B' & \Longrightarrow & B'^{\cap} \subseteq B^{\cap},
\end{array}
\right.
\tag*{(5.11)}
\]
i.e. les applications
\[
\begin{tikzcd}
\fP(\cP) \arrow[r, shift left, "\beta"] & \fP(\cQ) \arrow[l, shift left, "\alpha"]
\end{tikzcd}
\qquad
\begin{array}{l}
A \longmapsto A^{\cup}\\
B \longmapsto B^{\cap}
\end{array}
\tag*{(5.12)}
\]
sont décroissantes (ce qui implique que les deux composés $\alpha\beta = (A \longmapsto A^{\natural})$ et $\beta\alpha = (B \longmapsto B^{\flat})$ sont croissantes). Leurs restrictions aux sous-ensembles $\fP_{\cR}(\cP)$ et $\fP_{\sR}(\cQ)$ des parties $\cR$-closes de $\cP$ resp. de $\cQ$, sont, par définition de celles-ci comme
\[
\begin{array}{rcl}
\fP_{\cR}(\cP) & = & \{A \subseteq \cP \mid A^{\natural} \ (\overset{\text{déf}}{=} \alpha\beta(A)) = A\}\\
\fP_{\sR}(\cQ) & = & \{B \subseteq \cQ \mid B^{\flat} \ (\overset{\text{déf}}{=} \beta\alpha(B)) = B\},
\end{array}
\]
inverse l'une de l'autre. On trouve ainsi

\textbf{Proposition 5.1.} \textit{Soient $\cP$, $\cQ$ deux ensembles, $\cR \subseteq \cP \times \cQ$ une relation. Alors les applications $\alpha$, $\beta$ de (5.12) induisent des bijections inverses l'}
\msp{73}
\textit{une de l'autre, entre l'ensemble $\mathit{Cl}_{\cR}(\cP)$ des parties $\cR$-closes de $\cP$, et l'ensemble $\mathit{Cl}_{\sR}(\cQ)$ des parties ${}^{s}\cR$-closes (nous dirons $\cR$-closes par abus de langage) de $\cQ$. Ces bijections sont décroissantes, donc des antiisomorphismes pour les relations d'ordre naturelles (d'inclusion).}

\textbf{Corollaire 5.2.} \textit{Soit $(A_i)_i$ une famille de parties $\cR$-closes de $\cP$. Alors $\bigcap A_i$ est une partie $\cR$-close de $\cP$, et}
\[
(\bigcap_i A_i)^{\cup} = \underbrace{(\bigcup_i A_i^{\cup})^{\flat}}_{\substack{\text{sup des}\\ \text{parties } \cR\text{-closes}\\ A_i^{\cup} \text{ de } \cQ}} \quad .
\tag*{(5.2.1)}
\]

Mais la relation $\cR$ se `prolonge' en une relation $\bar{\cR}$ entre
\[
\bar{\cP} = \fP(\cP), \qquad \bar{\cQ} = \fP(\cQ)
\tag*{(5.13)}
\]
par la prescription
\[
\underbrace{\bar{\cR}}_{\subseteq \bar{\cP} \times \bar{\cQ}} = \{(A,B) \in \bar{\cP} \times \bar{\cQ} \mid B \subseteq A^{\cup},\ \text{i.e. } A \subseteq B^{\cap},\ \text{i.e. } A \times B \subseteq \cR\}.
\tag*{(5.14)}
\]
Les ensembles $\bar{\cP}$, $\bar{\cQ}$ sont stables par sup quelconques (réunion de parties de $\cP$ resp. $\cQ$), et la relation $\bar{\cR}$ sur ces ensembles ordonnés satisfait à la condition (D$_0$)
\printedpage{46}
% head: A. Grothendieck : Dérivateurs, Ch. II, Édition des Universités de Montpellier II et Paris VII
\msp{74}
ci-dessus (p. 69), de façon évidente : pour $A$ fixé, les $B$ tels que $(A,B) \in \bar{\cR}$ sont ceux qui sont contenus dans $A^{\cup}$, donc l'ensemble $\bar{\cR}(A)$ n'est autre que $\bar{\cQ}/A^{\cup}$, et dualement ${}^{s}\bar{\cR}(B)$ n'est autre que $\bar{\cP}/B^{\cap}$. Ainsi
\[
\beta(A) = A^{\cup}, \qquad \alpha(B) = B^{\cap}.
\tag*{(5.14)} % [sic] (5.14) -> numéro distinct, numéro déjà employé pour la formule (5.14) de la page précédente
\]
Si j'ai songé à introduire l'axiome (D$_0$), pour un couple d'ensembles ordonnés quelconques, c'est à cause des cas où ceux-ci ne sont \textit{pas} de la forme $\fP(\cP)$, $\fP(\cQ)$, et (alors même qu'ils le seraient) où la relation de base n'est pas déduite comme ci-dessus d'une relation quelconque entre ensembles $\cP$, $\cQ$.

Je reviens donc à l'axiomatique du début, avec $\cP$, $\cQ$ ordonnés, et $\cR$ satisfaisant (D). Dans tous les cas auxquels je peux songer, $\cP$ et $\cQ$ sont stables par sup quelconques, donc (D) équivaut à (D$_0$). Notons que l'application
\[
\cP \longrightarrow \fP(\cP), \qquad p \longmapsto \cP/p
\]
est injective, et de même pour $\cQ \longrightarrow \fP(\cQ)$. Ainsi, la connaisance de $A \longmapsto A^{\cup}$ équivaut à celle de % [sic] connaisance -> connaissance, faute d'orthographe
\[
\begin{array}{rcl}
A & \longmapsto & \beta(A)\\[0.5ex]
\beta : \fP(\cP) & \longrightarrow & \cQ,
\end{array}
\]
et de même celle de $B \longmapsto B^{\cap}$ équivaut à
\msp{75}
celle de
\[
\alpha : \fP(\cQ) \longrightarrow \cP.
\]
En fait, il suffit de se donner
\[
\begin{array}{rcccl}
\beta_0 : & \cP & \longrightarrow & \cQ & \qquad \text{applications}\\
\alpha_0 : & \cQ & \longrightarrow & \cP & \qquad \textit{décroissantes}
\end{array}
\tag*{(5.15)}
\]
définies par
\[
\left\{
\begin{array}{rcl}
\beta_0(p) & = & \beta(\{p\})\\[1ex]
\alpha_0(q) & = & \alpha(\{q\}),
\end{array}
\right.
\]
puisqu'on en déduit $\beta$, $\alpha$ par
\[
\left\{
\begin{array}{rcl}
\beta(A) & = & \inf\nolimits_{p \in A} \beta_0(p)\\[1ex]
\alpha(B) & = & \inf\nolimits_{q \in B} \alpha_0(q).
\end{array}
\right.
\tag*{(5.16)}
\]
La relation $\cR$ se déduit de la donnée de $\beta_0$, ou au choix de $\alpha_0$, par
\[
(p,q) \in \cR \iff q \leq \beta_0(p) \iff p \leq \alpha_0(q).
\tag*{(5.17)}
\]

\printedpage{47}
% head: A. Grothendieck : Dérivateurs, Ch. II, Édition des Universités de Montpellier II et Paris VII

Cela implique donc que $\alpha_0$, $\beta_0$ se déterminent mutuellement, par les formules
\[
\left\{
\begin{array}{rcl}
\alpha_0(q) & = & \sup \underbrace{\{p \in \cP \mid q \leq \beta_0(p)\}}_{\{q\}^{\cap}}\\[3ex]
\beta_0(p) & = & \sup \underbrace{\{q \in \cQ \mid p \leq \alpha_0(q)\}}_{\{p\}^{\cup}}.
\end{array}
\right.
\tag*{(5.18)}
\]
Reste à se demander quelles sont les applications $\alpha_0 : \cP \longrightarrow \cQ$, ou $\beta_0 : \cQ \longrightarrow \cP$, qui correspondent à une \textit{`relation dualisante'} % [sic] \alpha_0 : \cP \to \cQ, \beta_0 : \cQ \to \cP -> \beta_0 : \cP \to \cQ, \alpha_0 : \cQ \to \cP, les rôles de \alpha_0 et \beta_0 sont échangés par rapport à (5.15)
\msp{76}
$\cR$ sur $\cP \times \cQ$. Voici les conditions sur $\beta_0 : \cP \longrightarrow \cQ$ :
\begin{itemize}[leftmargin=4em, labelsep=0.8em]
\item[\circled{a}] $\beta_0$ est \textit{décroissante},
\end{itemize}
ce qui implique que la relation
\[
\underbrace{\cR_{\beta_0}}_{\subseteq \cP \times \cQ} = \{(p,q) \in \cP \times \cQ \mid q \leq \beta_0(p)\}
\]
satisfait à la condition (D$_0$) du côté gauche, et il faut encore la vérifier du côté droit, i.e. regarder, pour $q$ fixé, l'ensemble
\[
\{q\}^{\cap} = \{p \in \cP \mid q \leq \beta_0(p)\}.
\]
Comme $\beta_0$ est décroissant, on voit que $p \in \{q\}^{\cap}$, et $p' \leq p \Longrightarrow p' \in \{q\}^{\cap}$, i.e. $\{q\}^{\cap}$ est \textit{ouvert.} On doit vérifier encore (supposant l'existence des sup quelconques dans $\cP$) que cet ensemble est stable par sup quelconques - il en résultera qu'il est de la forme $\cP/\alpha(q)$. Cela équivaut à la condition
\begin{itemize}[leftmargin=4em, labelsep=0.8em]
\item[\circled{b}] $\beta_0$ transforme sup quelconques dans $\cP$ en inf dans $\cQ$.
\end{itemize}
Si on regarde $\beta_0$ comme un foncteur contravariant $\cP^{o} \longrightarrow \cQ$, cela signifie qu'il commute aux limites projectives, ou encore qu'il a un adjoint à gauche $\cQ \longrightarrow \cP^{o}$, lequel commute aux limites inductives, et correspond
\msp{77}
à l'application décroissante $\alpha_0 : \cQ \longrightarrow \cP$, transformant sup en inf.

Ici il est commode de considérer $\natural$ et $\flat$ comme des opérations sur $\cP$ resp. $\cQ$, par
\[
p^{\natural} = \alpha_0\beta_0(p), \qquad q^{\flat} = \beta_0\alpha_0(q).
\tag*{(5.19)}
\]
La relation entre ces endomorphismes, et l'opération de même nom sur $\fP(\cP)$, $\fP(\cQ)$, est donnée par
\[
\left\{
\begin{array}{rcl}
p^{\natural} & = & \text{Plus grand élément de } \{p\}^{\natural}\\[1ex]
A^{\natural} & = & \sup\nolimits_{p \in A} p^{\natural},
\end{array}
\right.
\tag*{(5.20)}
\]
\printedpage{48}
% head: A. Grothendieck : Dérivateurs, Ch. II, Édition des Universités de Montpellier II et Paris VII
et symétriquement du côté de $\cQ$.

On dit que $p$ est $\cR$-\textit{clos,} ou $\beta_0$-clos, ou $\alpha_0$-clos, si $p = p^{\natural}$, et de même $q \in \cQ$ est dit $\cR$-\textit{clos} etc. si $q = q^{\flat}$. Les applications
\[
\left\{
\begin{array}{ccccccc}
p & \longmapsto & p^{\natural} & : & \cP & \longrightarrow & \cP\\[1ex]
q & \longmapsto & q^{\flat} & : & \cQ & \longrightarrow & \cQ
\end{array}
\right.
\tag*{(5.21)}
\]
sont \textit{croissantes} et \textit{idempotentes.} On peut caractériser $p^{\natural}$ comme le plus petit élément clos de $\cP$ majorant $p$, et de même pour $q^{\flat}$. On trouve alors, par (5.19), que $\alpha_0$ et $\beta_0$ induisent des bijections inverses l'une de l'autre,
\msp{78}
entre l'ensemble des éléments clos de $\cP$, et l'ensemble des éléments clos de $\cQ$
\[
\begin{tikzcd}
\cP^{\natural} \arrow[r, shift left, "\bar{\beta}_0"] \arrow[r, phantom, "\scriptstyle\sim"] & \cQ^{\flat}. \arrow[l, shift left, "\bar{\alpha}_0"]
\end{tikzcd}
\tag*{(5.22)}
\]
Mais en pratique, les applications $\alpha_0$, $\beta_0$ ne sont pas obtenues directement, mais on dispose de la relation $\cR$, satisfaisant la condition (D$_0$), dont on déduit $\alpha_0$, $\beta_0$. Dans tous les cas auxquels je pense ici, ou bien $\cP$, $\cQ$ sont de la forme $\fP(P_0)$, $\fP(Q_0)$, et $\cR$ se déduit d'une relation $R_0 \subseteq P_0 \times Q_0$ comme ci-dessus (de loin le plus fréquent, il me semble), ou bien $\cP$ et $\cQ$ apparaissent comme sous-ensembles convenables d'ensembles $\fP(P_0)$, $\fP(Q_0)$, munis de l'ordre induit, et la relation $\cR$ sur $\cP \times \cQ$ s'obtient d'une façon plus complexe.

\dgrblock{dgrb129}{A join of closed elements need not be closed}
\begin{dcomment}{examples}{p77}{A join of closed elements need not be closed}
The ordered version allows the closure of a subset to be represented by one element, and Grothendieck wants to pass freely between these descriptions. Formula (5.20) loses a closure in making this passage. The difficulty is mathematical, not a transcription error: the edition prints the second equality exactly as transcribed.

Let $c=\alpha_0\beta_0$ and let $A\subseteq P$. The precise identities are
\[
 \{p\}^{\cup\cap}=\mathord\downarrow c(p),\qquad
 A^{\cup\cap}=\mathord\downarrow c\left(\bigvee_{p\in A}p\right)
 =\mathord\downarrow c\left(\bigvee_{p\in A}c(p)\right).
\]
Here $\mathord\downarrow x$ denotes the set of elements below $x$. To prove the middle equality, the polar $A^\cup$ has greatest element
$\bigwedge_{p\in A}\beta_0(p)=\beta_0(\bigvee A)$, and its polar has greatest element $\alpha_0\beta_0(\bigvee A)$. For the last equality, $p\leq c(p)\leq c(\bigvee A)$ for every $p\in A$; apply monotonicity and idempotence to the two resulting inequalities. Thus the supremum can be taken in the \emph{closed-element lattice}, where it means closure of the ambient supremum. It cannot in general be taken unchanged in $P$.

A small example displays both the omitted operation and its importance. Let $E=\mathbf F_2^2$, and take $P=\mathfrak P(E)$, $Q=\mathfrak P(E^*)$, with the relation that every vector in one subset is annihilated by every functional in the other. Double annihilation sends a subset to its linear span. To check this directly, a linear functional vanishing on a set vanishes on its span; if a vector lies outside that span, extend a basis of the span by that vector and prescribe a functional taking value $1$ on it and $0$ on the original basis.

Let $e_1,e_2$ be a basis and put $p_1=\{0,e_1\}$, $p_2=\{0,e_2\}$. Each $p_i$ is already closed, but their supremum in $P$ is
\[
 p_1\cup p_2=\{0,e_1,e_2\},
\]
which omits $e_1+e_2$ and is not closed. For the family $A=\{p_1,p_2\}\subseteq P$, the greatest element of $A^{\cup\cap}$ is all of $E$, whereas the supremum printed in (5.20) is only the three-element union. Moreover, $A^{\cup\cap}$ itself is a subset of $P$, namely all subsets of $E$, rather than an element of $P$. These two distinctions are exactly what the corrected formula records.

This correction does not disturb the subsequent statements that $c$ is monotone, extensive and idempotent, or that the closed elements on the two sides correspond anti-isomorphically. It supplies the correct join operation on those elements and will be needed again when the elements of $P$ are themselves classes of morphisms.
\end{dcomment}


\noindent\textbf{Exemples.}
\begin{enumerate}[label={\arabic*)}]
\item Relations d'\textit{orthogonalité} entre parties de deux modules sur un anneau, muni d'un accouplement de dualité. Dans le cas des \texttt{[espaces]} vectoriels de dimension finie et d'une vraie dualité, les parties closes sont les sous-espaces vectoriels, dans le cas infini, ce sont les sous-espaces vectoriels fermés pour la `topologie faible' définie par l'accouplement.
\item Relations d'\textit{orthogonalité} entre parties de deux groupes topologiques commutatifs localement compacts duales l'une de l'autre, au sens de Pontrjagin.
\msp{79}

Les parties closes sont les sous-groupes fermés.
\item Pour deux espaces vectoriels `en dualité' sur un corps valué complet, tel $\mathbf{R}$ où $\mathbf{C}$, la relation % [sic] où -> ou, accent grave fautif
\[
|\langle x, x' \rangle| \leq 1.
\]
Dans le cas réel ou complexe, les parties closes sont les parties convexes cerclées [\,?] faiblement fermées.
\item Pour les espaces vectoriels réels, la relation de semi-polarité
\[
\langle x, x' \rangle \geq -1.
\]
Les parties closes sont les parties convexes faiblement fermées contenant l'origine.
\end{enumerate}

\dgrblock{dgrb130}{What the double polar closes}
\begin{dcomment}{examples}{p79:Les parties closes sont les parties convexes faiblement fermées contenant l'origine.}{What the double polar closes}
The examples identify the content of a formal closure: the relation supplies tests, and the double polar retains everything that those tests cannot distinguish. The algebraic and analytic identifications require their own separation results; the abstract Galois calculation alone does not prove them.

For a separating dual pair of vector spaces $E,F$ over a field $k$, the annihilator of $A\subseteq E$ is $A^\perp=\{\varphi\in F:\langle a,\varphi\rangle=0\text{ for all }a\in A\}$. Its double annihilator is the closure of $\operatorname{span}(A)$ for the weak topology of the pairing, taking $k$ discrete in the purely algebraic formulation. Here is a direct finite-coordinate proof. If $x$ is not in the weak closure of a subspace $V$, there are $\varphi_1,\ldots,\varphi_n\in F$ for which the vector of values of $x$ is outside the image of $V$ in $k^n$. A linear functional on $k^n$ vanishes on that image but not on this vector. The corresponding linear combination of the $\varphi_i$ separates $x$ from $V$. Conversely every annihilating functional is weakly continuous, so its kernel contains the weak closure of $V$. For real or complex weak topologies the same argument applies because finite-dimensional linear subspaces are closed. In finite-dimensional perfect duality all subspaces are closed, giving the familiar double-annihilator identity.

For a Pontryagin dual pair $G,\widehat G$, the annihilator of $A\subseteq G$ consists of the characters trivial on $A$. It is unchanged upon replacing $A$ by the closed subgroup it generates. If $x$ lies outside that closed subgroup $H$, the locally compact abelian quotient $G/H$ has a continuous character nontrivial on $xH$. Composing it with the quotient map gives a character annihilating $A$ but not $x$. Thus the double annihilator is $\overline{\langle A\rangle}$. The separation step is the substantive theorem; \citet[Chapter 7, Proposition 38]{Morris1977} gives this annihilator statement and its proof from character separation.

For the absolute polar over $\mathbf R$ or $\mathbf C$, the answer is instead the weakly closed absolutely convex hull of $A\cup\{0\}$. One inclusion follows because each inequality $|\langle -,\varphi\rangle|\leq1$ defines a weakly closed absolutely convex set containing $0$. For the other, separate a point outside that hull by a weakly continuous linear functional; the absolute convexity allows multiplication by a scalar of modulus one, and rescaling gives a functional with absolute value at most $1$ on the hull and greater than $1$ at the point. The weak continuous dual is precisely the specified testing space. The Hahn--Banach separation theorem therefore supplies the desired test. In the real semi-polar convention $\langle x,\varphi\rangle\geq-1$, negate the separating functional; the double polar is the weakly closed convex hull of $A\cup\{0\}$, without the balancedness condition. These are the two forms of the bipolar theorem; compare \citet[Chapter IV, \S1, Theorem 1.5]{SchaeferWolff1999}.

The origin must not disappear from the statement. Under conventions allowing the empty set to count as convex and balanced, the list of absolute-polar closed sets must explicitly exclude it: $\varnothing^{\circ\circ}=\{0\}$ in a separating pair. In $\mathbf R^2$, for example, the absolute polar of $\{e_1,e_2\}$ is the square $\{|a|\leq1,|b|\leq1\}$, and its polar is the diamond $\{|x_1|+|x_2|\leq1\}$. The inequality follows by maximizing $ax_1+bx_2$ over the square; choosing the signs of $a,b$ gives equality. This makes the distinction between linear span, convex hull and absolutely convex hull visible in the same two-dimensional setting.
\end{dcomment}


% ---- from Der-ChII-234_p051-055.tex
\printedpage{49}
% head: A. Grothendieck : Dérivateurs, Ch. II, Édition des Universités de Montpellier II et Paris VII
% foot: Version 8 juillet 2001

\begin{itemize}[leftmargin=2.5em]
\item[5)] Soit $K$ un corps (pour simplifier), $L$ une extension galoisienne de groupe $G$. Relation entre $L$ et $G$
\[
\gamma\cdot x=x,\qquad \text{i.e. } x\in L^{\gamma}.
\]
Les parties closes de $L$ sont les sous-extensions, les parties closes de $G$ les sous-groupes de $G$.
\end{itemize}

\dgrblock{dgrb131}{The topology in the Galois example}
\begin{dcomment}{examples}{p79:les sous-groupes de $G$.}{The topology in the Galois example}
The fixed-point relation offers another instructive instance of the same issue: the relation always gives a closure, but the description of its fixed subsets changes when the extension becomes infinite. For a Galois extension $L/K$, the polar of a subset $S\subseteq L$ is the group of automorphisms fixing every element of $S$. Its double polar is the intermediate field $K(S)$. For a subset $H\subseteq G=\operatorname{Gal}(L/K)$, the double polar is the closure of the subgroup generated by $H$ in the Krull topology. Thus all subgroups are closed when the extension is finite; for an arbitrary Galois extension the manuscript's assertion must read \emph{closed subgroups}. See \citet[Proposition 7.12 and Theorem 7.13]{Milne2022FT}, where both directions of this correspondence are proved.

There is a concrete obstruction to omitting that qualification. Take $L=\overline{\mathbf F}_p$ and $K=\mathbf F_p$. The Galois group is $\widehat{\mathbf Z}$, with the Frobenius automorphism corresponding to $1$. The cyclic subgroup generated by Frobenius is isomorphic to $\mathbf Z$; its image surjects onto every finite quotient $\mathbf Z/n$, so it is dense, but it is not all of the uncountable group $\widehat{\mathbf Z}$. An element fixed by Frobenius satisfies $x^p=x$, hence lies in $\mathbf F_p$. The double polar of this cyclic subgroup is consequently the whole Galois group. It is not a closed subset for the formal polarity, even though it is a subgroup.

This correction is already part of classical infinite Galois theory, rather than a later repair peculiar to derivators. The nearby allusion to inseparable analogues requires a different testing object, such as a suitable restricted Lie algebra in a height-one purely inseparable extension. No precise inseparable theorem is stated here, so the audit records an allusion, not an unrestricted extension of the Galois correspondence.
\end{dcomment}


Il y a des analogues divers dans des cas inséparables (à coups de $p$-algèbres de Lie par exemple), et des énoncés de descente en tous genres.

\msp{page 80}Tous ces exemples sont de 1\textsuperscript{ère} espèce, les ensembles ordonnés qui interviennent sont de la forme $\Pf(P_0)$, $\Pf(Q_0)$, et la relation entre eux est déduite d'une relation $R_0$ entre $P_0$ et $Q_0$. Voici maintenant des exemples plus complexes.

\begin{itemize}[leftmargin=2.5em]
\item[6)] Soit $\mathcal{M}$ une catégorie (‘grosse’, le plus souvent \dots). On prend
\[
\mathcal{P}=\mathcal{Q}=\Pf(\Fl(\mathcal{M})).
\]
La relation $\mathcal{R}\subseteq\mathcal{P}\times\mathcal{Q}$ est la relation de Quillen, que j'écris, pour $\Phi,\Psi\in\Pf$
\[
\Psi\subseteq\Phi_{*},
\]
ou encore de façon équivalente
\[
\Phi\subseteq\Psi^{*},
\]
($\Psi$ a la LLP \texttt{[left lifting property]} pour $\Psi$, $\Phi$ la RLP \texttt{[right lifting property]} pour $\Psi$). Donc les applications $\alpha_0$, $\beta_0$ sont données par % [sic] Psi a la LLP pour Psi -> Psi a la LLP pour Phi, the relation is between Psi and Phi (cf. "Phi la RLP pour Psi")
\[
\beta_0(\Phi)=\Phi_{*},\qquad \alpha_0(\Psi)=\Psi^{*}.
\]
L'idée ici ne viendrait à personne d'introduire $\Pf(\mathcal{P})$, $\Pf(\mathcal{Q})$ - c'est déjà assez gros comme ça ! La notation pour $\Phi^{\natural}$ est ici $\tilde{\Phi}$, pour $\Psi^{\flat}$ c'est $\bar{\Psi}$. Les parties closes sont celles que j'appelle Q-closes à gauche et à droite.

La relation $\mathcal{R}$, que j'écris
\[
\Phi\longleftrightarrow\Psi,
\]
\msp{page 81}n'est pas déduite d'une relation entre flèches de $\mathcal{M}$, i.e. d'une partie de $\Fl(\mathcal{M})\times\Fl(\mathcal{M})$, mais est de nature plus complexe. Pour la décrire, il faut introduire l'ensemble $\mathcal{D}$ des diagrammes carrés commutatifs (traits pleins) de la forme
\[
\begin{tikzcd}
A \arrow[r] \arrow[d, "i"] & X \arrow[d, "p"] \\
B \arrow[r] \arrow[ur, dotted, "f"] & Y
\end{tikzcd}
\tag*{$(*)$}
\]
\printedpage{50}
% head: A. Grothendieck : Dérivateurs, Ch. II, Édition des Universités de Montpellier II et Paris VII
% foot: Version 8 juillet 2001
(qu'on peut interpréter comme $\Fl(\underline{\Fl}(\mathcal{M}))$, ensemble des flèches dans $\underline{\Fl}(\mathcal{M})$). On a une partie $\mathcal{D}_0$ de $\mathcal{D}$, formée des diagrammes pour lesquels \texttt{[il]} existe \texttt{[un]} $f$ avec deux triangles commutatifs (c'est donc l'image de $\underbrace{\Fl_3(\mathcal{M})}_{\mathrm{Hom}(\boldsymbol{\Delta}^3,\mathcal{M})}$ dans $\mathrm{Hom}(\boldsymbol{\Delta}^1\times\boldsymbol{\Delta}^1,\mathcal{M})$, transposée de $\boldsymbol{\Delta}^1\times\boldsymbol{\Delta}^1\longrightarrow\boldsymbol{\Delta}^3$ donné par $(0,0)\longmapsto 0$, $(0,1)\longmapsto 1$, $(1,0)\longmapsto 2$, $(1,1)\longmapsto 3$), et une application
\[
\mathcal{D}_0\subseteq\mathcal{D}\xrightarrow{\ \varphi\ }\Fl(\mathcal{M})\times\Fl(\mathcal{M}),
\]
associant à $(*)$ le couple $(i,j)$ \texttt{[plutôt} $(i,p)$\texttt{]} (source et but de $(*)$, considéré comme flèche dans $\underline{\Fl}(\mathcal{M})$ de $i$ dans $p$). La relation $\mathcal{R}$ est alors déduite de $\varphi$, $\mathcal{D}_0$ par la condition
\msp{page 82}
\[
\Phi\longleftrightarrow\Psi\qquad\Longleftrightarrow\qquad\varphi^{-1}(\Phi\times\Psi)\subseteq\mathcal{D}_0.
\]
Mais à vrai dire, je m'aperçois que c'est quand-même déduit d'une relation entre éléments de $\Fl(\mathcal{M})$, et cette relation entre $i : A\longrightarrow B$ et $p : X\longrightarrow Y$ peut s'écrire ainsi :
\begin{itemize}[leftmargin=3em]
\item[$(**)$] pour toutes flèches $A\longrightarrow X$, $B\longrightarrow Y$ rendant commutatif le diagramme $(*)$, la flèche $\varphi$ pointillée dans $(*)$ existe. % [sic] la flèche $\varphi$ -> la flèche $f$, the dotted arrow in $(*)$ is labelled $f$
\end{itemize}
\end{itemize}

Cela me fait penser que si on a une relation $\mathcal{R}$ entre $\Pf(P)$ et $\Pf(Q)$, satisfaisant $(\mathrm{D}_0)$, donc décrite par une application décroissante
\[
\Pf(P)\longrightarrow\Pf(Q)
\]
transformant sup en inf, la donné d'une telle application équivaut à celle d'une application quelconque % [sic] la donné -> la donnée, feminine noun
\[
P\longrightarrow\Pf(Q),
\]
ou ce qui revient au même, d'une relation $\mathcal{R}$ sur $P\times Q$. Donc cette structure doit être regardée comme ‘de première espèce’. La situation change à partir du moment où pour décrire la relation $\mathcal{R}$ avec les propriétés voulues, on se trouve contraint de remplacer $\Pf(P)$, $\Pf(Q)$ par des sous-ensembles convenables - quand on ne regarde que certaines parties de $P$, $Q$. C'est ce qu'on
\msp{page 83}va faire au § suivant. Le cas que je vais y traiter est le seul de son espèce que j'ai rencontré - un cas de ‘dualité formelle’ qui s'introduit de façon naturelle, et qui ne soit ‘de 1\textsuperscript{ère} espèce’.

\dgrblock{dgrb132}{Lifting gives a polarity before it gives factorization}
\begin{dcomment}{later}{p82}{Lifting gives a polarity before it gives factorization}
Grothendieck's second thought about the lifting example is the correct one. A lifting problem involves a square, but the assertion that \emph{every} square with a fixed left edge $i$ and right edge $p$ has a filler is a relation on the pair $(i,p)$. Write $i\mathbin{\square}p$ for that assertion. If $\Psi$ is a class of left maps and $\Phi$ a class of right maps, then
\[
 \Psi\subseteq{}^\square\Phi
 \quad\Longleftrightarrow\quad
 \Phi\subseteq\Psi^\square.
\]
The maps are order-reversing and their composites are closure operators by the preceding calculation. The text switches the order of the two factors in the square-based description: if $\varphi$ records the pair $(i,p)$, the corresponding condition is $\varphi^{-1}(\Psi\times\Phi)\subseteq\mathcal D_0$. This is a notational slip in the printed source, not an obstruction to the construction.

This closure formalism isolates the lifting half of Quillen's framework. A mutually closed pair $(\mathcal L,\mathcal R)$ need not supply a factorization of every arrow. That additional existence condition is what makes it a weak factorization system. In the notation used by \citet[Definition 11.2.1]{Riehl2014}, one requires both $\mathcal L={}^\square\mathcal R$, $\mathcal R=\mathcal L^\square$, and a factorization of every map as a left map followed by a right map. The small object argument supplies factorizations under separate generation and smallness hypotheses; its conclusion is not contained in the order-theoretic anti-isomorphism.

The distinction can be checked in a category with very few arrows. Let $\mathcal M$ be the ordered set with objects $0<1$ and its unique nonidentity arrow $a:0\to1$. A square with both vertical edges $a$ and identity horizontal edges has no filler, so $a$ is not left orthogonal to itself. Identities lift against every map. It follows by checking the only remaining pairs that $({}\!\operatorname{Fl}(\mathcal M),\operatorname{Is}(\mathcal M))$ and $(\operatorname{Is}(\mathcal M),\operatorname{Fl}(\mathcal M))$ are mutually closed; their factorizations happen to be the trivial ones. This example illustrates the elementary relation, without using any homotopy theory. In a general category one must still establish factorization independently.

Finally, every map from a powerset to a powerset that turns arbitrary unions into intersections is determined by its values on singletons:
\[
 \beta(A)=\bigcap_{p\in A}\beta(\{p\}).
\]
Defining $pRq$ by $q\in\beta(\{p\})$ recovers the relation. This is the exact reason the lifting example returns to Grothendieck's ``first species''. The transporter construction differs because its admissible right-hand classes must already be stable under pullback, while its left-hand classes will have to be stable under composition.
\end{dcomment}


\printedpage{51}
% head: A. Grothendieck : Dérivateurs, Ch. II, Édition des Universités de Montpellier II et Paris VII
% foot: Version 8 juillet 2001
\section*{6\quad Transporteurs, noyaux, stabiliseurs relatifs à $W$}

Soit $\mathcal{M}$ une catégorie, munie d'une partie
\[
W\subseteq\Fl(\mathcal{M}).
\tag*{(6.1)}
\]
Je pense au cas d'un localiseur, mais n'introduirai d'hypothèse sur $W$ qu'au besoin. Pour simplifier, je suppose $\mathcal{M}$ stable par produits fibrés.

Je pose
\[
\left\{
\begin{array}{rcl}
\mbox{(\footnotemark)}\mathcal{P} & = & \Pf(W)=\{A\in\Pf(\Fl(\mathcal{M}))\mid A\subseteq W\}\\[1ex]
\mathcal{Q} & = & \{L\in\Pf(\Fl(\mathcal{M}))\mid\
\parbox[c]{17em}{$L$ \textit{stable} par changement de\\
base (par qui je sous-entends aussi,\\
stable par isomorphie)\\
et par composition. $\}$.}
\end{array}
\right.
\tag*{(6.2)}
\]
\footnotetext{cf. rectification p. 86.}
J'introduis une relation
\[
\mathcal{R}\subseteq\mathcal{P}\times\mathcal{Q}
\]
de la façon suivante. Pour $A\in\mathcal{P}$, $L\in\mathcal{Q}$, on pose
\[
(A,L)\in\mathcal{R}\iff
\begin{array}[t]{l}
\text{Pour tout diagramme cartésien}\\
\begin{tikzcd}[ampersand replacement=\&]
X \arrow[d, "f"] \& X' \arrow[l] \arrow[d, "f'"] \\
Y \& Y' \arrow[l, "p"]
\end{tikzcd}\\
\text{avec $f\in A$, $p\in L$, on a $f'\in W$}\\
\text{(i.e. $A$ est stable par changement de base par $p\in L$).}
\end{array}
\tag*{(6.3)}
\]

\msp{page 84}Quand j'ordonne $\mathcal{P},\mathcal{Q}\subseteq\Pf(\Fl(\mathcal{M}))$ par la relation d'ordre induite par l'inclusion de parties de $\Fl(\mathcal{M})$, on trouve des sous-ensembles ordonnés de $\Pf(\Fl(\mathcal{M}))$, stables par réunions quelconques, et aussi par intersections quelconques, avec la seul réserve que dans le cas de $\mathcal{P}$, la famille d'indices doit être $\neq\emptyset$ (puisqu'alors on trouve comme intersection $\Fl(\mathcal{M})$ tout entier, qui n'est pas en général contenu dans $W$). % [sic] la seul réserve -> la seule réserve, feminine agreement

\dgrblock{dgrb133}{Why the admissible classes must be corrected}
\begin{dcomment}{exposition}{p83}{Why the admissible classes must be corrected}
Grothendieck now replaces the moving target $A$ in the stabilizer by the fixed target $W$. The resulting relation is suitable for a Galois connection, but the first choice of its two domains is too restrictive. Formula (6.2) requires the classes on the right to be closed under both pullback and composition. The transporter constructed on page 85 is closed under pullback, but it need not be closed under composition. The correction at (6.7) therefore removes composition from the definition of the right-hand domain and places it in the left-hand domain.

There is also a smaller order-theoretic point on page 84. Classes closed under composition are not in general closed under unions. In the category $0<1<2$, let $L_1$ consist of the identities and $0\to1$, and let $L_2$ consist of the identities and $1\to2$. Both classes are closed under pullback and composition, but their union omits $0\to2$. Thus the assertion about arbitrary unions is false with the initial definition of $\mathcal Q$. Intersections do preserve those properties; the appropriate join is obtained by closing a union under the required operations.

After (6.7), the framework is as follows. Assume $W$ contains isomorphisms and is closed under composition. Let $\mathcal P$ consist of the composition-closed classes $A$ with $\operatorname{Is}(\mathcal M)\subseteq A\subseteq W$. Let $\mathcal Q$ consist of the pullback-stable classes $L$ containing all isomorphisms. Both are complete lattices. Meets are intersections, taking the empty meet to be $W$ in $\mathcal P$ and all arrows in $\mathcal Q$. A join in $\mathcal P$ is the class of finite composites of arrows from the union; it remains in $W$. A nonempty join in $\mathcal Q$ is simply a union, while its empty join is the isomorphisms. These last empty-family conventions are necessary because every admissible class contains the isomorphisms.

For the size of these lattices, fix a universe large enough to contain the chosen collection of arrows, or work in a larger universe than that used for $\mathcal M$. The arguments concern classes invariant under isomorphism of arrows and universal quantification over pullback squares. This also explains the footnote on page 85: before invariance of $W$ has been assumed, choosing one representative of each pullback is insufficient. Quantifying over every cartesian square avoids that choice. Once $W$ contains isomorphisms and is composition-closed, invariance follows, and the familiar pullback-functor notation becomes harmless.
\end{dcomment}


On va vérifier la condition $(\mathrm{D}_0)$ de la page 69, pour $\mathcal{R}$ et les relations d'ordre. Il y a deux vérifications à faire.

1$^{o}$) \texttt{[où est ‘2$^{o}$’ ?]} Donnons nous $A\in\mathcal{P}$, i.e. $A\subseteq W$, et considérons l'ensemble $\mathcal{R}(A)$ des $L\in\mathcal{Q}$ tels que l'on ait $(A,L)\in\mathcal{R}$. Il est clair que c'est un ouvert (et de même, que ${}^{s}\mathcal{R}(L)$ est un ouvert de $\mathcal{Q}$ pour tout $L$ dans $\mathcal{Q}$) - i.e. $L\in\mathcal{R}(A)$, $L'\subseteq L$ implique % ?? uncertain: left superscript before R(L) (also on p. 53) -- italic s or upright s % [sic] ouvert de $\mathcal{Q}$ -> ouvert de $\mathcal{P}$, cf. "l'ouvert ${}^{s}\mathcal{R}(L)$ de $\mathcal{P}$" p. 53
\printedpage{52}
% head: A. Grothendieck : Dérivateurs, Ch. II, Édition des Universités de Montpellier II et Paris VII
% foot: Version 8 juillet 2001
$L'\in\mathcal{R}(A)$. Reste à trouver un plus grand élément de $\mathcal{R}(A)$. Notons que si $L\in\mathcal{R}(A)$, alors pour tout $p\in L$, $p : Y'\longrightarrow Y$, le changement de base
\[
p^{*} : \mathrm{Cat}/Y\longrightarrow\mathrm{Cat}/Y' : X\longmapsto X\times_{Y}Y'
\]
doit transformer flèches $f\in A$ en flèches $g\in W$, comme il résulte du fait que $L$ est stable par changement de base.
\[
\begin{tikzcd}
X_1 \arrow[d, "f"'] & X'_1 \arrow[l, "q_1"'] \arrow[d, "g"] \\
X_2 \arrow[d] & X'_2 \arrow[l, "q_2"'] \arrow[d] \\
Y & Y'. \arrow[l, "p"']
\end{tikzcd}
\]
Donc $L\subseteq\mathbf{L}_W(A)$, où on pose
\msp{page 85}
\[
\begin{aligned}
\mathbf{L}_W(A)=\{(p : Y'\longrightarrow Y)\in\Fl(\mathcal{M})\quad\mid\quad & p^{*} : \mathrm{Cat}/Y\longrightarrow\mathrm{Cat}/Y'\ \text{transforme}\\
& \text{flèches de $A$ en flèches \textit{de} $W$.}\}
\end{aligned}
\]
(\footnote{NB le dire plutôt en termes de diagrammes cartésiens, vu qu'on n'a pas fait d'hypothèse de stabilité de $W$ par isomorphismes.}).
\[
\begin{tikzcd}
X_1 \arrow[d, "f"'] \arrow[dr, phantom, "\text{cart.}"] & X'_1 \arrow[l] \arrow[d, "f'"] \\
X_2 \arrow[d] \arrow[dr, phantom, "\text{cart.}"] & X'_2 \arrow[l] \arrow[d] \\
Y & Y' \arrow[l, "p"']
\end{tikzcd}
\]
\[
f\in A\implies f'\in W.
\]
Il est clair d'autre part que cette partie de $\Fl(\mathcal{M})$ est stable par changement de base, donc $\mathbf{L}_W(A)\in\mathcal{Q}$, et que $(A,\mathbf{L}_W(A))\in\mathcal{R}$. Donc $\mathbf{L}_W(A)$ est un plus grand élément de $\mathcal{R}(A)$, donc on trouve
\[
\mathcal{R}=\{(A,L)\in\mathcal{P}\times\mathcal{Q}\mid L\subseteq\mathbf{L}_W(A)\}.
\tag*{(6.5)}
\]
\printedpage{53}
% head: A. Grothendieck : Dérivateurs, Ch. II, Édition des Universités de Montpellier II et Paris VII
% foot: Version 8 juillet 2001
\[
\begin{tikzcd}
X_1 \arrow[d, "f"'] \arrow[dr, phantom, "\text{cart.}"] & X'_1 \arrow[l, "q_1"'] \arrow[d, "f'"] \\
X_2 \arrow[d] \arrow[dr, phantom, "\text{cart.}"] & X'_2 \arrow[l, "q_2"'] \arrow[d] \\
Y_1 \arrow[d] \arrow[dr, phantom, "\text{cart.}"] & Y'_1 \arrow[l, "p_1"'] \arrow[d] \\
Y & Y'. \arrow[l, "p"']
\end{tikzcd}
\]
Inversement, partons de $L\in\mathcal{Q}$, et montrons que l'ouvert ${}^{s}\mathcal{R}(L)$ de $\mathcal{P}$ a un plus grand élément. Pour ceci, posons
\[
\begin{aligned}
\mathbf{A}_W(L)=\{(f : X\longrightarrow Y)\in W\subseteq\Fl(\mathcal{M})\quad\mid\quad\forall\qquad
&\begin{tikzcd}[ampersand replacement=\&]
X \arrow[d, "f"'] \arrow[dr, phantom, "\textit{cart.}"] \& X' \arrow[l, "q"'] \arrow[d, "f'"] \\
Y \& Y', \arrow[l, "p"]
\end{tikzcd}\\
&\text{avec $p\in L$, on a $f'\in W$}\}.
\end{aligned}
\tag*{(6.6)}
\]
(NB Quand $L$ contient les identités, comme c'est le cas notamment s'il est de la forme $\mathbf{L}_W(A)$ - alors il contient les isomorphismes, du moins si $A$ les contient - alors $\mathbf{A}_W(L)$ est l'ensemble de $f\in\Fl(\mathcal{M})$ satisfaisant la condition énoncée après $|$, sans avoir à spécifier $f\in W$.) On a $\mathbf{A}_W(L)\in\mathcal{P}$, et on vérifie que $(\mathbf{A}_W(L),L)\in\mathcal{R}$ (en utilisant la stabilité de $L$ par composition), enfin il est clair que $(A,L)\in\mathcal{R}$ implique $A\subseteq\mathbf{A}_W(L)$ (en utilisant $A\subseteq W$), donc $\mathbf{A}_W(L)$ est un plus grand élément de ${}^{s}\mathcal{R}(L)$.

\msp{page 86}Notons que si $W$ \textit{est stable par composition} (ce que nous allons supposer désormais), il en est de même de $\mathbf{A}_W(L)$ (ce qui utilise encore le fait que $L$ soit stable par changement de base). De plus, si $W$ \textit{contient les isomorphismes} (ce que nous supposerons également), $\mathbf{A}_W(L)$ les contient aussi. Quant à $\mathbf{L}_W(A)$, il contient les isomorphismes \textit{pourvu que} $A$ soit stable par isomorphismes ; quand $A$ est stable par composition et contient les isomorphismes, cette condition sera automatiquement satisfaite. Nous tenons à cette condition, c'est pourquoi du côté de $A$, dans la définition de $\mathcal{P}$, je vais introduire cette condition supplémentaire que $A$ soit \textit{stable par isomorphismes.} Donc je pose, en corrigeant (6.2),
\[
\left\{
\begin{array}{rcl}
\mathcal{P} & = & \{A\in\Pf(\Fl(\mathcal{M}))\quad\mid\quad\mathrm{Is}(\mathcal{M})\subseteq A\subseteq W,\quad A\text{ stable par composition}\}\\[1ex]
\mathcal{Q} & = & \{L\in\Pf(\Fl(\mathcal{M}))\quad\mid\quad\mathrm{Is}(\mathcal{M})\subseteq L,\quad L\text{ stable par changement de base}\}.
\end{array}
\right.
\tag*{(6.7)}
\]

\dgrblock{dgrb134}{The transporter and its adjoint}
\begin{dcomment}{exposition}{p85}{The transporter and its adjoint}
The corrected construction works because a larger testing class produces a smaller transporter. It is useful to give the complete argument, including the distinction between testing over one codomain and testing all arrows in a slice.

For $A\in\mathcal P$, define $T_W(A)$ to consist of the maps $p:Y'\to Y$ such that every arrow $f:X_1\to X_2$ in $A$, equipped with a map $X_2\to Y$, has $p^*f\in W$. Pullback pasting shows that $T_W(A)$ is stable under arbitrary base change: a test after pulling $p$ back along $Z\to Y$ is already a test for $p$ with its structure map factoring through $Z$. It contains the isomorphisms because $A\subseteq W$ and $W$ is invariant under isomorphism. Thus $T_W(A)\in\mathcal Q$.

For $L\in\mathcal Q$, define $N_W(L)$ to consist of the maps $f\in W$ whose pullbacks along all maps of $L$ into their codomain belong to $W$. The explicit requirement $f\in W$ is redundant here, since identities belong to $L$. Isomorphisms are in $N_W(L)$. For composable $f,g\in N_W(L)$ and $p\in L$ into the codomain of $g$, factor $p^*(gf)$ as the pullback of $f$ along $g^*p$, followed by $p^*g$. Pullback stability gives $g^*p\in L$, so both factors lie in $W$. Hence $N_W(L)$ is composition-closed and belongs to $\mathcal P$.

The defining relation is now exactly
\[
 A\subseteq N_W(L)\quad\Longleftrightarrow\quad L\subseteq T_W(A).
\]
For the forward direction, test $p\in L$ against $f\in A$ over $Y$. The base change $X_2\times_Y Y'\to X_2$ is in $L$, so the assumption on $f$ applies. Conversely, if $p$ is a map of $L$ into the actual codomain of $f$, this is a special case of the transporter test. This proves both implications and the maximality assertions on pages 85--86. Closure of $L$ under composition is unnecessary for this adjunction; its appearance in the text's proof of $(N_W(L),L)\in\mathcal R$ is redundant.

Both operations reverse inclusions, so the usual polarity argument now applies without qualification. Their composites are monotone, extensive and idempotent; their fixed classes correspond anti-isomorphically. The extreme pairs in (6.9)--(6.10) follow directly: $T_W(\operatorname{Is})$ is all arrows, $N_W(\operatorname{Fl})=W_{\mathrm{univ}}$, $N_W(\operatorname{Is})=W$, and $T_W(W)=\mathrm{Fib}_W$. In particular, every closed left class contains all universal $W$-equivalences, and every closed right class contains $\mathrm{Fib}_W$. This is the relation-theoretic content of \S6, established with only the hypotheses actually imposed there.
\end{dcomment}


% ---- from Der-ChII-234_p056-060.tex
\printedpage{54}
% head: A. Grothendieck : Dérivateurs, Ch. II, Édition des Universités de Montpellier II et Paris VII
On est bien sûr intéressé surtout aux éléments clos de $\mathcal{P}$, $\mathcal{Q}$, i.e. les parties $A$, $L$ de $\Fl(\mathcal{M})$, contenant les isomorphismes, stables par composition (et $L$ stable aussi par changement de base, $A \subseteq W$) telles que l'on ait
\[
A = \mathbf{A}_W(L), \qquad L = \mathbf{L}(A) \tag*{(6.8)}
\]
\texttt{[plutôt }$\mathbf{L}_W(A)$\texttt{]}.

\msp{page 87}

Je garderai les notations $\natural$, $\flat$ du paragraphe précédent. Quels sont les deux couples extrêmes ?

Partant de $A_0 = \Is(\mathcal{M})$ (plus petit des éléments de $\mathcal{P}$), on trouve $\mathbf{L}_W(A_0) = \Fl(\mathcal{M})$ (plus grand élément de $\mathcal{Q}$), et $A_0^{\natural} = \mathbf{A}_W(\Fl(\mathcal{M})) = W_{\mathrm{univ}}$, ensemble des flèches qui sont dans $W$ et le restent par tout changement de base :
\[
A_0^{\natural} = W_{\mathrm{univ}}, \qquad \beta_0(A_0) \overset{\text{déf}}{=} \mathbf{L}_W(A_0) = \Fl(\mathcal{M}), \qquad \text{\scriptsize où $A_0 = \Is(\mathcal{M})$ plus petit élément de $\mathcal{P}$.} \tag*{(6.9)}
\]
Inversement, partons de $L_0 = \Is(\mathcal{M})$, le plus petit élément de $\mathcal{Q}$. On a $\mathbf{A}_W(L_0) = W$ ($W$ étant stable par composition et contenant les isomorphismes), et $L_0^{\flat} = \mathbf{L}_W(W) = \Fib_W$, les $W$-fibrations :
\[
L_0^{\flat} = \Fib_W, \qquad \alpha_0(L_0) = W, \qquad \text{\scriptsize où $L_0 = \Is(\mathcal{M})$, plus petit élément de $\mathcal{Q}$.} \tag*{(6.10)}
\]
Donc les parties $\mathcal{R}$-closes à gauche de $\Fl(\mathcal{M})$ se promènent entre $W_{\mathrm{univ}}$ et $W$, les parties $\mathcal{R}$-closes à droite entre $\Fib_W$ et $\Fl(\mathcal{M})$.

\textbf{Proposition 6.1.} \textit{Soient $A \in \mathcal{P}$, $L \in \mathcal{Q}$ tels que $A = \alpha_0(L)$, $L = \beta_0(A)$ (i.e. $A = \mathbf{A}_W(L)$, $L = \mathbf{L}_W(A)$). Pour que $L$ soit stable par composition, il faut et il suffit que $A$ soit stable par changement de base par les flèches $f \in L$. (Plus généralement, si $L \in \mathcal{Q}$ est un élément quelconque de $\mathcal{Q}$, et si on suppose de} \msp{page 88} \textit{plus que $L$ est stable par composition, alors $A = \mathbf{A}_W(L)$ est stable par changement de base par $L$.)}

\dgrblock{dgrb135}{When a transporter becomes a stabilizer}
\begin{dcomment}{exposition}{p87:Pour que $L$ soit stable par composition, il faut et il suffit}{When a transporter becomes a stabilizer}
Proposition 6.1 identifies the extra condition that makes the formal correspondence compatible with iterated changes of base. Suppose $(A,L)$ is a closed pair, so $A=N_W(L)$ and $L=T_W(A)$. If $L$ is closed under composition, then $A$ is stable under base change along $L$: for $f\in A$ and $p\in L$, its pullback belongs to $W$, and each further test along $q\in L$ is a test of $f$ along $pq\in L$. Therefore $p^*f\in A$. When $f$ is an arrow over the target of $p$, use the pullback of $p$ along its codomain map; it belongs to $L$ by pullback stability. This proves the more general assertion in parentheses, since it uses only $A=N_W(L)$.

Conversely, assume $A$ is stable under base change along $L$. Given composable maps $q:Z\to Y'$ and $p:Y'\to Y$ in $L$, and an arrow $f\in A$ over $Y$, first pull back along $p$. The result is in $A$, and its pullback along $q$ is again in $A\subseteq W$. Thus $(pq)^*f\in W$, whence $pq\in T_W(A)=L$. This proves the proposition completely.

For a closed pair one can sharpen its wording without adding a hypothesis:
\[
 L\text{ is composition-closed}
 \quad\Longleftrightarrow\quad
 L=\operatorname{Stab}(A).
\]
Indeed, stability of $A$ under $L$ gives $L\subseteq\operatorname{Stab}(A)$, while the reverse inclusion always follows from $A\subseteq W$ and $L=T_W(A)$. Conversely every stabilizer is composition-closed. Thus the two names designate the same class exactly at the closed pairs satisfying Proposition 6.1. Without this extra condition the transporter may contain maps that preserve the $W$-tests but fail to preserve $A$ itself.

The word ``i.e.'' before (6.8) must accordingly be read with care. Closed pairs do not automatically have a composition-closed right member. Grothendieck is particularly interested in those that do; Proposition 6.1 states the criterion. The following finite example verifies that the distinction is real under the abstract hypotheses of this section.
\end{dcomment}


\dgrblock{dgrb136}{A closed transporter which does not compose}
\begin{dcomment}{development}{p87:Pour que $L$ soit stable par composition, il faut et il suffit}{A closed transporter which does not compose}
The formal correction of $\mathcal Q$ in (6.7) becomes necessary in a four-object category. The example also separates the universal quantifier in a transporter from the stronger demand that it stabilize its source class.

\textbf{A closed pair with a non-composition-closed right member.}\ours{}
Let $\mathcal M$ be the chain $0<1<2<3$, with pullbacks given by minima. Write $ij$ for its unique arrow $i\to j$ and $I$ for its identities. Set
\[
 W=I\cup\{02,03\},\qquad A=I\cup\{03\}.
\]
Then $W$ and $A$ contain the isomorphisms and are closed under composition. Moreover,
\[
 T_W(A)=\operatorname{Fl}(\mathcal M)\setminus\{13\},
 \qquad N_W(T_W(A))=A.
\]
Thus $(A,T_W(A))$ is a closed pair, but $12$ and $23$ belong to its right member while their composite $13$ does not.

To verify the transporter formula, the only nonidentity arrow of $A$ to test is $03$. It can be an arrow over $y$ only when $y=3$. Pullback along $x3$ gives $0x$, which lies in $W$ precisely for $x=0,2,3$. Therefore $13$ is the unique excluded map; all identity tests succeed. This proves the first formula, including all slice tests.

For the kernel formula, the identities and $03$ survive by the adjunction. The only remaining member of $W$ is $02$. Its pullback along $12$, a member of the transporter, is $01$, which is not in $W$. Hence $02$ does not survive, and the second formula follows. Finally the right member is closed as well: applying the triple-operation identity for the established Galois connection gives $T_WN_WT_W(A)=T_W(A)$. This completes all the claims.

Here $W$ is an abstract class of arrows with exactly the composition and isomorphism assumptions of \S6; it is not asserted to be a fundamental localizer on $\mathrm{Cat}$, or even to satisfy two-out-of-three. We checked the finite calculation exhaustively with the supplied script and proved it above. We did not find this specific example in the checked transporter and smooth-functor sources. It serves to explain the actual generality of Proposition 6.1, without making a claim about the special localizers used elsewhere in the manuscript.
\end{dcomment}


\dgrblock{dgrb137}{Intrinsic fibrations and sharp maps}
\begin{dcomment}{context}{p87}{Intrinsic fibrations and sharp maps}
Grothendieck wants the admissible changes of base to be determined by the weak equivalences themselves. This is why the extreme pair in (6.10) matters more than an isolated exercise with closure operations. Once $W$ has been fixed, a morphism $p:E\to B$ belongs to $\mathrm{Fib}_W$ precisely when
\[
 p^*:\mathcal M/B\longrightarrow\mathcal M/E
\]
preserves the arrows whose underlying morphisms belong to $W$. Both the test arrow and its map to $B$ vary. Restricting the test to arrows with target $B$ would omit the further base changes built into the definition.

The letter to Thomason of 2 April 1991 makes this intention explicit. In its first numbered discussion, Grothendieck proposes replacing a suitable proper fibration structure by the class determined by this preservation property; existence of enough such fibrations remains a separate factorization requirement \citep[\S1, printed p.~4]{Grothendieck1991Letter}. Thus the canonical class is available before one knows that it supports a useful model structure. The historical point is the choice of what should be intrinsic, not a rejection of the computational value of lifting and factorization axioms.

In a model category the same preservation property is the definition of a \emph{sharp map} in \citet[\S2]{Rezk1998Sharp}. In a right proper model category every model fibration is sharp, and sharpness is equivalent to every strict pullback square along the map being homotopy cartesian \citep[Propositions~2.2 and~2.7]{Rezk1998Sharp}. These statements explain both the usefulness and the limitations of the comparison: Grothendieck's $\mathrm{Fib}_W$ asks for a property of pullback with respect to $W$; a chosen Quillen fibration class supplies maps with that property under the additional hypothesis of right properness. An arbitrary transporter $\mathbf L_W(A)$ only promises to send the smaller class $A$ into $W$. A second pullback can therefore leave the class of arrows on which that promise applies. The extra stability in Proposition~6.1 is exactly what permits the second step.
\end{dcomment}


Terminologie : pour $A \in \mathcal{P}$, i.e. $A$ une partie de $\Fl(\mathcal{M})$ contenue dans $W$ contenant les isomorphismes et stable par composition, $\mathbf{L}_W(A)$ s'appelle le \textit{transporteur de $A$ dans $W$}. (Il y a une notion duale de \textit{cotransporteur.}) C'est un élement de $\mathcal{Q}$, donc une partie de $\Fl(\mathcal{M})$ contenant les isomorphismes et stable par changement de base. D'autre part, pour $L$ dans $\mathcal{Q}$, les $f \in \mathbf{A}_W(L)$ s'appellent flèches \textit{$W$-$L$-asphériques,} ou les \textit{$W$-$L$-équivalences,} ou les \textit{$W$-équivalences $L$-universelles,} et $\mathbf{A}_W(L)$ s'appelle le \textit{$L$-noyau de $W$}. On pourrait le noter plus simplement $W_L$ (quand cette notation ne conflicte pas avec $W_S \subseteq \Fl(\mathcal{M}/S)$ pour $S \in \mathrm{Ob}\,\mathcal{M}$). C'est une partie de $W$ qui (comme $W$ lui-même) contient les isomorphismes et est stable par composition. % [sic] élement -> élément, missing accent

Enfin, il y a lieu aussi d'introduire le \textit{stabiliseur} $\mathrm{Stab}(A)$, ou transporteur de $A$ dans $A$, qui ne dépend pas de la donné de $W$, mais seulement de la partie $A$ de $\Fl(\mathcal{M})$, contenant les isomorphismes et stable par composition. On pourrait aussi le noter $\Fib_A$, et appeler ses éléments des \textit{$A$-fibrations} ($A$ jouant à présent le rôle de $W$). On a $\mathrm{Stab}(A) \subseteq \mathbf{L}_W(A)$, et il semblerait qu'en % [sic] donné -> donnée, feminine noun
% foot: Version 8 juillet 2001
\printedpage{55}
% head: A. Grothendieck : Dérivateurs, Ch. II, Édition des Universités de Montpellier II et Paris VII
\msp{page 89}
pratique, ce sont ces stabiliseurs, plus que les transporteurs $\mathbf{L}_W(A)$, qui sont importants (quand les deux ne coïncident pas). C'est ainsi que dans Cat, si $W$ est un localiseur fondamental, les morphismes $W$-lisses (resp. $W$-propres) $\mathrm{Liss}_W$ ($\mathrm{Prop}_W$), qui sont les stabiliseurs des ensembles $W_{\mathrm{as}}$ ($W_{\mathrm{coas}}$) des ensembles de morphismes $W$-asphériques (resp. $W$-coasphériques), jouant un rôle essentiel, alors que je ne sais pas caractériser de façon sensible les flèches des transporteurs $\mathbf{L}_W(W_{\mathrm{as}})$, $\mathbf{L}_W(W_{\mathrm{coas}})$, lesquels contiennent, en plus des morphismes $W$-lisses (resp. $W$-propres) les $W$-fibrations $\Fib_W$. Je ne sais pas comment caractériser une classe de morphismes qui apparaît aussi hétéroclite! De même, il y a lieu d'introduire et de caractériser le stabiliseur de $W_{\mathrm{bias}}$ (ensemble des flèches $W$-biasphériques), comme les morphismes `$W$-nets' $\mathrm{Net}_W$, par une description similaire à celle des $\mathrm{Liss}_W$, $\mathrm{Prop}_W$, utilisant les bilocalisés $\backslash\underset{u}{X}/$, $\backslash\underset{v}{Y}/$ , au lieu des localisés $X/x$, $Y/y$ ou des colocalisés $x\backslash X$, $y\backslash Y$. Il y a lieu enfin de déterminer le stabiliseur de l'ensemble Equ des Cat-équivalences (ce qui ne dépend pas de la donné d'un $W$ - mais on a $\mathrm{Equ} \subseteq \mathcal{P}(W)$ pour tout $W$ fondamental). % [sic] donné -> donnée, feminine noun
% foot: Version 8 juillet 2001

\dgrblock{dgrb138}{The scope of the remaining transporter questions}
\begin{dcomment}{status}{p89:Il y a lieu enfin de déterminer le stabiliseur}{The scope of the remaining transporter questions}
The questions at the end of the section ask for usable local tests for classes already defined by universal preservation properties. For smooth maps, that programme was developed explicitly: \citet[Definition 4.6, Theorem 4.13, and Example 4.17]{Maltsiniotis2005Asphericite} defines preservation of asphericity under pullback, gives equivalent local conditions, and compares the result with smoothness for a fundamental localizer. Properness is obtained by opposite categories. This identifies a later theorem corresponding to the stabilizer side of Grothendieck's discussion.

The transporter $T_W(W_{\mathrm{asph}})$ asks only that pullbacks of aspherical maps be $W$-equivalences; its target condition is weaker than preservation of asphericity. The inclusions on page 89 follow directly: it contains $\operatorname{Stab}(W_{\mathrm{asph}})$ because $W_{\mathrm{asph}}\subseteq W$, and it contains $\operatorname{Stab}(W)$ because the latter preserves every arrow of $W$. The corresponding statements for coasphericity are dual. They do not identify a union of those classes, nor imply that the transporter is composition-closed.

We checked Maltsiniotis's asphericity paper and searched the literature under the manuscript's transporter, stabilizer and ``net'' terminology. We did not locate a theorem identifying the general biaspherical stabilizer with precisely the proposed bilocal lifting condition, or a full intrinsic description of both mixed transporters. This is a statement about the scope of our verification, not a claim that those problems remain open in every formulation. The separate question about the stabilizer of categorical equivalences has a definite answer, proved in the block ``The stabilizer of categorical equivalences''. The history of the intrinsic class $\mathrm{Fib}_W$ is discussed in ``Intrinsic fibrations and sharp maps''.
\end{dcomment}


\printedpage{56}
% head: A. Grothendieck : Dérivateurs, Ch. II, Édition des Universités de Montpellier II et Paris VII
\msp{page 90}

\dgrblock{dgrb139}{The stabilizer of categorical equivalences}
\begin{dcomment}{later}{p89}{The stabilizer of categorical equivalences}
The last question on this page has a concrete answer. The obstruction to preserving an equivalence under a strict pullback is an isomorphism of objects in the base which cannot be lifted. The relevant maps are the \emph{isofibrations}, the fibrations of the usual model structure on categories; see \citet[\S2 and Theorem~3.1]{Rezk1996Categories}. We give the argument in the notation of the manuscript, so that the answer also specifies the meaning of its stabilizer.

\textbf{Proposition.} Let $\mathrm{Equ}$ be the class of equivalences of small categories. A functor $p:E\to B$ belongs to $\mathrm{Stab}(\mathrm{Equ})$ if and only if, for every $e\in E$ and every isomorphism $\gamma:p(e)\to b$ in $B$, there are an object $e'$ and an isomorphism $\widetilde\gamma:e\to e'$ with $p(e')=b$ and $p(\widetilde\gamma)=\gamma$.

\emph{Proof.} Suppose first that this lifting property holds. Let $F:C\to D$ be an equivalence over $B$, with structure map $h:D\to B$. The pullback
\[
 F_E:C\times_B E\longrightarrow D\times_B E
\]
is fully faithful: a morphism in either pullback is a pair of morphisms with equal images in $B$, and full faithfulness of $F$ lifts the first member uniquely without changing that equality. To prove essential surjectivity, take $(d,e)$ in the target. Choose $c\in C$ and an isomorphism $\theta:F(c)\to d$. Apply the isomorphism-lifting property to $(h\theta)^{-1}:p(e)\to hF(c)$ and invert the resulting lift. This gives an isomorphism $\eta:e'\to e$ with $p(e')=hF(c)$ and $p(\eta)=h\theta$. The pair $(\theta,\eta)$ is an isomorphism $F_E(c,e')\to(d,e)$. Thus $F_E$ is an equivalence, as required for the stabilizer.

Conversely, suppose $p\in\mathrm{Stab}(\mathrm{Equ})$. Let $J$ be the groupoid with two objects $0,1$ and exactly one morphism between each ordered pair. The data $(e,\gamma)$ define a functor $J\to B$ taking $0$ to $p(e)$, $1$ to $b$, and $0\to1$ to $\gamma$. The inclusion $\{1\}\to J$ is an equivalence over $B$. Hence
\[
 E\times_B\{1\}\longrightarrow E\times_BJ
\]
is an equivalence. Essential surjectivity at $(e,0)$ supplies $(e',1)$ and an isomorphism between it and $(e,0)$. Orienting that isomorphism from $(e,0)$ to $(e',1)$ gives the required $\widetilde\gamma$. This proves the converse.\hfill$\square$

The proof also explains why the answer is sensitive to strict pullbacks. A point inclusion into $J$ is itself an equivalence, but pulling it back along the inclusion of the other point gives $\emptyset\to e$. An isofibration supplies the missing representatives over isomorphic objects. No assertion about the fundamental localizer $W$ is needed for this calculation. In particular, it answers the question about $\mathrm{Equ}$ without identifying the more delicate stabilizer of $W$-biaspherical maps with a familiar fibration class.
\end{dcomment}


\section*{7\quad Morphismes cofiltrants et localement cofiltrants}

Mais je ne veux pas m'attarder à présent sur ces développements, qui ne me paraissent pas urgents, et voudrais revenir aux questions liées à la notion de catégories filtrantes ou cofiltrantes.

\dgrblock{dgrb140}{The conditions that pass through a comma category}
\begin{dcomment}{exposition}{p90:Mais je ne veux pas m'attarder}{The conditions that pass through a comma category}
Grothendieck returns to filteredness because it supplies concrete conditions under which the cofinality arguments can be repeated after restricting a diagram. The adjective ``local'' has a precise force here. It does not mean merely that the set of objects in each connected component is directed. The commutativity of the triangles must survive the construction of a common predecessor.

Write $P(C)$ for the following condition: every cospan $a\xrightarrow{r}c\xleftarrow{s}b$ admits arrows $u:d\to a$, $v:d\to b$ such that $ru=sv$. Write $E(C)$ for the equalization condition: every parallel pair $r,s:a\to b$ admits $t:d\to a$ with $rt=st$. Finally, write $I(C)$ for nonemptiness together with the existence of a common predecessor for every pair of objects. These are respectively $\mathrm{PS}_g1$, $\mathrm{PS}_g2$, and $\mathrm{PS}_g1'$ in the text. Then $C$ is cofiltered precisely when $I(C)$ and $E(C)$ hold, and locally cofiltered precisely when $P(C)$ and $E(C)$ hold. The latter condition allows the empty category. Equivalently, every connected component is cofiltered; the elementary decomposition underlying this statement is recorded, with the opposite variance, in \citet[Lemmas 4.19.6 and 4.19.8]{StacksFiltered2026}.

The weaker terms need an explicit convention. In the working proof, ``weakly cofiltered'' means $I(C)$, whereas ``locally weakly cofiltered'' means $P(C)$. Footnote 64 proposes reserving the former term for nonempty connectedness together with $P(C)$. These conditions differ: the category with two parallel arrows $0\rightrightarrows1$ satisfies $I$ but fails $P$. The closing diagrams recognize the distinction and call $I$ irreducibility. We therefore retain $I$ and $P$ in statements where the terminology changes. Proposition 7.3 and its composition corollary hold with either weak convention, provided it is used consistently.

Here is the argument behind the proposition, including its three independent assertions. Suppose first that $P$ holds in $Y$ and in every $y\backslash X$. Given a cospan in $X$, complete its image in $Y$ to a commutative square with lower vertex $y$. The three objects of the cospan, equipped with the resulting maps from $y$, form a cospan in $y\backslash X$. Applying $P$ there gives the required square in $X$. If instead $I$ holds in $Y$ and in the comma categories, two objects of $X$ first acquire a common predecessor for their images in $Y$, and then acquire one in the corresponding comma category. Nonemptiness follows by choosing $y\in Y$ and an object of $y\backslash X$.

For equalization, start with $r,s:a\to b$ in $X$. Choose $v:y\to f(a)$ with $f(r)v=f(s)v$, using $E(Y)$. Then $r,s$ are parallel arrows in $y\backslash X$. An equalizing arrow there equalizes them in $X$. These arguments prove the cofiltered and locally cofiltered cases as well. For the revised weak convention, $P(Y)$ and connectedness imply $I(Y)$: the relation ``has a common predecessor'' is transitive by $P$, and contains every arrow and its reverse. The same observation applies to the comma categories, so the preceding argument also proves connectedness of $X$.

The composition statement now has a transparent reason. A discrete opfibration $Y'\to Y$ identifies $y'\backslash(X\times_Y Y')$ with $y'_{Y}\backslash X$, by unique lifting of arrows out of $y'$. This proves the co\'etale base-change assertion of Remark 7.2. Apply it to $z\backslash Y\to Y$, and then apply Proposition 7.3 to $z\backslash X\to z\backslash Y$. The printed assertion that \emph{\'etale} maps are locally ``filtrants'' has the variance reversed; Lemma 7.5 gives the correct locally cofiltered assertion below.
\end{dcomment}


\textbf{Definition 7.1.} Soit $f : X \longrightarrow Y$ une flèche dans Cat. On dit que $f$ est \textit{cofiltrante} (resp. \textit{localement cofiltrante}, resp. \textit{faiblement cofiltrante} (\footnote[64]{J'ai modifié la terminologie : maintenant faiblement cofiltrant $=$ 0-connexe $+$ localement faiblement cofiltrant, \textit{stable par passage} à $X'$ sur $X$, $X'$ étale et $\infty$-connexe. Revoir la terminologie dans le cas faiblement cofiltrant, qui est une notion traîtresse, faute d'être stable par passage \texttt{[à]} un $X'$ étale sur $X$.}), resp. \textit{localement faiblement cofiltrante} ou encore \textit{précofiltrante}) si pour tout $y$ dans $Y$, la catégorie colocalisé $y\backslash X$ est cofiltrante (resp. \dots). (Dualement, on dit que $f$ est \textit{filtrante}, resp. \dots, si les catégories $X/y$ sont filtrantes, resp. \dots, i.e. si $f^o : X^o \longrightarrow Y^o$ est cofiltrante resp. \dots.) % [sic] Definition -> Définition, missing accent; colocalisé -> colocalisée, agreement with catégorie

\textbf{Remarque 7.2.} Notions stables par changement de base coétale. Contient les isomorphismes et même les Cat-équivalences. Les morphismes \textit{étales} sont localement filtrants.

\textbf{Proposition 7.3.} \textit{Soit $f : X \longrightarrow Y$ un morphisme cofiltrant (resp. localement cofiltrant, resp. faiblement cofiltrant, resp. précofiltrant), et supposons que $Y$ soit cofiltrant (resp. localement cofiltrant, resp. faiblement cofiltrant, resp. précofiltrant). Alors $X$ lui-même est cofiltrant (resp. \dots).}

\textsc{Démonstration}. Il s'agit, dans les quatre cas envisagés, de vérifier pour $X$ la condition $\underbrace{\mathrm{PS}_g1' \;/\; \mathrm{PS}_g1}_{\text{deux derniers cas}} \;/\; \mathrm{PS}_g1' + \mathrm{PS}_g2 \;/\; \mathrm{PS}_g1 + \mathrm{PS}_g2$. Mais \msp{page 91} pour chacune des trois conditions $\mathrm{PS}_g1$, $\mathrm{PS}_g1'$, $\mathrm{PS}_g2$ séparément, on voit immédiatement que si elle est vérifiée pour $Y$ et pour les $y\backslash X$, elle l'est pour $X$, d'où la proposition.

\textbf{Corollaire 7.3.} \textit{Les quatre notions envisagées dans (7.1) sont stables par composition.}

\textsc{Démonstration} immédiate, en utilisant la remarque (7.2).
\[
\begin{tikzcd}
X \arrow[d] & z\backslash X \arrow[l] \arrow[d] \\
Y \arrow[d] & z\backslash Y \arrow[l] \arrow[d] \\
Z & z\backslash Z. \arrow[l]
\end{tikzcd}
\]
% foot: Version 8 juillet 2001

\printedpage{57}
% head: A. Grothendieck : Dérivateurs, Ch. II, Édition des Universités de Montpellier II et Paris VII
\textbf{Remarque 7.4.}
\begin{enumerate}
\item[1)] Notons que la notion de morphisme cofiltrant ou faiblement cofiltrant n'est \textit{pas} stable par changement de base quelconque, ni même par une immersion ouverte. Exemple : $e \longrightarrow \boldsymbol{\Delta}^1$, appliquant $e$ en 1, est cofiltrant,
\[
\begin{tikzcd}
 & e \arrow[d, mapsto] \\
0 \arrow[r] & 1
\end{tikzcd}
\]
mais sa fibre en 0 est vide, donc n'est ni cofiltrante, ni faiblement cofiltrante. Je n'ai pas vu d'exemple montrant que la notion de flèche \textit{localement} cofiltrante resp. précofiltrante n'est pas stable par changement de base, et ne voudrais pas m'y attarder, n'ayant pas besoin de connaître la réponse.
\item[2)] Pour que $X$ soit cofiltrante (resp. \dots), il faut et il suffit que $X \longrightarrow e$ le soit. Ainsi la notion absolue est ramenée à une notion relative.
\item[3)] Si $f : X \longrightarrow Y$ est cofiltrante, elle est \msp{page 92} $W$-coasphérique (donc $\in W$) pour tout $W$ satisfaisant l'axiome Loc (8) des limites inductives. En particulier, elle est $W_0$-asphérique, i.e. un foncteur cofiltrant (et même seulement faiblement cofiltrant) est \textit{cofinal.}
\end{enumerate}
La raison principale pour m'intéresser aux flèches cofiltrantes est contenue dans la

\dgrblock{dgrb141}{A finite answer to the base-change question}
\begin{dcomment}{development}{p91}{A finite answer to the base-change question}
Grothendieck's example separates cofilteredness from its local form only by making a fibre empty. To answer the question he leaves beside it, the pullback must instead create a connected comma category whose two branches have lost their common predecessor.

\noindent\textbf{A four-object counterexample.}\ours\par
The class of locally cofiltered functors, and the class of locally weakly cofiltered functors, are not stable under arbitrary base change, even among finite posets.

\emph{Proof.} Let $f:\{1\}\hookrightarrow[1]$ select the final object of the walking arrow. Both $0\backslash\{1\}$ and $1\backslash\{1\}$ are singleton categories, so $f$ is cofiltered. Let $D$ be the diamond poset
\[
\begin{tikzcd}[row sep=small,column sep=small]
 & b\arrow[dr] & \\
a\arrow[ur]\arrow[dr] & & d\\
 & c\arrow[ur] &
\end{tikzcd}
\]
and define $v:D\to[1]$ by $v(a)=0$ and $v(b)=v(c)=v(d)=1$. The pullback of $f$ is the full inclusion $j:U\hookrightarrow D$, where $U=\{b,c,d\}$ and $b<d$, $c<d$. Since $a$ is initial in $D$, there is an isomorphism $a\backslash U\cong U$. This category is connected, but the cospan $b\to d\leftarrow c$ has no common predecessor in $U$. Thus $P(a\backslash U)$ fails; in particular this comma category is not locally cofiltered. The same pullback disproves both proposed stability assertions.\hfill$\square$

This example answers the question for strict categorical pullbacks, which are the pullbacks used in the source. It does not conflict with Remark 7.2: $v$ is not a discrete opfibration, since the arrow $0\to1$ has several lifts with source $a$. Nor does it conflict with Proposition 7.4, which concerns an \emph{\'etale change of the source} followed by a factorization through an \emph{\'etale target}.

The same diamond makes the variance slip in Remark 7.2 visible. The lower open subposet $V=\{a,b,c\}\subset D$ defines a discrete fibration. Its comma category $V/d$ is $V$, which is connected and not filtered: $b$ and $c$ have no common upper bound there. Thus an \emph{\'etale} functor need not be locally filtered.

We checked the discussions of representable flatness in \citet[Remark 2.4]{Tendas2024} and \citet{Shulman2011Flat}, and the explicit filteredness criteria in \citet[Section 4.19]{StacksFiltered2026}. We did not locate this finite pullback construction there. The label records our supplied answer to the manuscript's question, without claiming priority for the example.
\end{dcomment}


\dgrblock{dgrb142}{The image that belongs in Proposition 7.4}
\begin{dcomment}{exposition}{p92:La raison principale pour m'intéresser}{The image that belongs in Proposition 7.4}
The proposition is intended to compare sections or cosections after passing to an open part of a slice. The target must therefore be the part reached by that open, rather than the part reached by the whole slice. In (c) the printed definition of $V$ omits the restriction to $U$. Its correct form is
\[
 V=\{z\in Y/f(x)\mid z\backslash U\ne\varnothing\}.
\]
Equivalently, $V$ is the lower open subcategory generated by $(f/x)(U)$. This is already forced by $U=\varnothing$: the desired target is then empty, whereas the target defined in print is all of $Y/f(x)$, because $(f/x)(x)=f(x)$ is its terminal object. More generally, the proof's assertion that $z\backslash U$ is nonempty ``par construction de $V$'' uses exactly the corrected definition. This is a slip in the edition's text, not a mistranscription; the correction is consequently explained here and retained as a recorded slip in the source.

Once this target is specified, the two comparison formulas have their expected variance. With $g:U\to V$, a covariant diagram $G:V\to\mathcal M$ has
\[
\mathop{\rm colim}_{U}Gg\ \cong\ \mathop{\rm colim}_{V}G,
\]
whereas a presheaf $F:V^{\rm op}\to\mathcal M$ has
\[
\mathop{\rm lim}_{V^{\rm op}}F\ \cong\
\mathop{\rm lim}_{U^{\rm op}}Fg^{\rm op}.
\]
Both statements follow from the same nonempty-connected comma criterion for finality. The second does not assert invariance of limits of \emph{covariant} diagrams under a final functor. The needed limits or colimits must exist in $\mathcal M$. The letter printed as $q$ in (c) is legible in the edition; it denotes the map called $g$ in the proof and in these formulas.

The homotopical qualification also has a precise scope. The cofiltered comma categories are $W$-aspherical under the stated axiom Loc~(8), so $g$ is $W$-coaspherical. Weak cofilteredness by itself gives connectedness, hence ordinary finality, but supplies no general assertion of contractibility. Footnote 65 correctly reserves the stronger conclusion for the cofiltered hypothesis. The phrase ``$W_0$-asph\'erique'' in Remark 7.4(3) must read ``$W_0$-coasph\'erique'' in the convention of this chapter.
\end{dcomment}


\textbf{Proposition 7.4.} \textit{Soit $f : X \longrightarrow Y$ dans} Cat \textit{localement cofiltrant (resp. localement faiblement cofiltrant). Alors :}
\begin{enumerate}
\item[\textit{a)}] \textit{Pour tout $X'$ étale sur $X$, le composé $f' : X' \longrightarrow X \longrightarrow Y$ est localement cofiltrant (resp. \dots), et pour toute factorisation de $f'$ en $X' \xrightarrow{g'} Y' \longrightarrow Y$ avec $Y' \longrightarrow Y$ étale, $g'$ est localement cofiltrant (resp. \dots).}
\item[\textit{b)}] \textit{$\forall\ x$ dans $X$, le morphisme sur les localisés}
\[
f/x : X/x \longrightarrow Y/x \qquad (y = f(x)) % [sic] Y/x -> Y/y, since y = f(x)
\]
\textit{est cofiltrant (resp. faiblement cofiltrant).}
\item[\textit{c)}] \textit{Si $x$, $y$ comme dans b), et $U$ est un ouvert de $X/x$, désignant par $V$ l'ouvert de $Y/y$ engendré par l'image de $f/x$ (i.e. l'ouvert formé des $z$ dans $Y/y$ tels que $z\backslash(X/x)$ soit non vide), la flèche induite $q : U \longrightarrow V$} \texttt{[est]} \textit{cofiltrante (resp. faiblement cofiltrante), a fortiori cofinale (et même $W$-coasphérique, pour tout $W$ satisfaisant l'axiome des limites Loc (8))} (\footnote[65]{Il suffit ici qu'on soit dans l'hypothèse respée ($f$ \textit{faiblement} localement cofiltrant) sauf pour le NB ici \texttt{[`le NB ici' est `(et même ...Loc (8))']}. % [sic] respée -> intended word unclear (perhaps « respectée » or « resp. »), not a French word
}). % doubt resolved: PDF prints q here; later g is an edition naming inconsistency (recorded in corrections.csv).

\msp{page 93}
\textit{Donc pout tout} cofaisceau \textit{$G$ sur $V$, à valeurs dans une catégorie $\mathcal{M}$ avec $\varinjlim$ , on a} \texttt{[un]} \textit{isomorphisme} % [sic] pout -> pour, misprint
\[
\varinjlim_U g^*(G) \xrightarrow{\ \sim\ } \varinjlim_V G,
\]
% foot: Version 8 juillet 2001
\printedpage{58}
% head: A. Grothendieck : Dérivateurs, Ch. II, Édition des Universités de Montpellier II et Paris VII
\textit{et pour tout préfaisceau $F$ sur $V$, à valeurs dans une catégorie $\mathcal{M}$ avec $\varprojlim$ , on a} \texttt{[un]} \textit{isomorphisme}
\[
\varprojlim_U g^*(F) \xleftarrow{\ \sim\ } \varprojlim_V F.
\]
\end{enumerate}

\textsc{Démonstration}. Je vais d'abord expliciter le

\textbf{Lemme 7.5.} \textit{Soit $h : X' \longrightarrow X$ un morphisme étale dans} Cat. \textit{Alors $h$ est localement cofiltrant. De façon plus précise, pour tout $x$ dans $X$, on a}
\[
\pi_0(x\backslash X') \xrightarrow{\ \sim\ } \mathrm{Ob}\,X'_x,
\]
\textit{toute composante connexe de $x\backslash X'$ contient un élément et un seul de la catégorie discrète $X'_x \subseteq x\backslash X'$. Si $x'$ est dans $X'_x$ (i.e. $x'$ objet de $X'$ tel que $h(x') = x$), c'est un objet} initial \textit{de la composante connexe qui le contient. (Donc cette composante est cofiltrante.)}

Ceci prouve (et au delà) la première assertion de 7.4, compte tenu de 7.3 (stabilité par composition), ou parce qu'une \msp{page 94} catégorie étale (en l'occurrence $y\backslash X'$) au dessus d'une catégorie localement cofiltrante (resp. localement faiblement cofiltrante) (en l'occurrence $y\backslash X$) est itou.
\[
\begin{tikzcd}[row sep=large]
y\backslash X' \arrow[r] \arrow[d, "\text{ét.}"'] \arrow[dr, phantom, "\text{cart.}"] & X' \arrow[d, "\text{ét.}"] \\
y\backslash X \arrow[r] \arrow[d] \arrow[dr, phantom, "\text{cart.}"] & X \arrow[d] \\
y\backslash Y \arrow[r] & Y.
\end{tikzcd}
\]
Pour la deuxième assertion dans a), on peut maintenant supposer $X' = X$. Il faut prouver que pour tout $y'$ dans $Y'$, $y'\backslash X \overset{\text{déf}}{=} X \times_{Y'} y'\backslash Y'$ est localement cofiltrante (resp. localement faiblement cofiltrante). Soit $y$ dans $Y$ l'image de $y'$, et considérons le diagramme de trois carrés cartésiens ci-contre.
\[
\begin{tikzcd}[row sep=large]
y'\backslash X \arrow[r] \arrow[d, "\text{ét.}"] \arrow[dr, phantom, "\textit{cart.}"] & y\backslash X \arrow[r] \arrow[d, "\text{ét.}"] \arrow[dr, phantom, "\textit{cart.}"] & X \arrow[d, "\text{ét.}"] \\
y'\backslash Y' \arrow[r] & y\backslash Y' \arrow[r] \arrow[d] \arrow[dr, phantom, "\textit{cart.}"] & Y' \arrow[d] \\
 & y\backslash Y \arrow[r] & Y
\end{tikzcd}
\]
Le Lemme 7.5 implique que $y'\backslash Y' \longrightarrow y\backslash Y'$ est l'inclusion dans $y\backslash Y'$ d'une composante connexe de $y\backslash Y'$, c'est donc une immersion ouverte. Donc $y'\backslash X \longrightarrow y\backslash X$ est une immersion
% foot: Version 8 juillet 2001

% ---- from Der-ChII-234_p061-065.tex
%======================================================================
\printedpage{59}
% head: A. Grothendieck : Dérivateurs, Ch. II, Édition des Universités de Montpellier II et Paris VII

ouverte. Comme $y\backslash X$ est localement cofiltrant (resp. localement faiblement cofiltrant) par hypothèse, il en est donc de même de $y\backslash X$. Cela prouve d).

Appliquant a) à la situation $X/x \longrightarrow Y/y$ dans b), on voit que les colocalisés $z\backslash(X/x)$ de
\margmark{!} % margin: "!" with a triple vertical rule alongside the lines "X/x sur Y/y ... donc est 0-connexe."
$X/x$ sur $Y/y$ ($z \in \Ob\,(Y/y)$) sont localement cofiltrants (resp. faiblement). Mais comme $X/x$ a un objet final au dessus de l'objet final de $Y/y$, il s'ensuit que $z\backslash(X/x)$ a un objet final, donc est $0$-\textit{connexe.}

\msp{page 95}
Donc ces catégories, étant localement cofiltrantes (resp. faiblement), sont en fait cofiltrantes (resp. faiblement cofiltrantes).

Prouvons c). Notons que pour $z$ dans $V$, $z\backslash U$ s'identifie à l'image inverse, par $z\backslash(X/x) \longrightarrow X/x$, de l'ouvert $U$ de $X/x$, c'est donc un \textit{ouvert} de $z\backslash(X/x)$, non vide par construction de $V$. Mais un ouvert non vide d'une catégorie cofiltrante (resp. faiblement cofiltrante) est encore cofiltrant (resp. faiblement cofiltrant). Cela signifie que $g : U \longrightarrow V$ est cofiltrant (resp. faiblement cofiltrant), a fortiori cofinal (et $W$-coasphérique, pour tout $W$ satisfaisant Loc (8)). Cela achève la démonstration.

\dgrblock{dgrb143}{Why the restriction argument works}
\begin{dcomment}{exposition}{p95:Cela achève la démonstration.}{Why the restriction argument works}
The mechanism of Proposition 7.4 is the decomposition in Lemma 7.5, and it is useful to see why that decomposition is canonical. An \emph{\'etale} functor $h:E\to C$ is a discrete fibration. Given $(e,\alpha:c\to h(e))$ in $c\backslash E$, lift $\alpha$ uniquely to an arrow $e_0\to e$ with $h(e_0)=c$. This assigns an object $e_0$ of the discrete fibre to each object of the comma category. Every arrow of the comma category preserves $e_0$, by uniqueness of the lift. For each fixed $e_0$, the object $(e_0,1_c)$ is initial in the corresponding full subcategory. Thus
\[
c\backslash E\ \cong\
\coprod_{e_0\in h^{-1}(c)}(e_0\backslash E),
\]
and each displayed summand has an initial object. This proves all assertions of the lemma, including the empty-fibre case.

For the first assertion of 7.4(a), $y\backslash X'\to y\backslash X$ is a discrete fibration. Conditions $P$ and $E$ pass to its total category: lift a commutative cone in the base to the prescribed object above its target; uniqueness makes the lifted arrows commute. For $E$, faithfulness of a discrete fibration makes the lifted equalization an actual equality. This proves preservation of local cofilteredness and its $P$-only variant.

For the second assertion, let $X\xrightarrow{g}Y'\xrightarrow{h}Y$ be the factorization, with $h$ a discrete fibration. An object $y'\in Y'$ over $y$ determines the summand of $y\backslash Y'$ whose initial object is $(y',1_y)$. Its inverse image under $y\backslash X\to y\backslash Y'$ is exactly $y'\backslash X$. Hence $y'\backslash X$ is a union of connected components of $y\backslash X$, so it inherits $P$ and $E$. This argument also explains the second diagram: the labels ``\'et.'' on its three vertical arrows are not justified, since $g:X\to Y'$ is arbitrary. The relevant map is the horizontal component inclusion and its pullback. The claimed result does not use those incorrect vertical labels.

Now set $y=f(x)$. The category $X/x$ is \emph{\'etale} over $X$, and $Y/y$ is \emph{\'etale} over $Y$. Part (a) therefore makes $f/x:X/x\to Y/y$ locally cofiltered. For $z=(z_0\to y)$, its comma category $z\backslash(X/x)$ has a terminal object: take $(x,1_x)$ with the specified arrow $z_0\to y$. Every other object includes its unique morphism to this object as part of its structure. The comma category is consequently nonempty and connected. Combined with local cofilteredness this proves (b), and the same reasoning gives connectedness plus $P$ in the weaker case.

For (c), use the corrected $V$ above. The category $z\backslash U$ is a nonempty lower open full subcategory of $z\backslash(X/x)$. A nonempty lower open subcategory of a cofiltered category is cofiltered: for two objects, any common predecessor still lies in the open, and the same is true of the source of an equalizing arrow. For connectedness plus $P$, the common-predecessor argument proves connectedness and the square condition restricts. These observations prove (c), and also prove Remark 7.5(1) for an arbitrary cofiltered $f:X\to Y$ and an open $U\subset X$. In that remark the defining condition for $V$ must likewise be $y\backslash U\ne\varnothing$, as its preceding words ``engendr\'e par $f(U)$'' already require.
\end{dcomment}


\noindent\textbf{Remarque 7.5.}
\begin{itemize}
\item[1)] Le même argument que dans c) montre que si $f : X \longrightarrow Y$ est cofiltrant (resp. faiblement cofiltrant), alors pour tout ouvert $U$ dans $X$, si on désigne par $V$ l'ouvert de $Y$ engendré par $f(U)$, i.e. l'ouvert formé des $y$ dans $Y$ tels que $y\backslash X \neq \emptyset$, alors la flèche induite $g : U \longrightarrow V$ est cofiltrante (resp. faiblement cofiltrante), et a fortiori $g$ est cofinal (et plus généralement, si $f$ est cofiltrant, $g$ est $W$-coasphérique pour tout $W$ satisfaisant la condition Loc (8)).

\msp{page 96}
\item[2)] Un cas particulier intéressant de 7.4 est celui où on suppose que les $y\backslash X$ ($y \in \Ob\,Y$) sont sommes directes de catégories ayant chacune un \textit{objet initial.} Alors on voit que cette condition est stable par les opérations de 7.4 a) [\,?], et par conséquent, les colocalisés $z\backslash(X/x)$ des morphismes localisés en haut $X/x \longrightarrow Y/y$ sont des catégories à objet initial. Ceci implique alors qu'il en est de même des colocalisés $z\backslash U$ ($z \in \Ob\,V$) dans $g : U \longrightarrow V$. Il s'ensuit que $g$ est $W$-coasphérique pour \textit{tout} localiseur fondamental $W$ (sans avoir à se soucier de l'axiome des limites).

La condition envisagée sur $f : X \longrightarrow Y$ est satisfaite notamment dans les deux cas suivants (qui me paraissent de loin les plus importants) : a) $f$ est étale et b) $X$ est discrète.
\end{itemize}
% foot: Version 8 juillet 2001

%======================================================================
\clearpage
\printedpage{60}
% head: A. Grothendieck : Dérivateurs, Ch. II, Édition des Universités de Montpellier II et Paris VII

\msp{page 97}

\dgrblock{dgrb144}{Initial objects remove the limit axiom}
\begin{dcomment}{exposition}{p96}{Initial objects remove the limit axiom}
Grothendieck singles out components with initial objects because here the cofinality argument needs no passage to filtered limits or colimits. The bracketed question in Remark 7.5(2) can be settled directly. Suppose each $y\backslash X$ is a disjoint union of categories with initial objects. If $E\to C$ is a discrete fibration and every component of $C$ has initial object $c_0$, each object of $E$ receives a unique lift of the arrow from the relevant $c_0$. The lifts with sources over $c_0$ partition $E$ into full subcategories, each having its source as an initial object. Consequently the initial-object condition is preserved by the first operation in 7.4(a). The second operation selects a union of components, by the preceding decomposition, and therefore preserves it as well.

For $f/x:X/x\to Y/f(x)$ the comma categories have terminal objects and hence just one component. Their initial-object decomposition therefore consists of a single summand, so they actually have initial objects. Passing to a nonempty lower open $z\backslash U$ preserves this initial object: if an open contains any object, it contains the source of every arrow into that object, in particular the initial object. This proves the strengthened assertion in full.

The examples in the text follow without an additional hypothesis. For a discrete fibration $f$, Lemma 7.5 supplies the initial objects. If $X$ is discrete, every $y\backslash X$ is discrete, so all of its components are singleton categories. The strengthened conclusion can therefore apply to a functor which is far from final before one restricts to the open image.

For a fundamental localizer $W$, a category with initial object is $W$-aspherical. One can justify the initial-object version from the terminal-object axiom: the unique map to the terminal category and its section have composites related to the identity by the natural transformation from the constant initial-object functor. Fundamental localizers identify naturally transformed functors after passage to homotopy, as follows from the interval projection $C\times[1]\to C$ and the terminal-object axiom. Thus all the relevant comma categories are $W$-aspherical for every fundamental localizer. This is exactly where the appeal to Loc~(8) disappears in the manuscript's argument; mere weak cofilteredness would not give this conclusion.

The later theory removes the same qualification from the full cofiltered case. A cofiltered category has weakly contractible nerve: any finite collection of simplices admits a compatible cone after finitely many common-predecessor and equalization operations, giving the required null homotopies. \citet[Theorem 2.2.11, p.~136]{Cisinski2004Minimal} proves that the weak equivalences detected by nerves form the minimal fundamental localizer. They therefore belong to every fundamental localizer. Consequently every cofiltered comma category is $W$-aspherical for every fundamental $W$, and a cofiltered functor is $W$-coaspherical without an additional Loc~(8) assumption. The source's conditional assertion remains correct; the stronger present statement uses this later minimality theorem.
\end{dcomment}


\begin{center}
{\small
\begin{tikzcd}[ampersand replacement=\&, column sep=small, row sep=large]
\left\{\begin{tabular}{@{}c@{\ }l@{}}
 & Top $X$\\
 & localement totalement 0-connexe\\
 & et 0-connexe\\
$\Longleftrightarrow$ & totalement 0-connexe\\
 & et localement irréductible\\
$\Longleftrightarrow$ & $X$ \textit{cofiltrante,}\\
 & i.e. $\PS^*1' + \PS^*2$\\
$\Longleftrightarrow$ & $X$ 0-connexe\\
 & $+ \PS^*1 + \PS^*2$\\
$\Longleftrightarrow$ & $\left\{\begin{tabular}{@{}l@{}}
$\forall\, x \in \Ob\,X$, $X/x$ et\\
$X/x \times x$ sont irréductible,\\ % [sic] irréductible -> irréductibles, agreement with "sont"
$X$ 0-connexe
\end{tabular}\right.$\\
$\Longleftrightarrow$ & $\left\{\begin{tabular}{@{}l@{}}
$X$ 0-connexe,\\
les $X/a \times_x b$ non vides\\
(et même 0-connexes),\\
les $X/a \times_{x\times x} b$ 0-connexes
\end{tabular}\right.$
\end{tabular}\right.
\;
\left[\begin{tabular}{@{}l@{}}
stable par\\
passage\\
à $X' \longrightarrow X$\\
étale sur $X$, si\\
$X'$ 0-connexe,\\
\textit{mais} si les\\
$X/x$ sont cofiltrantes\\
et $X$ 0-connexe,\\
cela n'implique\\
\textit{pas} $X$ cofiltrante\\
Ex. $B_G$
\end{tabular}\right]
\arrow[rrrd, Rightarrow, "{\tl{stricte\\(ex. $\bullet\ \bullet$)}}"]
\arrow[rrdd, Rightarrow, "{\tl{stricte\\(ex. $\boldsymbol{\Delta}$)}}"]
\arrow[ddd, Rightarrow, "{\tl{stricte (ex.\\$B_G$)}}"']
\& \& \& |[draw,rounded corners=6pt]| \varepsilon = \bullet\ \bullet \\
\& \& \&
\begin{tabular}{@{}l@{}}
\begin{tabular}{@{}c@{\ }l@{}}
 & Top $X$ localement\\
 & totalement 0-connexe\\
$\Longleftrightarrow$ & les composantes connexes de $X$\\
 & sont cofiltrantes\\
$\Longleftrightarrow$ & $\PS^*1 + \PS^*2$\\
 & (X \textit{localement cofiltrante}) % [sic] X (upright roman) -> $X$ (italic), inconsistent with all other occurrences
\end{tabular}\\[2ex]
$\left[\begin{tabular}{@{}l@{}}
Stable par passage à $X' \longrightarrow X$\\
étale, mais si les $X'/x$\\
y satisfont, pas forcément $X'$!
\end{tabular}\right]$
\end{tabular}
\arrow[d, Rightarrow] \\
\& \&
\begin{tabular}{@{}l@{}}
\begin{tabular}{@{}c@{\ }l@{}}
 & Top $X$ totalement 0-connexe\\
 & (on dit $X$ \textit{totalement} 0-\textit{connexe})\\
$\Longleftrightarrow$ & $X^o \longrightarrow X^o \times X^o$ est cofinal\\
 & \textit{et} $X \neq \emptyset$\\
$\Longleftrightarrow$ & $\forall\, x, y \in \Ob\,X$, $x \times y$ 0-connexe de $X^\wedge$\\
 & i.e. $X/x \times y$ 0-connexe\\
 & \textit{et} $X \neq \emptyset$
\end{tabular}\\[2ex]
$\left[\begin{tabular}{@{}l@{}}
NB n'est pas stable par passage aux $X/x$ (\footnotemark[66]), et si\\
les $X/x$ sont totalement 0-connexes et $X \neq \emptyset$,\\
il ne s'ensuit pas que\\
$X$ le soit (\footnotemark[67]).
\end{tabular}\right]$
\end{tabular}
\arrow[dddl, Rightarrow, "{\tl{stricte (ex. $B_G$, ou\\$0 \rightrightarrows 1$)}}"]
\&
\tl{Top $X$\\localement irréductible} \\
\begin{tabular}{@{}c@{\ }l@{}}
 & Top $X$ localement\\
 & irréductible\\
 & $+$ 0-connexe\\
$\Longleftrightarrow$ & $X$ \textit{pseudofiltrante}\\
 & $+$ 0-\textit{connexe,}\\
 & i.e. $\PS^*1$\\
 & $+$ 0-connexe
\end{tabular}
\arrow[ddr, Rightarrow, "{\tl{stricte\\(ex. $0 \rightrightarrows 1$)}}"']
\&
\left[\begin{tabular}{@{}l@{}}
$\PS^*1$ ne dépend\\
que des catégories\\
induites $X/x$,\\
donc stable par passage\\
aux $X' \longrightarrow X$ étales,\\
donc la condition\\
$\PS^*1 +{}$0-connexe stable\\
par passage aux\\
revêtements\\
$X' \longrightarrow X$ 0-\textit{connexes.}
\end{tabular}\right]
\& \& \\
|[draw,rounded corners=6pt]| B_G \& \& \& \\
\&
\begin{tabular}{@{}c@{\ }l@{}}
 & Top $X$ irréductible\\
 & ($X$ \textit{irréductible})\\
$\Longleftrightarrow$ & $\PS^*1'$
\end{tabular}
\& \& \\
|[draw,rounded corners=6pt]| 0 \rightrightarrows 1
\&
\left[\begin{tabular}{@{}l@{}}
ne dépend que \texttt{[de]} $\mathrm{Ord}\,X$ ou de $\mathrm{Esp}(X)$,\\
signifie que $\mathrm{Ord}\,X$ cofiltrant,\\
ou $\mathrm{Esp}(X)$ irréductible.\\
N'est pas stable par passage à\\
$X' \longrightarrow X$ étale 0-connexe,\\
p.ex. par passage à un $X/x$,\\
p.ex. $X = (0 \rightrightarrows 1)$ est irréductible,\\
mais $X/1 = \left(\begin{tabular}{@{}l@{}}$\cdot$\\ $\;\searrow$\\ $\qquad\cdot$\\ $\;\nearrow$\\ $\cdot$\end{tabular}\right)$ ne l'est\\
pas.
\end{tabular}\right]
\& \& \\
\end{tikzcd}
}
\end{center}
% footnote marks (66), (67) are set inside the diagram; texts below
\footnotetext[66]{\textbf{Ex.} \circled{$\boldsymbol{\Delta}$}.}
\footnotetext[67]{\textbf{Ex.} $B_G$.}
% foot: Version 8 juillet 2001

%======================================================================
\clearpage
\printedpage{61}
% head: A. Grothendieck : Dérivateurs, Ch. II, Édition des Universités de Montpellier II et Paris VII

\msp{page 98}

\dgrblock{dgrb145}{Reading the six classes without the changing terminology}
\begin{dcomment}{examples}{p97}{Reading the six classes without the changing terminology}
The diagrams collect the distinctions developed earlier in the chapter. Their content is more stable than their terminology: ``pseudo-filtrante'' and ``pseudo-cofiltrante'' sometimes omit the direction or change meaning. We use $P,E,I$ as above and write $T$ for the assertion that $C^{\rm op}$ is sifted. Explicitly, $T$ means $C\ne\varnothing$ and the category of cones $c\to a$, $c\to b$ is connected for every $a,b$. This is Grothendieck's total $0$-connectedness, as discussed already in the commentary at manuscript p.~18 (``The ordinary sifted condition and its variance''). The six classes are
\[
\begin{array}{c|l}
C_{\rm cof}&I+E\quad(\text{equivalently }P+E+\text{connected, nonempty})\\
C_{\rm loc}&P+E\\
C_{\rm weak}&P+\text{connected, nonempty}\\
C_{\rm pre}&P\\
C_{\rm irr}&I\\
C_{\rm tot}&T.
\end{array}
\]
There are inclusions $C_{\rm cof}\subset C_{\rm loc},C_{\rm weak},C_{\rm tot}$, $C_{\rm loc}\subset C_{\rm pre}$, $C_{\rm weak}\subset C_{\rm pre},C_{\rm irr}$, and $C_{\rm tot}\subset C_{\rm irr}$. The two middle classes $C_{\rm weak}$ and $C_{\rm tot}$ meet exactly in $C_{\rm cof}$, by the criterion established earlier in the chapter. The following table makes the five separating examples completely explicit; $G$ is nontrivial and $A$ is the parallel-arrow category.
\[
\begin{array}{c|cccccc}
 &C_{\rm cof}&C_{\rm loc}&C_{\rm weak}&C_{\rm pre}&C_{\rm irr}&C_{\rm tot}\\\hline
B_G&0&0&1&1&1&0\\
\boldsymbol\Delta&0&0&0&0&1&1\\
A=(0\rightrightarrows1)&0&0&0&0&1&0\\
\varepsilon=\{0,1\}&0&1&0&1&0&0\\
B_G\amalg e&0&0&0&1&0&0
\end{array}
\]

For $B_G$, every cospan can be completed because its arrows are invertible: for $r,s\in G$, use $u=r^{-1}s$ and $v=1$, obtaining $ru=sv$. Equalization fails for a nonidentity $g$ and $1$, since $gt=t$ would imply $g=1$. The common-cone category of the unique object with itself has objects $(r,s)\in G^2$, with simultaneous right multiplication by $t\in G$ as morphisms. The quantity $sr^{-1}$ labels its connected components, giving $|G|$ components. Thus $T$ fails. This is the same obstruction as the diagonal action on $G\times G$; see \citet[Lemma 4.19.4]{StacksFiltered2026}. Every slice $B_G/*$ nevertheless has exactly one arrow between any pair of its objects and is cofiltered. This verifies the diagram's warning that cofiltered slices do not imply a locally cofiltered category.

For $A$, the object $0$ is a common predecessor of any pair of objects, so $I$ holds. The cospan given by the two distinct arrows $r,s:0\to1$ cannot be completed to a commutative square: the only arrow into $0$ is its identity. Thus $P$ fails. The cone category for the pair $(0,1)$ has exactly two objects and no nonidentity arrows, so $T$ fails too. The slice $A/1$ has the shape $r\to1\leftarrow s$ and is not irreducible. This simultaneously explains the failure of irreducibility to survive a connected \emph{\'etale} change of source.

For $\boldsymbol\Delta$, the category of nonempty finite ordinals and monotone maps, every pair admits a common predecessor $[0]$. A cone from $[k]$ to $[m]$ and $[n]$ is a finite monotone string in the rectangular poset $[m]\times[n]$. It is connected, in its cone category, to each of its vertices by the corresponding map $[0]\to[k]$. Any two vertices can be joined by moving horizontally and vertically through the rectangle, each move being represented by a $[1]$-string. Thus the cone category is nonempty and connected, proving $T$. The two endpoint maps $[0]\rightrightarrows[1]$ cannot be equalized even by two different maps into their sources: every nonempty ordinal has a unique map to $[0]$, and the two resulting maps to $[1]$ remain distinct. Hence $P$ fails. In particular $\Delta/[1]$ is not irreducible: the two endpoint objects have no common predecessor. Total connectedness therefore does not pass to slices.

For $\varepsilon$, the conditions $P$ and $E$ hold separately on each singleton component, while all four classes requiring connectedness fail. For $B_G\amalg e$, $P$ still holds componentwise, $E$ fails in the group component, and connectedness fails globally. These calculations also settle the abbreviated ``Itou pour $\varepsilon$'' on printed page 62: it can refer to exclusion from the four connected classes, but cannot refer to failure of local cofilteredness, since $\varepsilon$ is locally cofiltered.

Finally, the product notation on page 97 takes place in the presheaf category. Thus $C/(a\times_x b)$ means the category of elements of $h_a\times_{h_x}h_b$, and $C/(x\times x)$ means that of $h_x\times h_x$; it is not a product of a slice with an object. The irreducibility of every $C/x$ is $P$. Together with irreducibility of every $C/(x\times x)$ it forces $E$: compare the object defined by $r,s:a\to x$ with $(x,1_x,1_x)$ and take a common predecessor. Conversely a cofiltered category provides cones and equalizations for these finite families, so the displayed categories of elements are cofiltered. This explains the strengthened representable tests in the first diagram, without repeating the longer topos arguments of the commentary on the first part of the chapter.
\end{dcomment}


\begin{center}
{\small
\begin{tikzcd}[ampersand replacement=\&, column sep=small, row sep=large]
\& \tl{Top $X$ localement totalement 0-connexe\\$+$ 0-connexe\\($X$ \textit{cofiltrante})}
\arrow[ddl, Rightarrow] \arrow[ddr, Rightarrow] \arrow[rrdd, Rightarrow]
\& \& \\
\& \& \& |[draw,rounded corners=6pt]| \varepsilon = \bullet\ \bullet \\
\tl{Top $X$ localement irréductible\\$+$ 0-connexe\\($X$ \textit{pseudofiltrante} et 0-\textit{connexe})\\$\PS^*1 +{}$ 0-connexe}
\arrow[ddr, Rightarrow]
\& \&
\tl{Top $X$ totalement 0-connexe\\($X$ \textit{totalement} 0-\textit{connexe})\\$X^o \longrightarrow X^o \times X^o$ cofinal}
\arrow[ddl, Rightarrow]
\&
\tl{Top $X$ localement totalement\\0-connexe\\($X$ \textit{localement cofiltrante})\\$\PS^*1 + \PS^*2$}
\arrow[dd, Rightarrow] \\
|[draw,rounded corners=6pt]| B_G \& \& |[draw,circle,inner sep=1pt]| \boldsymbol{\Delta} \& \\
\& \tl{Top $X$ irréductible\\($X$ \textit{irréductible})\\$\PS^*1'$}
\& \&
\tl{Top $X$ localement irréductible\\($X$ \textit{pseudo-cofiltrante} [\,?])\\$\PS^*1$} \\
\& |[draw,rounded corners=6pt]| 0 \rightrightarrows 1 \& \&
|[draw,rounded corners=12pt]| \begin{tabular}{@{}c@{}}$B_G \amalg e$\\avec $G \neq 1$\end{tabular} \\
\end{tikzcd}
}
\end{center}

\begin{center}
{\small
\begin{tikzcd}[ampersand replacement=\&, column sep=small, row sep=large]
\text{(\footnotemark[68]).}
\& \tl{$X$ \textit{cofiltrante}\\$\PS^*1' + \PS^*2$\\$= \PS^*1 + \PS^*2 +{}$ 0-connexe}
\arrow[ddl, Rightarrow] \arrow[ddr, Rightarrow] \arrow[rrdd, Rightarrow]
\& \& \\
\& \& \& |[draw,rounded corners=6pt]| \bullet\ \bullet \\
\tl{$X$ \textit{pseudo-cofiltrante} et 0-\textit{connexe}\\($\PS^*1 +{}$ 0-connexe)}
\arrow[ddr, Rightarrow] \arrow[ddrrr, Rightarrow]
\& \&
\tl{$X$ \textit{totalement} 0-\textit{connexe}\\($X^o \longrightarrow X^o \times X^o$ cofinal)}
\arrow[ddl, Rightarrow]
\&
\tl{$X$ \textit{localement cofiltrante}\\$\PS^*1 + \PS^*2$}
\arrow[dd, Rightarrow] \\
|[draw,rounded corners=6pt]| B_G \& \& |[draw,circle,inner sep=1pt]| \boldsymbol{\Delta} \& \\
\& \tl{$X$ \textit{irréductible}\\$\PS^*1'$}
\& \&
\tl{$X$ \textit{pseudo-cofiltrante}\\$\PS^*1$} \\
\& |[draw,rounded corners=6pt]| \bullet \rightrightarrows \bullet \& \&
|[draw,rounded corners=6pt]| B_G \amalg e \\
\end{tikzcd}
}
\end{center}
\footnotetext[68]{\textbf{NB} Les cinq catégories typiques $B_G$ ($G \neq 1$), $\boldsymbol{\Delta}$, $0 \rightrightarrows 1$, $\varepsilon = \bullet\ \bullet$, $B_G \amalg e$ ($G \neq 1$) sont telles qu'elles appartiennent à la classe envisagée $C$ d'objets de Cat, mais à aucune des cinq autres classes sauf celles qui sont reliées à $C$ par $C \Longrightarrow C'$ sur le diagramme d'implications (ou par une implication composée).}
% foot: Version 8 juillet 2001

%======================================================================
\clearpage
\printedpage{62}
% head: A. Grothendieck : Dérivateurs, Ch. II, Édition des Universités de Montpellier II et Paris VII

\noindent\begin{tabular}{@{}p{2.8cm}@{}p{11.5cm}@{}}
$\bullet \rightrightarrows \bullet$ & n'est ni pseudo-cofiltrante, ni totalement 0-connexe.\\[1ex]
$B_G$ & n'est ni localement cofiltrante, ni totalement 0-connexe.\\[1ex]
$\boldsymbol{\Delta}$ & n'est pas pseudo-filtrante.\\[1ex]
$B_G \amalg e$ & n'est pas localement cofiltante, car $B_G$ ne l'est pas, ni appartient à une des 4 classes d'objets connexes.\\[1ex] % [sic] cofiltante -> cofiltrante, letter r missing
Itou pour & $\varepsilon = \bullet\ \bullet$
\end{tabular}

\msp{page 99}

\begin{center}
{\small
\begin{tikzcd}[ampersand replacement=\&, column sep=small, row sep=large]
\& \tl{$X$ \textit{cofiltrante}\\($\PS^*1' + \PS^*2$\\$\Longleftrightarrow \PS^*1 + \PS^*2$\\$+$ 0-connexe)}
\arrow[ddl, Rightarrow, "{\tls{$\PS^*1'$\\($\Longrightarrow$ 0-connexe)}}"]
\arrow[dd, Rightarrow, "{\tls{$\PS^*1$\\(les $x\backslash X$ satisfont\\$\PS^*1'$)}}"]
\arrow[ddr, Rightarrow]
\& \\
% layout: the two small labels stand between the arrows; attribution to the arrows as in this layout
\& \& |[draw,rounded corners=6pt]| \varepsilon = \bullet\ \bullet \\
\tl{$X$ totalement 0-connexe\\($X^o \longrightarrow X^o \times X^o$ cofinal)}
\arrow[ddd, Rightarrow]
\&
\tl{$X$ \textit{pseudo-cofiltrante}\\0-connexe\\($\PS^*1 +{}$ 0-connexe)\\(manque $\PS^*2$)}
\arrow[dddl, Rightarrow] \arrow[dddr, Rightarrow]
\&
\tl{$X$ \textit{localement cofiltrante}\\(composantes connexes filtrantes)\\($\PS^*1 + \PS^*2$)}
\arrow[ddd, Rightarrow] \\
|[draw,circle,inner sep=1pt]| \boldsymbol{\Delta} \& |[draw,rounded corners=6pt]| B_G \& \\
\& G \neq \{1\} \& \\
\tl{$X$ irréductible\\($\PS^*1'$)}
\& \&
\tl{$X$ \textit{pseudo-cofiltrante}\\ou encore \textit{localement irréductible}\\(composantes connexes\\pseudo-filtrantes)\\($\PS^*1$)} \\
|[draw,rounded corners=6pt]| \bullet \rightrightarrows \bullet \& \& |[draw,rounded corners=6pt]| B_G \amalg e \\
\& \& G \neq \{1\} \\
\end{tikzcd}
}
\end{center}

et dualement on a

\begin{center}
{\small
\begin{tikzcd}[ampersand replacement=\&, column sep=small, row sep=huge]
\& \text{$X$ filtrante}
\arrow[dl, Rightarrow] \arrow[d, Rightarrow] \arrow[dr, Rightarrow]
\& \\
\text{$X^o$ totalement 0-connexe}
\arrow[d, Rightarrow]
\&
\tl{$X$ pseudo-filtrante\\0-connexe}
\arrow[dl, Rightarrow] \arrow[dr, Rightarrow]
\&
\text{$X$ localement filtrante}
\arrow[d, Rightarrow] \\
\text{$X$ coirréductible} \& \& \text{$X$ pseudo-cofiltrante} \\
\end{tikzcd}
}
\end{center}
% foot: Version 8 juillet 2001

%======================================================================
\clearpage
\printedpage{63}
% head: A. Grothendieck : Dérivateurs, Ch. II, Édition des Universités de Montpellier II et Paris VII

donc on a une terminologie directe sauf pour `$X^o$ totalement 0-connexe' - je n'ose dire `$X$ totalement co-connexe'!!
% foot: Version 8 juillet 2001


\dgrblock{dgrb146}{Representable flatness and the pointwise Kan extension}
\begin{dcomment}{later}{p99}{Representable flatness and the pointwise Kan extension}
Grothendieck's cofiltered functors have a standard modern name: they are \emph{representably flat}. For each $y\in Y$, the covariant set-valued functor $X\to\mathrm{Set}$, $x\mapsto Y(y,f(x))$, has category of elements $y\backslash X$; requiring this category to be cofiltered is precisely the definition. The adjective matters, since flatness for a functor into a topos has another, internal interpretation. \citet[Remark 2.4]{Tendas2024} explicitly distinguishes these notions, following the discussion of \citet{Shulman2011Flat}.

One consequence places the closing section directly in the theory of presheaf topoi. For $F:X^{\rm op}\to\mathrm{Set}$, the left Kan extension along $f^{\rm op}$ is
\[
(f_!F)(y)=\mathop{\rm colim}_{(y\backslash X)^{\rm op}}F(x).
\]
If $f$ is cofiltered, the indexing category is filtered. Finite limits of presheaves are computed pointwise, and filtered colimits of sets commute with finite limits; hence $f_!$ preserves finite limits. Consequently the adjunction $f_!\dashv f^*$ is a geometric morphism in the direction determined by this inverse-image functor. This is an additional geometric adjunction alongside the usual $f^*\dashv f_*$.

If $f$ is only locally cofiltered, the same formula is a coproduct, indexed by $\pi_0(y\backslash X)$, of filtered colimits. It preserves finite \emph{connected} limits. To check the distinction, a compatible tuple in a connected finite diagram of such coproducts must lie throughout in one summand; within that summand the filtered-colimit argument applies. The terminal presheaf instead has image
\[
(f_!1)(y)\cong\pi_0(y\backslash X),
\]
which need not be a singleton. This separates the local condition from representable flatness exactly where the manuscript separates local cofilteredness from cofilteredness. The general connected-limit calculation appears in \citet[Lemma 4.19.9]{StacksFiltered2026}.

The last diagrams thus organize three different demands on colimits: finite limits, finite connected limits, and finite products. The product condition is the sifted condition on $X^{\rm op}$, whose later systematic treatment is \citet[Section 1]{AdamekRosickyVitale2010}. The displayed examples show why none of these demands can be substituted for another without its accompanying connectedness or square condition. The missing nonemptiness qualification in the abbreviated diagonal criteria on pages 98--99 should always be read from page 97: the empty category has a final diagonal vacuously, but is not sifted.
\end{dcomment}
\begin{otherlanguage}{english}
\dgroverview{II-b}{Chapter II \textperiodcentered\ \textup{II-b} \textperiodcentered\ Overview}
% Overview by GPT-6 Astra (OpenAI, ChatGPT Pro mode).
\begingroup
\selectlanguage{english}
\noindent\textbf{Reconstruction, preservation, and the geometry of indexing categories.}

These pages complete the second chapter by asking how much of a category is
visible to its diagram theories. The starting point is the cofinality theorem:
a restriction of diagrams preserves a specified limit or colimit when the
appropriate comma categories are connected. The next question is more demanding.
When does restriction embed the whole category of diagrams fully faithfully?
Grothendieck then turns from individual functors to classes of functors. He
tries to recover the classes controlling change of base from a chosen class
of weak equivalences, and develops the order-theoretic machinery needed to
state such a recovery. The closing discussion of cofiltered categories gives
concrete tests and counterexamples for this programme. These are successive
parts of one investigation: which local tests retain exactly the information
that a theory of coefficients can detect?

For ordinary presheaves, the answer to the full-faithfulness question is a
two-sided cofinality criterion. Given $u:A\to B$ and an arrow
$\gamma:b_0\to b_1$, form the category of factorizations
$b_0\to u(a)\to b_1$ whose composite is $\gamma$. Restriction on
set-valued presheaves is fully faithful precisely when every such category
is nonempty and connected. The counit evaluated on a representable presheaf
has these connected-component sets as its fibres, and representables generate
all presheaves by colimits. This is the decisive argument of Proposition 4.1.
In later terminology the functors are lax epimorphisms, or absolutely dense
functors; the equivalence of these formulations is established by
\citet{AdamekElBashirSobralVelebil2001}. The commentary on the first part of the
chapter already treats the general criterion. Here the commentary supplies the
missing distinction between its two adjunction tests and checks the direction
of the comma-category maps.

The right-adjoint unit behaves differently. Its being invertible on every
representable says that the restricted Yoneda functor
$B\to\operatorname{Set}^{A^{\mathrm{op}}}$ is fully faithful: this is
density, in the usual sense \citep[Definition 6.5.11]{Riehl2016}. It does not
imply full faithfulness of restriction on all presheaves. The full category
of sets on a singleton and a two-element set, with $u$ selecting the
singleton, already gives a counterexample. All maps of the original category
are recovered from maps of sets of elements, but the identity of the
two-element set cannot factor through the singleton. Thus the hesitation on
manuscript pages 54--56 concerns a genuine obstruction. The commentary also
computes the four-object example on pages 57--58: every factorization
category is contractible, while one strict pullback is a disconnected
two-point category. The example separates reconstruction of diagrams from
universal preservation under strict base change.

For an arbitrary derivator $\D$, one must distinguish full faithfulness for
that particular coefficient theory from full faithfulness for every
derivator. Modern homotopical epimorphisms impose the latter quantifier.
\citet[Definition 6.1 and Theorem 6.10, arXiv v3]{GrothStovicek2016Trees}
characterize them by weak contractibility of the same factorization
categories. Grothendieck's fixed-$\D$ problem is subtler. Global diagrams
test only certain diagrams on a slice, and his proof of Proposition 4.4
carefully isolates the additional detection needed to recover all the slice
tests. A conservative left Kan extension from the slice suffices; essential
surjectivity of restriction to the slice is a stronger sufficient condition.
The pointwise Kan-extension and detection axioms of a derivator perform
different jobs in this argument \citep[Definition 1.5]{Groth2013}.

The developments in the commentary make two limits of these sufficient
conditions explicit. First, in the derivator represented by the complete
lattice $\{0<1\}$, diagrams on the free category generated by a directed
two-cycle are constant. Evaluation at one vertex is therefore fully
faithful. Nevertheless, the identity of the other vertex has no
factorization through it, and its factorization category is empty. This
answers the intended question about $\D$-connectedness versus
$\D$-biasphericity. Second, the parenthetical faithfulness assertion in
Corollary 4.7 fails for coherent derived diagrams. A discrete fibration from
a three-object fork to an interval sends a nonsplit extension of quiver
representations to a split extension. Direct sums in the coefficient
derived category are faithful, but the Kan extension is not. The
conservativity assertion survives: pointwise detection proves it. Each
development is given with a complete proof and a statement of the sources
checked, without a claim of priority.

The cancelled passage at the end of \S4 must be read with comparable care.
Let $S(A)$ be the maps whose pullback functors preserve a class $A$, and let
$N(L)$ be the maps whose pullbacks along a testing class $L$ lie in $W$.
Under the requisite identity and stability hypotheses, the inclusions and
the two idempotence calculations in the passage survive. The missing
ingredient is order reversal of $S$: enlarging the preserved class changes
both the inputs and the required outputs of a pullback test. It need not
give a smaller stabilizer. With $W$ the maximal localizer on categories and
$A$ the categorical equivalences, two point inclusions into the walking
isomorphism directly refute the stronger identity used in the printed
justification. This diagnosis preserves the valid calculations while
explaining why they do not constitute a Galois connection.

Section 5 supplies that connection in its proper generality. Any relation
between two sets defines two antitone polar operations. The composites are
closure operators, and their closed subsets form anti-isomorphic complete
lattices. When the sets are ordered, a further condition makes polar lower
sets principal and hence represents their closures by elements. Formula
(5.20) slips between these levels: a closed subset is not its greatest
element, and a join of closed elements need not be closed in the ambient
lattice. The corrected formula takes the closure of the join. Double
annihilation in $\mathbf F_2^2$ gives a short counterexample to the printed
equality. The classical examples also require their proper separation or
topological hypotheses. In infinite Galois theory the fixed subgroups are
the Krull-closed subgroups; the dense cyclic Frobenius subgroup in
$\widehat{\mathbf Z}$ shows why the qualification cannot be dropped
\citep[Proposition 7.12 and Theorem 7.13]{Milne2022FT}. A lifting relation is
another polarity, but a mutually closed pair of lifting classes still
requires a separate factorization theorem to become a weak factorization
system \citep[Definition 11.2.1]{Riehl2014}.

Section 6 repairs the earlier construction by fixing its target class.
The transporter $T_W(A)$ consists of maps whose pullback functors send
arrows of $A$ to arrows of $W$. Testing includes every arrow in the slice,
with any structure map from its codomain to the base. It is not enough to
test arrows whose codomain happens to equal the base. With $A$ closed under
composition and $L$ stable under pullback, both containing isomorphisms,
the equivalence
\[
 A\subseteq N_W(L)\quad\Longleftrightarrow\quad L\subseteq T_W(A)
\]
is a genuine antitone adjunction. The correction at (6.7) removes
composition closure from the right-hand domain, where it is not automatic.
A four-object chain in the commentary exhibits a closed transporter that
does not compose, under exactly the abstract assumptions of \S6. The
example is not claimed to use a fundamental localizer on $\mathrm{Cat}$.
Proposition 6.1 then identifies the additional condition: for a closed pair,
the right class is composition-closed exactly when it equals the stabilizer
of its left class.

This formalism has a direct historical continuation. Grothendieck's letter
to Thomason singles out a fibration class defined intrinsically from weak
equivalences, while treating factorization as another question
\citep[\S1, pp.~3--4]{Grothendieck1991Letter}. Rezk's sharp maps use the same
preservation condition. In a right proper model category they characterize
the maps all of whose strict pullback squares are homotopy cartesian
\citep[Propositions 2.2 and 2.7]{Rezk1998Sharp}. Maltsiniotis develops local
tests for the corresponding preservation of aspherical maps
\citep[Definition 4.6 and Theorem 4.13]{Maltsiniotis2005Asphericite}.
One question on page 89 has a particularly concrete answer: the stabilizer
of categorical equivalences is exactly the isofibrations. The commentary gives
both directions of the proof, using the walking isomorphism for necessity,
and relates the result to the categorical model structure
\citep[\S2 and Theorem 3.1]{Rezk1996Categories}. For the general mixed
transporters and the proposed bilocal description of the biaspherical
stabilizer, what was verified is kept apart from what was not. Failure to find
a precise matching theorem is not presented as proof that the questions
are open today.

The final section turns the abstract machinery into the geometry of
predecessors. A cofiltered functor is tested by cofiltered comma categories;
the local version allows several components and asks each to be
cofiltered. Over a poset, these conditions become statements about lower
open subsets. Proposition 7.4 is essentially sound, but its target lower
open must be generated by the image of the chosen open subset: taking the
whole principal lower set makes the printed condition false even for the
empty subset. The commentary proves the corrected criterion and the
initial-object assertion of Remark 7.5. It also exhibits a cofiltered map
whose arbitrary base change is not even locally weakly cofiltered, while
respecting the stated preservation under co-étale base change. For
fundamental localizers, the later minimality theorem removes the extra
Loc(8) qualification attached here to the asphericity of cofiltered
categories \citep[Theorem 2.2.11]{Cisinski2004Minimal}.

The final diagrams distinguish common predecessors, the filling of
cospans, equalization of parallel arrows, and connectedness of the
categories of cones over every finite family of objects. Their implications
become transparent when these requirements are
kept separate. Groups, the simplex category, a parallel-arrow category,
and disjoint unions give the required separating examples. The modern
interpretation of total connectedness is siftedness of the opposite
category; cofilteredness adds the equalization needed for commuting with
all finite limits, rather than only finite products
\citep{AdamekRosickyVitale2010,StacksFiltered2026,Tendas2024}. In this way the
chapter ends where its larger project began: the shape of an indexing
category governs which constructions a theory of diagrams can perform
and which information restriction can retain. The derived-diagram
motivation already appears in \emph{Pursuing Stacks}, \S69
\citep{Grothendieck1983}; here it is being refined into local, verifiable
conditions on functors.

\endgroup
\end{otherlanguage}
\dgrchunkend
